\subsection*{Textbook Formal Inventory}
This subsection records the exact formal material in \file{HAETAE_Rejection.ec}.  It is generated from the proof-surface EasyCrypt file, so the line numbers and statements are meant to be read next to the checked source.

\noindent\textbf{Chapter role.} rejection-sampling predicates, accepted-attempt structure, and loss obligations.

\noindent\textbf{Inventory size.} 12 context declarations, 101 definitions/game objects, 139 lemmas or relational theorems, and 0 explicit Hoare/Phoare judgments were found.

\subsubsection*{Theory Context, Imports, and Quantified Modules}
\paragraph{Context 1 (line 1)}
\noindent\textbf{EasyCrypt name.}
\begin{lstlisting}[language=EasyCrypt,basicstyle=\ttfamily\scriptsize]
imported theories
\end{lstlisting}
\noindent\textbf{Checked EasyCrypt statement.}
\begin{lstlisting}[language=EasyCrypt,basicstyle=\ttfamily\scriptsize]
require import AllCore Distr IntDiv List Real StdOrder RealSeries.
\end{lstlisting}
\noindent\textbf{Formal mathematical reading.}
Mathematical context. This line imports or specializes earlier theories, making their carrier sets, functions, modules, and theorems available to the current file.

\noindent\textbf{Role in the proof.}
Use in this chapter. It fixes the proof context, imported mathematical language, or quantified module parameters for the following statements.

\paragraph{Context 2 (line 2)}
\noindent\textbf{EasyCrypt name.}
\begin{lstlisting}[language=EasyCrypt,basicstyle=\ttfamily\scriptsize]
imported theories
\end{lstlisting}
\noindent\textbf{Checked EasyCrypt statement.}
\begin{lstlisting}[language=EasyCrypt,basicstyle=\ttfamily\scriptsize]
require import HAETAE_Params HAETAE_Algebra HAETAE_Distributions.
\end{lstlisting}
\noindent\textbf{Formal mathematical reading.}
Mathematical context. This line imports or specializes earlier theories, making their carrier sets, functions, modules, and theorems available to the current file.

\noindent\textbf{Role in the proof.}
Use in this chapter. It fixes the proof context, imported mathematical language, or quantified module parameters for the following statements.

\paragraph{Context 3 (line 3)}
\noindent\textbf{EasyCrypt name.}
\begin{lstlisting}[language=EasyCrypt,basicstyle=\ttfamily\scriptsize]
imported theories
\end{lstlisting}
\noindent\textbf{Checked EasyCrypt statement.}
\begin{lstlisting}[language=EasyCrypt,basicstyle=\ttfamily\scriptsize]
require import HAETAE_Events HAETAE_Reductions.
\end{lstlisting}
\noindent\textbf{Formal mathematical reading.}
Mathematical context. This line imports or specializes earlier theories, making their carrier sets, functions, modules, and theorems available to the current file.

\noindent\textbf{Role in the proof.}
Use in this chapter. It fixes the proof context, imported mathematical language, or quantified module parameters for the following statements.

\paragraph{Context 4 (line 4)}
\noindent\textbf{EasyCrypt name.}
\begin{lstlisting}[language=EasyCrypt,basicstyle=\ttfamily\scriptsize]
imported theories
\end{lstlisting}
\noindent\textbf{Checked EasyCrypt statement.}
\begin{lstlisting}[language=EasyCrypt,basicstyle=\ttfamily\scriptsize]
require import HAETAE_Transcript.
\end{lstlisting}
\noindent\textbf{Formal mathematical reading.}
Mathematical context. This line imports or specializes earlier theories, making their carrier sets, functions, modules, and theorems available to the current file.

\noindent\textbf{Role in the proof.}
Use in this chapter. It fixes the proof context, imported mathematical language, or quantified module parameters for the following statements.

\paragraph{Context 5 (line 6)}
\noindent\textbf{EasyCrypt name.}
\begin{lstlisting}[language=EasyCrypt,basicstyle=\ttfamily\scriptsize]
HAETAE_Rejection
\end{lstlisting}
\noindent\textbf{Checked EasyCrypt statement.}
\begin{lstlisting}[language=EasyCrypt,basicstyle=\ttfamily\scriptsize]
theory HAETAE_Rejection.
\end{lstlisting}
\noindent\textbf{Formal mathematical reading.}
Mathematical context. This begins a namespace for the formal development in this file.

\noindent\textbf{Role in the proof.}
Use in this chapter. It fixes the proof context, imported mathematical language, or quantified module parameters for the following statements.

\paragraph{Context 6 (line 8)}
\noindent\textbf{EasyCrypt name.}
\begin{lstlisting}[language=EasyCrypt,basicstyle=\ttfamily\scriptsize]
opened namespace
\end{lstlisting}
\noindent\textbf{Checked EasyCrypt statement.}
\begin{lstlisting}[language=EasyCrypt,basicstyle=\ttfamily\scriptsize]
import HAETAE_Params.
\end{lstlisting}
\noindent\textbf{Formal mathematical reading.}
Mathematical context. This line imports or specializes earlier theories, making their carrier sets, functions, modules, and theorems available to the current file.

\noindent\textbf{Role in the proof.}
Use in this chapter. It fixes the proof context, imported mathematical language, or quantified module parameters for the following statements.

\paragraph{Context 7 (line 9)}
\noindent\textbf{EasyCrypt name.}
\begin{lstlisting}[language=EasyCrypt,basicstyle=\ttfamily\scriptsize]
opened namespace
\end{lstlisting}
\noindent\textbf{Checked EasyCrypt statement.}
\begin{lstlisting}[language=EasyCrypt,basicstyle=\ttfamily\scriptsize]
import HAETAE_Algebra.
\end{lstlisting}
\noindent\textbf{Formal mathematical reading.}
Mathematical context. This line imports or specializes earlier theories, making their carrier sets, functions, modules, and theorems available to the current file.

\noindent\textbf{Role in the proof.}
Use in this chapter. It fixes the proof context, imported mathematical language, or quantified module parameters for the following statements.

\paragraph{Context 8 (line 10)}
\noindent\textbf{EasyCrypt name.}
\begin{lstlisting}[language=EasyCrypt,basicstyle=\ttfamily\scriptsize]
opened namespace
\end{lstlisting}
\noindent\textbf{Checked EasyCrypt statement.}
\begin{lstlisting}[language=EasyCrypt,basicstyle=\ttfamily\scriptsize]
import HAETAE_Distributions.
\end{lstlisting}
\noindent\textbf{Formal mathematical reading.}
Mathematical context. This line imports or specializes earlier theories, making their carrier sets, functions, modules, and theorems available to the current file.

\noindent\textbf{Role in the proof.}
Use in this chapter. It fixes the proof context, imported mathematical language, or quantified module parameters for the following statements.

\paragraph{Context 9 (line 11)}
\noindent\textbf{EasyCrypt name.}
\begin{lstlisting}[language=EasyCrypt,basicstyle=\ttfamily\scriptsize]
opened namespace
\end{lstlisting}
\noindent\textbf{Checked EasyCrypt statement.}
\begin{lstlisting}[language=EasyCrypt,basicstyle=\ttfamily\scriptsize]
import HAETAE_Events.
\end{lstlisting}
\noindent\textbf{Formal mathematical reading.}
Mathematical context. This line imports or specializes earlier theories, making their carrier sets, functions, modules, and theorems available to the current file.

\noindent\textbf{Role in the proof.}
Use in this chapter. It fixes the proof context, imported mathematical language, or quantified module parameters for the following statements.

\paragraph{Context 10 (line 12)}
\noindent\textbf{EasyCrypt name.}
\begin{lstlisting}[language=EasyCrypt,basicstyle=\ttfamily\scriptsize]
opened namespace
\end{lstlisting}
\noindent\textbf{Checked EasyCrypt statement.}
\begin{lstlisting}[language=EasyCrypt,basicstyle=\ttfamily\scriptsize]
import HAETAE_Reductions.
\end{lstlisting}
\noindent\textbf{Formal mathematical reading.}
Mathematical context. This line imports or specializes earlier theories, making their carrier sets, functions, modules, and theorems available to the current file.

\noindent\textbf{Role in the proof.}
Use in this chapter. It fixes the proof context, imported mathematical language, or quantified module parameters for the following statements.

\paragraph{Context 11 (line 13)}
\noindent\textbf{EasyCrypt name.}
\begin{lstlisting}[language=EasyCrypt,basicstyle=\ttfamily\scriptsize]
opened namespace
\end{lstlisting}
\noindent\textbf{Checked EasyCrypt statement.}
\begin{lstlisting}[language=EasyCrypt,basicstyle=\ttfamily\scriptsize]
import HAETAE_Transcript.
\end{lstlisting}
\noindent\textbf{Formal mathematical reading.}
Mathematical context. This line imports or specializes earlier theories, making their carrier sets, functions, modules, and theorems available to the current file.

\noindent\textbf{Role in the proof.}
Use in this chapter. It fixes the proof context, imported mathematical language, or quantified module parameters for the following statements.

\paragraph{Context 12 (line 14)}
\noindent\textbf{EasyCrypt name.}
\begin{lstlisting}[language=EasyCrypt,basicstyle=\ttfamily\scriptsize]
opened namespace
\end{lstlisting}
\noindent\textbf{Checked EasyCrypt statement.}
\begin{lstlisting}[language=EasyCrypt,basicstyle=\ttfamily\scriptsize]
import RealOrder.
\end{lstlisting}
\noindent\textbf{Formal mathematical reading.}
Mathematical context. This line imports or specializes earlier theories, making their carrier sets, functions, modules, and theorems available to the current file.

\noindent\textbf{Role in the proof.}
Use in this chapter. It fixes the proof context, imported mathematical language, or quantified module parameters for the following statements.

\subsubsection*{Definitions, Carrier Sets, Operations, Games, and Procedures}
\begin{formaldefinition}[EasyCrypt line 16]
\noindent\textbf{EasyCrypt name.}
\begin{lstlisting}[language=EasyCrypt,basicstyle=\ttfamily\scriptsize]
signing_transcript_record
\end{lstlisting}
\noindent\textbf{Checked EasyCrypt statement.}
\begin{lstlisting}[language=EasyCrypt,basicstyle=\ttfamily\scriptsize]
type signing_transcript_record = message * context * signature * transcript.
\end{lstlisting}
\noindent\textbf{Formal mathematical reading.}
Mathematical definition. This is a definitional type abbreviation.  The exact displayed EasyCrypt statement gives the right-hand side; later proof steps may unfold it without adding an assumption.

\noindent\textbf{Role in the proof.}
Use in this chapter. It defines transcript/extraction data for the Fiat-Shamir forking and special-soundness argument.

\end{formaldefinition}

\begin{formaldefinition}[EasyCrypt line 17]
\noindent\textbf{EasyCrypt name.}
\begin{lstlisting}[language=EasyCrypt,basicstyle=\ttfamily\scriptsize]
signing_attempt_sample
\end{lstlisting}
\noindent\textbf{Checked EasyCrypt statement.}
\begin{lstlisting}[language=EasyCrypt,basicstyle=\ttfamily\scriptsize]
type signing_attempt_sample = polyvecl * polyveck * polyveck.
\end{lstlisting}
\noindent\textbf{Formal mathematical reading.}
Mathematical definition. This is a definitional type abbreviation.  The exact displayed EasyCrypt statement gives the right-hand side; later proof steps may unfold it without adding an assumption.

\noindent\textbf{Role in the proof.}
Use in this chapter. It contributes a sampler or sampler-derived object to the distributional model.

\end{formaldefinition}

\begin{formaldefinition}[EasyCrypt line 18]
\noindent\textbf{EasyCrypt name.}
\begin{lstlisting}[language=EasyCrypt,basicstyle=\ttfamily\scriptsize]
signing_attempt_state
\end{lstlisting}
\noindent\textbf{Checked EasyCrypt statement.}
\begin{lstlisting}[language=EasyCrypt,basicstyle=\ttfamily\scriptsize]
type signing_attempt_state =
  random_coins * signing_attempt_sample * polyveck * poly * challenge *
  polyvecl * polyveck * hint_t * signature.
\end{lstlisting}
\noindent\textbf{Formal mathematical reading.}
Mathematical definition. This is a definitional type abbreviation.  The exact displayed EasyCrypt statement gives the right-hand side; later proof steps may unfold it without adding an assumption.

\noindent\textbf{Role in the proof.}
Use in this chapter. It is part of the exact formal vocabulary used by subsequent lemmas and games.

\end{formaldefinition}

\begin{formaldefinition}[EasyCrypt line 22]
\noindent\textbf{EasyCrypt name.}
\begin{lstlisting}[language=EasyCrypt,basicstyle=\ttfamily\scriptsize]
signing_attempt_state_coins
\end{lstlisting}
\noindent\textbf{Checked EasyCrypt statement.}
\begin{lstlisting}[language=EasyCrypt,basicstyle=\ttfamily\scriptsize]
op signing_attempt_state_coins
   (st : signing_attempt_state) : random_coins =
  st.`1.
\end{lstlisting}
\noindent\textbf{Formal mathematical reading.}
Mathematical definition. This is a total EasyCrypt operation.  The exact displayed statement gives the domain, codomain, and defining equation when present.

\noindent\textbf{Role in the proof.}
Use in this chapter. It contributes a sampler or sampler-derived object to the distributional model.

\end{formaldefinition}

\begin{formaldefinition}[EasyCrypt line 25]
\noindent\textbf{EasyCrypt name.}
\begin{lstlisting}[language=EasyCrypt,basicstyle=\ttfamily\scriptsize]
signing_attempt_state_sample
\end{lstlisting}
\noindent\textbf{Checked EasyCrypt statement.}
\begin{lstlisting}[language=EasyCrypt,basicstyle=\ttfamily\scriptsize]
op signing_attempt_state_sample
   (st : signing_attempt_state) : signing_attempt_sample =
  st.`2.
\end{lstlisting}
\noindent\textbf{Formal mathematical reading.}
Mathematical definition. This is a total EasyCrypt operation.  The exact displayed statement gives the domain, codomain, and defining equation when present.

\noindent\textbf{Role in the proof.}
Use in this chapter. It contributes a sampler or sampler-derived object to the distributional model.

\end{formaldefinition}

\begin{formaldefinition}[EasyCrypt line 28]
\noindent\textbf{EasyCrypt name.}
\begin{lstlisting}[language=EasyCrypt,basicstyle=\ttfamily\scriptsize]
signing_attempt_sample_y1
\end{lstlisting}
\noindent\textbf{Checked EasyCrypt statement.}
\begin{lstlisting}[language=EasyCrypt,basicstyle=\ttfamily\scriptsize]
op signing_attempt_sample_y1 (s : signing_attempt_sample) : polyvecl =
  s.`1.
\end{lstlisting}
\noindent\textbf{Formal mathematical reading.}
Mathematical definition. This is a total EasyCrypt operation.  The exact displayed statement gives the domain, codomain, and defining equation when present.

\noindent\textbf{Role in the proof.}
Use in this chapter. It contributes a sampler or sampler-derived object to the distributional model.

\end{formaldefinition}

\begin{formaldefinition}[EasyCrypt line 30]
\noindent\textbf{EasyCrypt name.}
\begin{lstlisting}[language=EasyCrypt,basicstyle=\ttfamily\scriptsize]
signing_attempt_sample_y2
\end{lstlisting}
\noindent\textbf{Checked EasyCrypt statement.}
\begin{lstlisting}[language=EasyCrypt,basicstyle=\ttfamily\scriptsize]
op signing_attempt_sample_y2 (s : signing_attempt_sample) : polyveck =
  s.`2.
\end{lstlisting}
\noindent\textbf{Formal mathematical reading.}
Mathematical definition. This is a total EasyCrypt operation.  The exact displayed statement gives the domain, codomain, and defining equation when present.

\noindent\textbf{Role in the proof.}
Use in this chapter. It contributes a sampler or sampler-derived object to the distributional model.

\end{formaldefinition}

\begin{formaldefinition}[EasyCrypt line 32]
\noindent\textbf{EasyCrypt name.}
\begin{lstlisting}[language=EasyCrypt,basicstyle=\ttfamily\scriptsize]
signing_attempt_sample_commitment_raw
\end{lstlisting}
\noindent\textbf{Checked EasyCrypt statement.}
\begin{lstlisting}[language=EasyCrypt,basicstyle=\ttfamily\scriptsize]
op signing_attempt_sample_commitment_raw
   (s : signing_attempt_sample) : polyveck =
  s.`3.
\end{lstlisting}
\noindent\textbf{Formal mathematical reading.}
Mathematical definition. This is a total EasyCrypt operation.  The exact displayed statement gives the domain, codomain, and defining equation when present.

\noindent\textbf{Role in the proof.}
Use in this chapter. It contributes a sampler or sampler-derived object to the distributional model.

\end{formaldefinition}

\begin{formaldefinition}[EasyCrypt line 35]
\noindent\textbf{EasyCrypt name.}
\begin{lstlisting}[language=EasyCrypt,basicstyle=\ttfamily\scriptsize]
signing_attempt_state_sampled_y1
\end{lstlisting}
\noindent\textbf{Checked EasyCrypt statement.}
\begin{lstlisting}[language=EasyCrypt,basicstyle=\ttfamily\scriptsize]
op signing_attempt_state_sampled_y1
   (st : signing_attempt_state) : polyvecl =
  signing_attempt_sample_y1 (signing_attempt_state_sample st).
\end{lstlisting}
\noindent\textbf{Formal mathematical reading.}
Mathematical definition. This is a total EasyCrypt operation.  The exact displayed statement gives the domain, codomain, and defining equation when present.

\noindent\textbf{Role in the proof.}
Use in this chapter. It contributes a sampler or sampler-derived object to the distributional model.

\end{formaldefinition}

\begin{formaldefinition}[EasyCrypt line 38]
\noindent\textbf{EasyCrypt name.}
\begin{lstlisting}[language=EasyCrypt,basicstyle=\ttfamily\scriptsize]
signing_attempt_state_sampled_y2
\end{lstlisting}
\noindent\textbf{Checked EasyCrypt statement.}
\begin{lstlisting}[language=EasyCrypt,basicstyle=\ttfamily\scriptsize]
op signing_attempt_state_sampled_y2
   (st : signing_attempt_state) : polyveck =
  signing_attempt_sample_y2 (signing_attempt_state_sample st).
\end{lstlisting}
\noindent\textbf{Formal mathematical reading.}
Mathematical definition. This is a total EasyCrypt operation.  The exact displayed statement gives the domain, codomain, and defining equation when present.

\noindent\textbf{Role in the proof.}
Use in this chapter. It contributes a sampler or sampler-derived object to the distributional model.

\end{formaldefinition}

\begin{formaldefinition}[EasyCrypt line 41]
\noindent\textbf{EasyCrypt name.}
\begin{lstlisting}[language=EasyCrypt,basicstyle=\ttfamily\scriptsize]
signing_attempt_state_sampled_y
\end{lstlisting}
\noindent\textbf{Checked EasyCrypt statement.}
\begin{lstlisting}[language=EasyCrypt,basicstyle=\ttfamily\scriptsize]
op signing_attempt_state_sampled_y
   (st : signing_attempt_state) : polyveck =
  signing_attempt_sample_commitment_raw (signing_attempt_state_sample st).
\end{lstlisting}
\noindent\textbf{Formal mathematical reading.}
Mathematical definition. This is a total EasyCrypt operation.  The exact displayed statement gives the domain, codomain, and defining equation when present.

\noindent\textbf{Role in the proof.}
Use in this chapter. It contributes a sampler or sampler-derived object to the distributional model.

\end{formaldefinition}

\begin{formaldefinition}[EasyCrypt line 44]
\noindent\textbf{EasyCrypt name.}
\begin{lstlisting}[language=EasyCrypt,basicstyle=\ttfamily\scriptsize]
signing_attempt_state_highbits
\end{lstlisting}
\noindent\textbf{Checked EasyCrypt statement.}
\begin{lstlisting}[language=EasyCrypt,basicstyle=\ttfamily\scriptsize]
op signing_attempt_state_highbits
   (st : signing_attempt_state) : polyveck =
  st.`3.
\end{lstlisting}
\noindent\textbf{Formal mathematical reading.}
Mathematical definition. This is a total EasyCrypt operation.  The exact displayed statement gives the domain, codomain, and defining equation when present.

\noindent\textbf{Role in the proof.}
Use in this chapter. It is part of the exact formal vocabulary used by subsequent lemmas and games.

\end{formaldefinition}

\begin{formaldefinition}[EasyCrypt line 47]
\noindent\textbf{EasyCrypt name.}
\begin{lstlisting}[language=EasyCrypt,basicstyle=\ttfamily\scriptsize]
signing_attempt_state_lowbits
\end{lstlisting}
\noindent\textbf{Checked EasyCrypt statement.}
\begin{lstlisting}[language=EasyCrypt,basicstyle=\ttfamily\scriptsize]
op signing_attempt_state_lowbits
   (st : signing_attempt_state) : poly =
  st.`4.
\end{lstlisting}
\noindent\textbf{Formal mathematical reading.}
Mathematical definition. This is a total EasyCrypt operation.  The exact displayed statement gives the domain, codomain, and defining equation when present.

\noindent\textbf{Role in the proof.}
Use in this chapter. It is part of the exact formal vocabulary used by subsequent lemmas and games.

\end{formaldefinition}

\begin{formaldefinition}[EasyCrypt line 50]
\noindent\textbf{EasyCrypt name.}
\begin{lstlisting}[language=EasyCrypt,basicstyle=\ttfamily\scriptsize]
signing_attempt_state_challenge
\end{lstlisting}
\noindent\textbf{Checked EasyCrypt statement.}
\begin{lstlisting}[language=EasyCrypt,basicstyle=\ttfamily\scriptsize]
op signing_attempt_state_challenge
   (st : signing_attempt_state) : challenge =
  st.`5.
\end{lstlisting}
\noindent\textbf{Formal mathematical reading.}
Mathematical definition. This is a total EasyCrypt operation.  The exact displayed statement gives the domain, codomain, and defining equation when present.

\noindent\textbf{Role in the proof.}
Use in this chapter. It is part of the exact formal vocabulary used by subsequent lemmas and games.

\end{formaldefinition}

\begin{formaldefinition}[EasyCrypt line 53]
\noindent\textbf{EasyCrypt name.}
\begin{lstlisting}[language=EasyCrypt,basicstyle=\ttfamily\scriptsize]
signing_attempt_state_response
\end{lstlisting}
\noindent\textbf{Checked EasyCrypt statement.}
\begin{lstlisting}[language=EasyCrypt,basicstyle=\ttfamily\scriptsize]
op signing_attempt_state_response
   (st : signing_attempt_state) : polyvecl =
  st.`6.
\end{lstlisting}
\noindent\textbf{Formal mathematical reading.}
Mathematical definition. This is a total EasyCrypt operation.  The exact displayed statement gives the domain, codomain, and defining equation when present.

\noindent\textbf{Role in the proof.}
Use in this chapter. It is part of the exact formal vocabulary used by subsequent lemmas and games.

\end{formaldefinition}

\begin{formaldefinition}[EasyCrypt line 56]
\noindent\textbf{EasyCrypt name.}
\begin{lstlisting}[language=EasyCrypt,basicstyle=\ttfamily\scriptsize]
signing_attempt_state_response_aux
\end{lstlisting}
\noindent\textbf{Checked EasyCrypt statement.}
\begin{lstlisting}[language=EasyCrypt,basicstyle=\ttfamily\scriptsize]
op signing_attempt_state_response_aux
   (st : signing_attempt_state) : polyveck =
  st.`7.
\end{lstlisting}
\noindent\textbf{Formal mathematical reading.}
Mathematical definition. This is a total EasyCrypt operation.  The exact displayed statement gives the domain, codomain, and defining equation when present.

\noindent\textbf{Role in the proof.}
Use in this chapter. It is part of the exact formal vocabulary used by subsequent lemmas and games.

\end{formaldefinition}

\begin{formaldefinition}[EasyCrypt line 59]
\noindent\textbf{EasyCrypt name.}
\begin{lstlisting}[language=EasyCrypt,basicstyle=\ttfamily\scriptsize]
signing_attempt_state_hint
\end{lstlisting}
\noindent\textbf{Checked EasyCrypt statement.}
\begin{lstlisting}[language=EasyCrypt,basicstyle=\ttfamily\scriptsize]
op signing_attempt_state_hint
   (st : signing_attempt_state) : hint_t =
  st.`8.
\end{lstlisting}
\noindent\textbf{Formal mathematical reading.}
Mathematical definition. This is a total EasyCrypt operation.  The exact displayed statement gives the domain, codomain, and defining equation when present.

\noindent\textbf{Role in the proof.}
Use in this chapter. It is part of the exact formal vocabulary used by subsequent lemmas and games.

\end{formaldefinition}

\begin{formaldefinition}[EasyCrypt line 62]
\noindent\textbf{EasyCrypt name.}
\begin{lstlisting}[language=EasyCrypt,basicstyle=\ttfamily\scriptsize]
signing_attempt_state_signature
\end{lstlisting}
\noindent\textbf{Checked EasyCrypt statement.}
\begin{lstlisting}[language=EasyCrypt,basicstyle=\ttfamily\scriptsize]
op signing_attempt_state_signature
   (st : signing_attempt_state) : signature =
  st.`9.
\end{lstlisting}
\noindent\textbf{Formal mathematical reading.}
Mathematical definition. This is a total EasyCrypt operation.  The exact displayed statement gives the domain, codomain, and defining equation when present.

\noindent\textbf{Role in the proof.}
Use in this chapter. It is part of the exact formal vocabulary used by subsequent lemmas and games.

\end{formaldefinition}

\begin{formaldefinition}[EasyCrypt line 65]
\noindent\textbf{EasyCrypt name.}
\begin{lstlisting}[language=EasyCrypt,basicstyle=\ttfamily\scriptsize]
signature_rejection_branch_byte
\end{lstlisting}
\noindent\textbf{Checked EasyCrypt statement.}
\begin{lstlisting}[language=EasyCrypt,basicstyle=\ttfamily\scriptsize]
op signature_rejection_branch_byte (sig : signature) : int =
  sig.`6.
\end{lstlisting}
\noindent\textbf{Formal mathematical reading.}
Mathematical definition. This is a total EasyCrypt operation.  The exact displayed statement gives the domain, codomain, and defining equation when present.

\noindent\textbf{Role in the proof.}
Use in this chapter. It names a bad event or rejection condition whose probability must be zero or explicitly bounded.

\end{formaldefinition}

\begin{formaldefinition}[EasyCrypt line 67]
\noindent\textbf{EasyCrypt name.}
\begin{lstlisting}[language=EasyCrypt,basicstyle=\ttfamily\scriptsize]
signing_attempt_state_reject2_branch_byte
\end{lstlisting}
\noindent\textbf{Checked EasyCrypt statement.}
\begin{lstlisting}[language=EasyCrypt,basicstyle=\ttfamily\scriptsize]
op signing_attempt_state_reject2_branch_byte
   (st : signing_attempt_state) : int =
  signature_rejection_branch_byte (signing_attempt_state_signature st).
\end{lstlisting}
\noindent\textbf{Formal mathematical reading.}
Mathematical definition. This is a total EasyCrypt operation.  The exact displayed statement gives the domain, codomain, and defining equation when present.

\noindent\textbf{Role in the proof.}
Use in this chapter. It is part of the exact formal vocabulary used by subsequent lemmas and games.

\end{formaldefinition}

\begin{formaldefinition}[EasyCrypt line 70]
\noindent\textbf{EasyCrypt name.}
\begin{lstlisting}[language=EasyCrypt,basicstyle=\ttfamily\scriptsize]
signing_attempt_state_reject2_branch_bit
\end{lstlisting}
\noindent\textbf{Checked EasyCrypt statement.}
\begin{lstlisting}[language=EasyCrypt,basicstyle=\ttfamily\scriptsize]
op signing_attempt_state_reject2_branch_bit
   (st : signing_attempt_state) : bool =
  (signing_attempt_state_reject2_branch_byte st %/ 2) %% 2 = 1.
\end{lstlisting}
\noindent\textbf{Formal mathematical reading.}
Mathematical definition. This is a Boolean predicate.  The exact displayed EasyCrypt statement gives its arguments; mathematically it is an event, invariant, support condition, or well-formedness predicate over those arguments.

\noindent\textbf{Role in the proof.}
Use in this chapter. It is part of the exact formal vocabulary used by subsequent lemmas and games.

\end{formaldefinition}

\begin{formaldefinition}[EasyCrypt line 73]
\noindent\textbf{EasyCrypt name.}
\begin{lstlisting}[language=EasyCrypt,basicstyle=\ttfamily\scriptsize]
polyvecl_double
\end{lstlisting}
\noindent\textbf{Checked EasyCrypt statement.}
\begin{lstlisting}[language=EasyCrypt,basicstyle=\ttfamily\scriptsize]
op polyvecl_double (md : mode) (xs : polyvecl) : polyvecl =
  mkseq
    (fun i => poly_double (nth poly_zero xs i))
    (mode_l md).
\end{lstlisting}
\noindent\textbf{Formal mathematical reading.}
Mathematical definition. This is a total EasyCrypt operation.  The exact displayed statement gives the domain, codomain, and defining equation when present.

\noindent\textbf{Role in the proof.}
Use in this chapter. It is part of the exact formal vocabulary used by subsequent lemmas and games.

\end{formaldefinition}

\begin{formaldefinition}[EasyCrypt line 77]
\noindent\textbf{EasyCrypt name.}
\begin{lstlisting}[language=EasyCrypt,basicstyle=\ttfamily\scriptsize]
signing_attempt_state_reject2_sampled_y1
\end{lstlisting}
\noindent\textbf{Checked EasyCrypt statement.}
\begin{lstlisting}[language=EasyCrypt,basicstyle=\ttfamily\scriptsize]
op signing_attempt_state_reject2_sampled_y1
   (st : signing_attempt_state) : polyvecl =
  signing_attempt_state_sampled_y1 st.
\end{lstlisting}
\noindent\textbf{Formal mathematical reading.}
Mathematical definition. This is a total EasyCrypt operation.  The exact displayed statement gives the domain, codomain, and defining equation when present.

\noindent\textbf{Role in the proof.}
Use in this chapter. It contributes a sampler or sampler-derived object to the distributional model.

\end{formaldefinition}

\begin{formaldefinition}[EasyCrypt line 80]
\noindent\textbf{EasyCrypt name.}
\begin{lstlisting}[language=EasyCrypt,basicstyle=\ttfamily\scriptsize]
signing_attempt_state_reject2_sampled_y2
\end{lstlisting}
\noindent\textbf{Checked EasyCrypt statement.}
\begin{lstlisting}[language=EasyCrypt,basicstyle=\ttfamily\scriptsize]
op signing_attempt_state_reject2_sampled_y2
   (st : signing_attempt_state) : polyveck =
  signing_attempt_state_sampled_y2 st.
\end{lstlisting}
\noindent\textbf{Formal mathematical reading.}
Mathematical definition. This is a total EasyCrypt operation.  The exact displayed statement gives the domain, codomain, and defining equation when present.

\noindent\textbf{Role in the proof.}
Use in this chapter. It contributes a sampler or sampler-derived object to the distributional model.

\end{formaldefinition}

\begin{formaldefinition}[EasyCrypt line 83]
\noindent\textbf{EasyCrypt name.}
\begin{lstlisting}[language=EasyCrypt,basicstyle=\ttfamily\scriptsize]
signing_attempt_state_reject2_balance_response
\end{lstlisting}
\noindent\textbf{Checked EasyCrypt statement.}
\begin{lstlisting}[language=EasyCrypt,basicstyle=\ttfamily\scriptsize]
op signing_attempt_state_reject2_balance_response
   (md : mode) (st : signing_attempt_state) : polyvecl =
  polyvecl_sub md
    (polyvecl_double md (signing_attempt_state_response st))
    (signing_attempt_state_reject2_sampled_y1 st).
\end{lstlisting}
\noindent\textbf{Formal mathematical reading.}
Mathematical definition. This is a total EasyCrypt operation.  The exact displayed statement gives the domain, codomain, and defining equation when present.

\noindent\textbf{Role in the proof.}
Use in this chapter. It is part of the exact formal vocabulary used by subsequent lemmas and games.

\end{formaldefinition}

\begin{formaldefinition}[EasyCrypt line 88]
\noindent\textbf{EasyCrypt name.}
\begin{lstlisting}[language=EasyCrypt,basicstyle=\ttfamily\scriptsize]
signing_attempt_state_reject2_balance_aux
\end{lstlisting}
\noindent\textbf{Checked EasyCrypt statement.}
\begin{lstlisting}[language=EasyCrypt,basicstyle=\ttfamily\scriptsize]
op signing_attempt_state_reject2_balance_aux
   (md : mode) (st : signing_attempt_state) : polyveck =
  polyveck_sub md
    (polyveck_double md (signing_attempt_state_response_aux st))
    (signing_attempt_state_reject2_sampled_y2 st).
\end{lstlisting}
\noindent\textbf{Formal mathematical reading.}
Mathematical definition. This is a total EasyCrypt operation.  The exact displayed statement gives the domain, codomain, and defining equation when present.

\noindent\textbf{Role in the proof.}
Use in this chapter. It is part of the exact formal vocabulary used by subsequent lemmas and games.

\end{formaldefinition}

\begin{formaldefinition}[EasyCrypt line 93]
\noindent\textbf{EasyCrypt name.}
\begin{lstlisting}[language=EasyCrypt,basicstyle=\ttfamily\scriptsize]
signing_attempt_state_lowbits_carrier
\end{lstlisting}
\noindent\textbf{Checked EasyCrypt statement.}
\begin{lstlisting}[language=EasyCrypt,basicstyle=\ttfamily\scriptsize]
op signing_attempt_state_lowbits_carrier
   (st : signing_attempt_state) : poly =
  nth poly_zero (signing_attempt_state_sampled_y st) 0.
\end{lstlisting}
\noindent\textbf{Formal mathematical reading.}
Mathematical definition. This is a total EasyCrypt operation.  The exact displayed statement gives the domain, codomain, and defining equation when present.

\noindent\textbf{Role in the proof.}
Use in this chapter. It is part of the exact formal vocabulary used by subsequent lemmas and games.

\end{formaldefinition}

\begin{formaldefinition}[EasyCrypt line 96]
\noindent\textbf{EasyCrypt name.}
\begin{lstlisting}[language=EasyCrypt,basicstyle=\ttfamily\scriptsize]
signing_attempt_state_token
\end{lstlisting}
\noindent\textbf{Checked EasyCrypt statement.}
\begin{lstlisting}[language=EasyCrypt,basicstyle=\ttfamily\scriptsize]
op signing_attempt_state_token
   (st : signing_attempt_state) : sig_token =
  ( signing_attempt_state_response st,
    signing_attempt_state_response_aux st,
    signing_attempt_state_hint st).
\end{lstlisting}
\noindent\textbf{Formal mathematical reading.}
Mathematical definition. This is a total EasyCrypt operation.  The exact displayed statement gives the domain, codomain, and defining equation when present.

\noindent\textbf{Role in the proof.}
Use in this chapter. It names a well-formedness, support, or validity predicate needed by later preservation lemmas.

\end{formaldefinition}

\begin{formaldefinition}[EasyCrypt line 102]
\noindent\textbf{EasyCrypt name.}
\begin{lstlisting}[language=EasyCrypt,basicstyle=\ttfamily\scriptsize]
signing_attempt_state_programmed_lowbits
\end{lstlisting}
\noindent\textbf{Checked EasyCrypt statement.}
\begin{lstlisting}[language=EasyCrypt,basicstyle=\ttfamily\scriptsize]
op signing_attempt_state_programmed_lowbits
   (st : signing_attempt_state) : poly =
  mkseq
    (fun i =>
      if i = 0 then signing_entropy_token_of_coins
        (signing_attempt_state_coins st)
      else poly_coeff (signing_attempt_state_lowbits_carrier st) i)
    n.
\end{lstlisting}
\noindent\textbf{Formal mathematical reading.}
Mathematical definition. This is a total EasyCrypt operation.  The exact displayed statement gives the domain, codomain, and defining equation when present.

\noindent\textbf{Role in the proof.}
Use in this chapter. It is part of the exact formal vocabulary used by subsequent lemmas and games.

\end{formaldefinition}

\begin{formaldefinition}[EasyCrypt line 111]
\noindent\textbf{EasyCrypt name.}
\begin{lstlisting}[language=EasyCrypt,basicstyle=\ttfamily\scriptsize]
signing_attempt_state_lowbits_consistent
\end{lstlisting}
\noindent\textbf{Checked EasyCrypt statement.}
\begin{lstlisting}[language=EasyCrypt,basicstyle=\ttfamily\scriptsize]
op signing_attempt_state_lowbits_consistent
   (st : signing_attempt_state) : bool =
  signing_attempt_state_lowbits st =
    signing_attempt_state_programmed_lowbits st.
\end{lstlisting}
\noindent\textbf{Formal mathematical reading.}
Mathematical definition. This is a Boolean predicate.  The exact displayed EasyCrypt statement gives its arguments; mathematically it is an event, invariant, support condition, or well-formedness predicate over those arguments.

\noindent\textbf{Role in the proof.}
Use in this chapter. It is part of the exact formal vocabulary used by subsequent lemmas and games.

\end{formaldefinition}

\begin{formaldefinition}[EasyCrypt line 116]
\noindent\textbf{EasyCrypt name.}
\begin{lstlisting}[language=EasyCrypt,basicstyle=\ttfamily\scriptsize]
signing_attempt_state_programmed_challenge
\end{lstlisting}
\noindent\textbf{Checked EasyCrypt statement.}
\begin{lstlisting}[language=EasyCrypt,basicstyle=\ttfamily\scriptsize]
op signing_attempt_state_programmed_challenge
   (md : mode) (pk : pkey) (m : message) (ctx : context)
   (st : signing_attempt_state) : challenge =
  challenge_hash md
    (signing_attempt_state_response_aux st)
    (signing_attempt_state_lowbits st)
    (message_hash pk ctx m).
\end{lstlisting}
\noindent\textbf{Formal mathematical reading.}
Mathematical definition. This is a total EasyCrypt operation.  The exact displayed statement gives the domain, codomain, and defining equation when present.

\noindent\textbf{Role in the proof.}
Use in this chapter. It is part of the exact formal vocabulary used by subsequent lemmas and games.

\end{formaldefinition}

\begin{formaldefinition}[EasyCrypt line 124]
\noindent\textbf{EasyCrypt name.}
\begin{lstlisting}[language=EasyCrypt,basicstyle=\ttfamily\scriptsize]
signing_attempt_state_challenge_consistent
\end{lstlisting}
\noindent\textbf{Checked EasyCrypt statement.}
\begin{lstlisting}[language=EasyCrypt,basicstyle=\ttfamily\scriptsize]
op signing_attempt_state_challenge_consistent
   (md : mode) (pk : pkey) (m : message) (ctx : context)
   (st : signing_attempt_state) : bool =
  signing_attempt_state_challenge st =
    signing_attempt_state_programmed_challenge md pk m ctx st.
\end{lstlisting}
\noindent\textbf{Formal mathematical reading.}
Mathematical definition. This is a Boolean predicate.  The exact displayed EasyCrypt statement gives its arguments; mathematically it is an event, invariant, support condition, or well-formedness predicate over those arguments.

\noindent\textbf{Role in the proof.}
Use in this chapter. It is part of the exact formal vocabulary used by subsequent lemmas and games.

\end{formaldefinition}

\begin{formaldefinition}[EasyCrypt line 130]
\noindent\textbf{EasyCrypt name.}
\begin{lstlisting}[language=EasyCrypt,basicstyle=\ttfamily\scriptsize]
signing_attempt_state_reconstructed_highbits
\end{lstlisting}
\noindent\textbf{Checked EasyCrypt statement.}
\begin{lstlisting}[language=EasyCrypt,basicstyle=\ttfamily\scriptsize]
op signing_attempt_state_reconstructed_highbits
   (md : mode) (pk : pkey) (st : signing_attempt_state) : polyveck =
  reconstructed_highbits md
    (signing_attempt_state_response_aux st)
    (public_key_challenge_term md pk (signing_attempt_state_challenge st)).
\end{lstlisting}
\noindent\textbf{Formal mathematical reading.}
Mathematical definition. This is a total EasyCrypt operation.  The exact displayed statement gives the domain, codomain, and defining equation when present.

\noindent\textbf{Role in the proof.}
Use in this chapter. It is part of the exact formal vocabulary used by subsequent lemmas and games.

\end{formaldefinition}

\begin{formaldefinition}[EasyCrypt line 136]
\noindent\textbf{EasyCrypt name.}
\begin{lstlisting}[language=EasyCrypt,basicstyle=\ttfamily\scriptsize]
signing_attempt_state_public_equation_consistent
\end{lstlisting}
\noindent\textbf{Checked EasyCrypt statement.}
\begin{lstlisting}[language=EasyCrypt,basicstyle=\ttfamily\scriptsize]
op signing_attempt_state_public_equation_consistent
   (md : mode) (pk : pkey) (st : signing_attempt_state) : bool =
  signing_attempt_state_highbits st =
    signing_attempt_state_reconstructed_highbits md pk st.
\end{lstlisting}
\noindent\textbf{Formal mathematical reading.}
Mathematical definition. This is a Boolean predicate.  The exact displayed EasyCrypt statement gives its arguments; mathematically it is an event, invariant, support condition, or well-formedness predicate over those arguments.

\noindent\textbf{Role in the proof.}
Use in this chapter. It is part of the exact formal vocabulary used by subsequent lemmas and games.

\end{formaldefinition}

\begin{formaldefinition}[EasyCrypt line 141]
\noindent\textbf{EasyCrypt name.}
\begin{lstlisting}[language=EasyCrypt,basicstyle=\ttfamily\scriptsize]
signing_attempt_state_signature_of_fields
\end{lstlisting}
\noindent\textbf{Checked EasyCrypt statement.}
\begin{lstlisting}[language=EasyCrypt,basicstyle=\ttfamily\scriptsize]
op signing_attempt_state_signature_of_fields
   (pk : pkey) (m : message) (ctx : context)
   (st : signing_attempt_state) : signature =
  ( signing_attempt_state_highbits st,
    signing_attempt_state_lowbits st,
    message_hash pk ctx m,
    signing_attempt_state_challenge st,
    signing_attempt_state_token st,
    0,
    0).
\end{lstlisting}
\noindent\textbf{Formal mathematical reading.}
Mathematical definition. This is a total EasyCrypt operation.  The exact displayed statement gives the domain, codomain, and defining equation when present.

\noindent\textbf{Role in the proof.}
Use in this chapter. It is part of the exact formal vocabulary used by subsequent lemmas and games.

\end{formaldefinition}

\begin{formaldefinition}[EasyCrypt line 152]
\noindent\textbf{EasyCrypt name.}
\begin{lstlisting}[language=EasyCrypt,basicstyle=\ttfamily\scriptsize]
signing_attempt_state_fields_match_signature
\end{lstlisting}
\noindent\textbf{Checked EasyCrypt statement.}
\begin{lstlisting}[language=EasyCrypt,basicstyle=\ttfamily\scriptsize]
op signing_attempt_state_fields_match_signature
   (pk : pkey) (m : message) (ctx : context)
   (st : signing_attempt_state) : bool =
  signing_attempt_state_signature st =
    signing_attempt_state_signature_of_fields pk m ctx st.
\end{lstlisting}
\noindent\textbf{Formal mathematical reading.}
Mathematical definition. This is a Boolean predicate.  The exact displayed EasyCrypt statement gives its arguments; mathematically it is an event, invariant, support condition, or well-formedness predicate over those arguments.

\noindent\textbf{Role in the proof.}
Use in this chapter. It is part of the exact formal vocabulary used by subsequent lemmas and games.

\end{formaldefinition}

\begin{formaldefinition}[EasyCrypt line 158]
\noindent\textbf{EasyCrypt name.}
\begin{lstlisting}[language=EasyCrypt,basicstyle=\ttfamily\scriptsize]
signing_attempt_state_field_accepts
\end{lstlisting}
\noindent\textbf{Checked EasyCrypt statement.}
\begin{lstlisting}[language=EasyCrypt,basicstyle=\ttfamily\scriptsize]
op signing_attempt_state_field_accepts
   (md : mode) (pk : pkey) (m : message) (ctx : context)
   (st : signing_attempt_state) : bool =
  signing_attempt_state_lowbits_consistent st /\
  signing_attempt_state_challenge_consistent md pk m ctx st /\
  signing_attempt_state_public_equation_consistent md pk st /\
  signing_attempt_state_fields_match_signature pk m ctx st.
\end{lstlisting}
\noindent\textbf{Formal mathematical reading.}
Mathematical definition. This is a Boolean predicate.  The exact displayed EasyCrypt statement gives its arguments; mathematically it is an event, invariant, support condition, or well-formedness predicate over those arguments.

\noindent\textbf{Role in the proof.}
Use in this chapter. It is part of the exact formal vocabulary used by subsequent lemmas and games.

\end{formaldefinition}

\begin{formaldefinition}[EasyCrypt line 166]
\noindent\textbf{EasyCrypt name.}
\begin{lstlisting}[language=EasyCrypt,basicstyle=\ttfamily\scriptsize]
signing_attempt_state_of_coins
\end{lstlisting}
\noindent\textbf{Checked EasyCrypt statement.}
\begin{lstlisting}[language=EasyCrypt,basicstyle=\ttfamily\scriptsize]
op signing_attempt_state_of_coins
   (md : mode) (sk : skey) (m : message) (ctx : context)
   (coins : random_coins) : signing_attempt_state =
  ( coins,
    ( signing_sample_y1 md sk m ctx coins,
      signing_sample_y2 md sk m ctx coins,
      commitment_raw md sk m ctx coins),
    commitment_highbits md sk m ctx coins,
    commitment_lowbits md sk m ctx coins,
    commitment_challenge md sk m ctx coins,
    response_vector md sk m ctx coins,
    response_aux_vector md sk m ctx coins,
    hint_vector md sk m ctx coins,
    sign_internal md sk m ctx coins).
\end{lstlisting}
\noindent\textbf{Formal mathematical reading.}
Mathematical definition. This is a total EasyCrypt operation.  The exact displayed statement gives the domain, codomain, and defining equation when present.

\noindent\textbf{Role in the proof.}
Use in this chapter. It contributes a sampler or sampler-derived object to the distributional model.

\end{formaldefinition}

\begin{formaldefinition}[EasyCrypt line 181]
\noindent\textbf{EasyCrypt name.}
\begin{lstlisting}[language=EasyCrypt,basicstyle=\ttfamily\scriptsize]
signing_attempt_sample_of_pair
\end{lstlisting}
\noindent\textbf{Checked EasyCrypt statement.}
\begin{lstlisting}[language=EasyCrypt,basicstyle=\ttfamily\scriptsize]
op signing_attempt_sample_of_pair
   (md : mode) (sk : skey) (m : message) (ctx : context)
   (coins : random_coins) (sample : signing_sample_pair) :
   signing_attempt_sample =
  (signing_sample_pair_y1 sample,
   signing_sample_pair_y2 sample,
   commitment_raw md sk m ctx coins).
\end{lstlisting}
\noindent\textbf{Formal mathematical reading.}
Mathematical definition. This is a total EasyCrypt operation.  The exact displayed statement gives the domain, codomain, and defining equation when present.

\noindent\textbf{Role in the proof.}
Use in this chapter. It contributes a sampler or sampler-derived object to the distributional model.

\end{formaldefinition}

\begin{formaldefinition}[EasyCrypt line 189]
\noindent\textbf{EasyCrypt name.}
\begin{lstlisting}[language=EasyCrypt,basicstyle=\ttfamily\scriptsize]
signing_attempt_state_of_sample_pair
\end{lstlisting}
\noindent\textbf{Checked EasyCrypt statement.}
\begin{lstlisting}[language=EasyCrypt,basicstyle=\ttfamily\scriptsize]
op signing_attempt_state_of_sample_pair
   (md : mode) (sk : skey) (m : message) (ctx : context)
   (coins : random_coins) (sample : signing_sample_pair) :
   signing_attempt_state =
  ( coins,
    signing_attempt_sample_of_pair md sk m ctx coins sample,
    commitment_highbits md sk m ctx coins,
    commitment_lowbits md sk m ctx coins,
    commitment_challenge md sk m ctx coins,
    response_vector md sk m ctx coins,
    response_aux_vector md sk m ctx coins,
    hint_vector md sk m ctx coins,
    sign_internal md sk m ctx coins).
\end{lstlisting}
\noindent\textbf{Formal mathematical reading.}
Mathematical definition. This is a total EasyCrypt operation.  The exact displayed statement gives the domain, codomain, and defining equation when present.

\noindent\textbf{Role in the proof.}
Use in this chapter. It contributes a sampler or sampler-derived object to the distributional model.

\end{formaldefinition}

\begin{formaldefinition}[EasyCrypt line 203]
\noindent\textbf{EasyCrypt name.}
\begin{lstlisting}[language=EasyCrypt,basicstyle=\ttfamily\scriptsize]
signing_attempt_state_valid_signature_accepts
\end{lstlisting}
\noindent\textbf{Checked EasyCrypt statement.}
\begin{lstlisting}[language=EasyCrypt,basicstyle=\ttfamily\scriptsize]
op signing_attempt_state_valid_signature_accepts
   (md : mode) (st : signing_attempt_state) : bool =
  valid_signature md (signing_attempt_state_signature st).
\end{lstlisting}
\noindent\textbf{Formal mathematical reading.}
Mathematical definition. This is a Boolean predicate.  The exact displayed EasyCrypt statement gives its arguments; mathematically it is an event, invariant, support condition, or well-formedness predicate over those arguments.

\noindent\textbf{Role in the proof.}
Use in this chapter. It names a well-formedness, support, or validity predicate needed by later preservation lemmas.

\end{formaldefinition}

\begin{formaldefinition}[EasyCrypt line 207]
\noindent\textbf{EasyCrypt name.}
\begin{lstlisting}[language=EasyCrypt,basicstyle=\ttfamily\scriptsize]
signing_attempt_state_response_norm_accepts
\end{lstlisting}
\noindent\textbf{Checked EasyCrypt statement.}
\begin{lstlisting}[language=EasyCrypt,basicstyle=\ttfamily\scriptsize]
op signing_attempt_state_response_norm_accepts
   (md : mode) (st : signing_attempt_state) : bool =
  response_norm_ok md (signing_attempt_state_signature st).
\end{lstlisting}
\noindent\textbf{Formal mathematical reading.}
Mathematical definition. This is a Boolean predicate.  The exact displayed EasyCrypt statement gives its arguments; mathematically it is an event, invariant, support condition, or well-formedness predicate over those arguments.

\noindent\textbf{Role in the proof.}
Use in this chapter. It is part of the exact formal vocabulary used by subsequent lemmas and games.

\end{formaldefinition}

\begin{formaldefinition}[EasyCrypt line 211]
\noindent\textbf{EasyCrypt name.}
\begin{lstlisting}[language=EasyCrypt,basicstyle=\ttfamily\scriptsize]
signing_attempt_state_hint_norm_accepts
\end{lstlisting}
\noindent\textbf{Checked EasyCrypt statement.}
\begin{lstlisting}[language=EasyCrypt,basicstyle=\ttfamily\scriptsize]
op signing_attempt_state_hint_norm_accepts
   (md : mode) (st : signing_attempt_state) : bool =
  verify_norm_ok md (signing_attempt_state_signature st).
\end{lstlisting}
\noindent\textbf{Formal mathematical reading.}
Mathematical definition. This is a Boolean predicate.  The exact displayed EasyCrypt statement gives its arguments; mathematically it is an event, invariant, support condition, or well-formedness predicate over those arguments.

\noindent\textbf{Role in the proof.}
Use in this chapter. It is part of the exact formal vocabulary used by subsequent lemmas and games.

\end{formaldefinition}

\begin{formaldefinition}[EasyCrypt line 215]
\noindent\textbf{EasyCrypt name.}
\begin{lstlisting}[language=EasyCrypt,basicstyle=\ttfamily\scriptsize]
signing_attempt_state_accepts
\end{lstlisting}
\noindent\textbf{Checked EasyCrypt statement.}
\begin{lstlisting}[language=EasyCrypt,basicstyle=\ttfamily\scriptsize]
op signing_attempt_state_accepts
   (md : mode) (st : signing_attempt_state) : bool =
  signing_attempt_state_valid_signature_accepts md st /\
  signing_attempt_state_response_norm_accepts md st /\
  signing_attempt_state_hint_norm_accepts md st.
\end{lstlisting}
\noindent\textbf{Formal mathematical reading.}
Mathematical definition. This is a Boolean predicate.  The exact displayed EasyCrypt statement gives its arguments; mathematically it is an event, invariant, support condition, or well-formedness predicate over those arguments.

\noindent\textbf{Role in the proof.}
Use in this chapter. It is part of the exact formal vocabulary used by subsequent lemmas and games.

\end{formaldefinition}

\begin{formaldefinition}[EasyCrypt line 221]
\noindent\textbf{EasyCrypt name.}
\begin{lstlisting}[language=EasyCrypt,basicstyle=\ttfamily\scriptsize]
signing_attempt_reject1_bound
\end{lstlisting}
\noindent\textbf{Checked EasyCrypt statement.}
\begin{lstlisting}[language=EasyCrypt,basicstyle=\ttfamily\scriptsize]
op signing_attempt_reject1_bound (md : mode) : int =
  mode_b1sq md * mode_ln md * mode_ln md.
\end{lstlisting}
\noindent\textbf{Formal mathematical reading.}
Mathematical definition. This is a total EasyCrypt operation.  The exact displayed statement gives the domain, codomain, and defining equation when present.

\noindent\textbf{Role in the proof.}
Use in this chapter. It names a numerical term that appears in a game-hop or final security inequality.

\end{formaldefinition}

\begin{formaldefinition}[EasyCrypt line 224]
\noindent\textbf{EasyCrypt name.}
\begin{lstlisting}[language=EasyCrypt,basicstyle=\ttfamily\scriptsize]
signing_attempt_reject2_bound
\end{lstlisting}
\noindent\textbf{Checked EasyCrypt statement.}
\begin{lstlisting}[language=EasyCrypt,basicstyle=\ttfamily\scriptsize]
op signing_attempt_reject2_bound (md : mode) : int =
  mode_b0sq md * mode_ln md * mode_ln md.
\end{lstlisting}
\noindent\textbf{Formal mathematical reading.}
Mathematical definition. This is a total EasyCrypt operation.  The exact displayed statement gives the domain, codomain, and defining equation when present.

\noindent\textbf{Role in the proof.}
Use in this chapter. It names a numerical term that appears in a game-hop or final security inequality.

\end{formaldefinition}

\begin{formaldefinition}[EasyCrypt line 227]
\noindent\textbf{EasyCrypt name.}
\begin{lstlisting}[language=EasyCrypt,basicstyle=\ttfamily\scriptsize]
signing_attempt_reject2_balance_norm_sq
\end{lstlisting}
\noindent\textbf{Checked EasyCrypt statement.}
\begin{lstlisting}[language=EasyCrypt,basicstyle=\ttfamily\scriptsize]
op signing_attempt_reject2_balance_norm_sq
   (md : mode) (z : polyvecl) (zaux : polyveck)
   (y1 : polyvecl) (y2 : polyveck) : int =
  polyvecl_norm_sq md (polyvecl_sub md (polyvecl_double md z) y1) +
  polyveck_norm_sq md (polyveck_sub md (polyveck_double md zaux) y2).
\end{lstlisting}
\noindent\textbf{Formal mathematical reading.}
Mathematical definition. This is a total EasyCrypt operation.  The exact displayed statement gives the domain, codomain, and defining equation when present.

\noindent\textbf{Role in the proof.}
Use in this chapter. It is part of the exact formal vocabulary used by subsequent lemmas and games.

\end{formaldefinition}

\begin{formaldefinition}[EasyCrypt line 233]
\noindent\textbf{EasyCrypt name.}
\begin{lstlisting}[language=EasyCrypt,basicstyle=\ttfamily\scriptsize]
signing_attempt_reject2_concrete_aborts
\end{lstlisting}
\noindent\textbf{Checked EasyCrypt statement.}
\begin{lstlisting}[language=EasyCrypt,basicstyle=\ttfamily\scriptsize]
op signing_attempt_reject2_concrete_aborts
   (md : mode) (branch : bool) (z : polyvecl) (zaux : polyveck)
   (y1 : polyvecl) (y2 : polyveck) : bool =
  branch /\
  signing_attempt_reject2_balance_norm_sq md z zaux y1 y2 <
  signing_attempt_reject2_bound md.
\end{lstlisting}
\noindent\textbf{Formal mathematical reading.}
Mathematical definition. This is a Boolean predicate.  The exact displayed EasyCrypt statement gives its arguments; mathematically it is an event, invariant, support condition, or well-formedness predicate over those arguments.

\noindent\textbf{Role in the proof.}
Use in this chapter. It is part of the exact formal vocabulary used by subsequent lemmas and games.

\end{formaldefinition}

\begin{formaldefinition}[EasyCrypt line 240]
\noindent\textbf{EasyCrypt name.}
\begin{lstlisting}[language=EasyCrypt,basicstyle=\ttfamily\scriptsize]
signing_attempt_state_reject1_norm_sq
\end{lstlisting}
\noindent\textbf{Checked EasyCrypt statement.}
\begin{lstlisting}[language=EasyCrypt,basicstyle=\ttfamily\scriptsize]
op signing_attempt_state_reject1_norm_sq
   (md : mode) (st : signing_attempt_state) : int =
  response_norm_sq md (signing_attempt_state_signature st).
\end{lstlisting}
\noindent\textbf{Formal mathematical reading.}
Mathematical definition. This is a total EasyCrypt operation.  The exact displayed statement gives the domain, codomain, and defining equation when present.

\noindent\textbf{Role in the proof.}
Use in this chapter. It is part of the exact formal vocabulary used by subsequent lemmas and games.

\end{formaldefinition}

\begin{formaldefinition}[EasyCrypt line 244]
\noindent\textbf{EasyCrypt name.}
\begin{lstlisting}[language=EasyCrypt,basicstyle=\ttfamily\scriptsize]
signing_attempt_state_reject1_accepts
\end{lstlisting}
\noindent\textbf{Checked EasyCrypt statement.}
\begin{lstlisting}[language=EasyCrypt,basicstyle=\ttfamily\scriptsize]
op signing_attempt_state_reject1_accepts
   (md : mode) (st : signing_attempt_state) : bool =
  0 <= signing_attempt_state_reject1_norm_sq md st /\
  signing_attempt_state_reject1_norm_sq md st <=
    signing_attempt_reject1_bound md.
\end{lstlisting}
\noindent\textbf{Formal mathematical reading.}
Mathematical definition. This is a Boolean predicate.  The exact displayed EasyCrypt statement gives its arguments; mathematically it is an event, invariant, support condition, or well-formedness predicate over those arguments.

\noindent\textbf{Role in the proof.}
Use in this chapter. It is part of the exact formal vocabulary used by subsequent lemmas and games.

\end{formaldefinition}

\begin{formaldefinition}[EasyCrypt line 250]
\noindent\textbf{EasyCrypt name.}
\begin{lstlisting}[language=EasyCrypt,basicstyle=\ttfamily\scriptsize]
signing_attempt_state_balance_norm_sq
\end{lstlisting}
\noindent\textbf{Checked EasyCrypt statement.}
\begin{lstlisting}[language=EasyCrypt,basicstyle=\ttfamily\scriptsize]
op signing_attempt_state_balance_norm_sq
   (md : mode) (st : signing_attempt_state) : int =
  signing_attempt_reject2_balance_norm_sq md
    (signing_attempt_state_response st)
    (signing_attempt_state_response_aux st)
    (signing_attempt_state_reject2_sampled_y1 st)
    (signing_attempt_state_reject2_sampled_y2 st).
\end{lstlisting}
\noindent\textbf{Formal mathematical reading.}
Mathematical definition. This is a total EasyCrypt operation.  The exact displayed statement gives the domain, codomain, and defining equation when present.

\noindent\textbf{Role in the proof.}
Use in this chapter. It is part of the exact formal vocabulary used by subsequent lemmas and games.

\end{formaldefinition}

\begin{formaldefinition}[EasyCrypt line 258]
\noindent\textbf{EasyCrypt name.}
\begin{lstlisting}[language=EasyCrypt,basicstyle=\ttfamily\scriptsize]
signing_attempt_state_bimodal_reject2_branch
\end{lstlisting}
\noindent\textbf{Checked EasyCrypt statement.}
\begin{lstlisting}[language=EasyCrypt,basicstyle=\ttfamily\scriptsize]
op signing_attempt_state_bimodal_reject2_branch
   (st : signing_attempt_state) : bool =
  signing_attempt_state_reject2_branch_bit st.
\end{lstlisting}
\noindent\textbf{Formal mathematical reading.}
Mathematical definition. This is a Boolean predicate.  The exact displayed EasyCrypt statement gives its arguments; mathematically it is an event, invariant, support condition, or well-formedness predicate over those arguments.

\noindent\textbf{Role in the proof.}
Use in this chapter. It is part of the exact formal vocabulary used by subsequent lemmas and games.

\end{formaldefinition}

\begin{formaldefinition}[EasyCrypt line 262]
\noindent\textbf{EasyCrypt name.}
\begin{lstlisting}[language=EasyCrypt,basicstyle=\ttfamily\scriptsize]
signing_attempt_state_reject2_under_bound
\end{lstlisting}
\noindent\textbf{Checked EasyCrypt statement.}
\begin{lstlisting}[language=EasyCrypt,basicstyle=\ttfamily\scriptsize]
op signing_attempt_state_reject2_under_bound
   (md : mode) (st : signing_attempt_state) : bool =
  signing_attempt_state_balance_norm_sq md st <
  signing_attempt_reject2_bound md.
\end{lstlisting}
\noindent\textbf{Formal mathematical reading.}
Mathematical definition. This is a Boolean predicate.  The exact displayed EasyCrypt statement gives its arguments; mathematically it is an event, invariant, support condition, or well-formedness predicate over those arguments.

\noindent\textbf{Role in the proof.}
Use in this chapter. It names a numerical term that appears in a game-hop or final security inequality.

\end{formaldefinition}

\begin{formaldefinition}[EasyCrypt line 267]
\noindent\textbf{EasyCrypt name.}
\begin{lstlisting}[language=EasyCrypt,basicstyle=\ttfamily\scriptsize]
signing_attempt_state_reject2_aborts
\end{lstlisting}
\noindent\textbf{Checked EasyCrypt statement.}
\begin{lstlisting}[language=EasyCrypt,basicstyle=\ttfamily\scriptsize]
op signing_attempt_state_reject2_aborts
   (md : mode) (st : signing_attempt_state) : bool =
  signing_attempt_reject2_concrete_aborts md
    (signing_attempt_state_bimodal_reject2_branch st)
    (signing_attempt_state_response st)
    (signing_attempt_state_response_aux st)
    (signing_attempt_state_reject2_sampled_y1 st)
    (signing_attempt_state_reject2_sampled_y2 st).
\end{lstlisting}
\noindent\textbf{Formal mathematical reading.}
Mathematical definition. This is a Boolean predicate.  The exact displayed EasyCrypt statement gives its arguments; mathematically it is an event, invariant, support condition, or well-formedness predicate over those arguments.

\noindent\textbf{Role in the proof.}
Use in this chapter. It is part of the exact formal vocabulary used by subsequent lemmas and games.

\end{formaldefinition}

\begin{formaldefinition}[EasyCrypt line 276]
\noindent\textbf{EasyCrypt name.}
\begin{lstlisting}[language=EasyCrypt,basicstyle=\ttfamily\scriptsize]
signing_attempt_state_reject2_accepts
\end{lstlisting}
\noindent\textbf{Checked EasyCrypt statement.}
\begin{lstlisting}[language=EasyCrypt,basicstyle=\ttfamily\scriptsize]
op signing_attempt_state_reject2_accepts
   (md : mode) (st : signing_attempt_state) : bool =
  ! signing_attempt_state_reject2_aborts md st.
\end{lstlisting}
\noindent\textbf{Formal mathematical reading.}
Mathematical definition. This is a Boolean predicate.  The exact displayed EasyCrypt statement gives its arguments; mathematically it is an event, invariant, support condition, or well-formedness predicate over those arguments.

\noindent\textbf{Role in the proof.}
Use in this chapter. It is part of the exact formal vocabulary used by subsequent lemmas and games.

\end{formaldefinition}

\begin{formaldefinition}[EasyCrypt line 280]
\noindent\textbf{EasyCrypt name.}
\begin{lstlisting}[language=EasyCrypt,basicstyle=\ttfamily\scriptsize]
signing_attempt_state_pack_accepts
\end{lstlisting}
\noindent\textbf{Checked EasyCrypt statement.}
\begin{lstlisting}[language=EasyCrypt,basicstyle=\ttfamily\scriptsize]
op signing_attempt_state_pack_accepts
   (md : mode) (st : signing_attempt_state) : bool =
  signing_attempt_state_valid_signature_accepts md st /\
  0 < mode_crypto_bytes md.
\end{lstlisting}
\noindent\textbf{Formal mathematical reading.}
Mathematical definition. This is a Boolean predicate.  The exact displayed EasyCrypt statement gives its arguments; mathematically it is an event, invariant, support condition, or well-formedness predicate over those arguments.

\noindent\textbf{Role in the proof.}
Use in this chapter. It is part of the exact formal vocabulary used by subsequent lemmas and games.

\end{formaldefinition}

\begin{formaldefinition}[EasyCrypt line 285]
\noindent\textbf{EasyCrypt name.}
\begin{lstlisting}[language=EasyCrypt,basicstyle=\ttfamily\scriptsize]
signing_attempt_state_reference_rejection_accepts
\end{lstlisting}
\noindent\textbf{Checked EasyCrypt statement.}
\begin{lstlisting}[language=EasyCrypt,basicstyle=\ttfamily\scriptsize]
op signing_attempt_state_reference_rejection_accepts
   (md : mode) (st : signing_attempt_state) : bool =
  signing_attempt_state_reject1_accepts md st /\
  signing_attempt_state_reject2_accepts md st /\
  signing_attempt_state_pack_accepts md st /\
  signing_attempt_state_hint_norm_accepts md st.
\end{lstlisting}
\noindent\textbf{Formal mathematical reading.}
Mathematical definition. This is a Boolean predicate.  The exact displayed EasyCrypt statement gives its arguments; mathematically it is an event, invariant, support condition, or well-formedness predicate over those arguments.

\noindent\textbf{Role in the proof.}
Use in this chapter. It names a bad event or rejection condition whose probability must be zero or explicitly bounded.

\end{formaldefinition}

\begin{formaldefinition}[EasyCrypt line 292]
\noindent\textbf{EasyCrypt name.}
\begin{lstlisting}[language=EasyCrypt,basicstyle=\ttfamily\scriptsize]
signing_attempt_state_reference_rejection_aborts
\end{lstlisting}
\noindent\textbf{Checked EasyCrypt statement.}
\begin{lstlisting}[language=EasyCrypt,basicstyle=\ttfamily\scriptsize]
op signing_attempt_state_reference_rejection_aborts
   (md : mode) (st : signing_attempt_state) : bool =
  ! signing_attempt_state_reference_rejection_accepts md st.
\end{lstlisting}
\noindent\textbf{Formal mathematical reading.}
Mathematical definition. This is a Boolean predicate.  The exact displayed EasyCrypt statement gives its arguments; mathematically it is an event, invariant, support condition, or well-formedness predicate over those arguments.

\noindent\textbf{Role in the proof.}
Use in this chapter. It names a bad event or rejection condition whose probability must be zero or explicitly bounded.

\end{formaldefinition}

\begin{formaldefinition}[EasyCrypt line 296]
\noindent\textbf{EasyCrypt name.}
\begin{lstlisting}[language=EasyCrypt,basicstyle=\ttfamily\scriptsize]
signing_attempt_state_reject1_aborts
\end{lstlisting}
\noindent\textbf{Checked EasyCrypt statement.}
\begin{lstlisting}[language=EasyCrypt,basicstyle=\ttfamily\scriptsize]
op signing_attempt_state_reject1_aborts
   (md : mode) (st : signing_attempt_state) : bool =
  ! signing_attempt_state_reject1_accepts md st.
\end{lstlisting}
\noindent\textbf{Formal mathematical reading.}
Mathematical definition. This is a Boolean predicate.  The exact displayed EasyCrypt statement gives its arguments; mathematically it is an event, invariant, support condition, or well-formedness predicate over those arguments.

\noindent\textbf{Role in the proof.}
Use in this chapter. It is part of the exact formal vocabulary used by subsequent lemmas and games.

\end{formaldefinition}

\begin{formaldefinition}[EasyCrypt line 300]
\noindent\textbf{EasyCrypt name.}
\begin{lstlisting}[language=EasyCrypt,basicstyle=\ttfamily\scriptsize]
signing_attempt_state_pack_aborts
\end{lstlisting}
\noindent\textbf{Checked EasyCrypt statement.}
\begin{lstlisting}[language=EasyCrypt,basicstyle=\ttfamily\scriptsize]
op signing_attempt_state_pack_aborts
   (md : mode) (st : signing_attempt_state) : bool =
  ! signing_attempt_state_pack_accepts md st.
\end{lstlisting}
\noindent\textbf{Formal mathematical reading.}
Mathematical definition. This is a Boolean predicate.  The exact displayed EasyCrypt statement gives its arguments; mathematically it is an event, invariant, support condition, or well-formedness predicate over those arguments.

\noindent\textbf{Role in the proof.}
Use in this chapter. It is part of the exact formal vocabulary used by subsequent lemmas and games.

\end{formaldefinition}

\begin{formaldefinition}[EasyCrypt line 304]
\noindent\textbf{EasyCrypt name.}
\begin{lstlisting}[language=EasyCrypt,basicstyle=\ttfamily\scriptsize]
signing_attempt_state_verification_accepts
\end{lstlisting}
\noindent\textbf{Checked EasyCrypt statement.}
\begin{lstlisting}[language=EasyCrypt,basicstyle=\ttfamily\scriptsize]
op signing_attempt_state_verification_accepts
   (md : mode) (pk : pkey) (m : message) (ctx : context)
   (st : signing_attempt_state) : bool =
  signing_attempt_state_field_accepts md pk m ctx st /\
  signing_attempt_state_accepts md st.
\end{lstlisting}
\noindent\textbf{Formal mathematical reading.}
Mathematical definition. This is a Boolean predicate.  The exact displayed EasyCrypt statement gives its arguments; mathematically it is an event, invariant, support condition, or well-formedness predicate over those arguments.

\noindent\textbf{Role in the proof.}
Use in this chapter. It is part of the exact formal vocabulary used by subsequent lemmas and games.

\end{formaldefinition}

\begin{formaldefinition}[EasyCrypt line 310]
\noindent\textbf{EasyCrypt name.}
\begin{lstlisting}[language=EasyCrypt,basicstyle=\ttfamily\scriptsize]
signing_attempt_state_rejection_accepts
\end{lstlisting}
\noindent\textbf{Checked EasyCrypt statement.}
\begin{lstlisting}[language=EasyCrypt,basicstyle=\ttfamily\scriptsize]
op signing_attempt_state_rejection_accepts
   (md : mode) (pk : pkey) (m : message) (ctx : context)
   (st : signing_attempt_state) : bool =
  signing_attempt_state_verification_accepts md pk m ctx st /\
  signing_attempt_state_reference_rejection_accepts md st.
\end{lstlisting}
\noindent\textbf{Formal mathematical reading.}
Mathematical definition. This is a Boolean predicate.  The exact displayed EasyCrypt statement gives its arguments; mathematically it is an event, invariant, support condition, or well-formedness predicate over those arguments.

\noindent\textbf{Role in the proof.}
Use in this chapter. It names a bad event or rejection condition whose probability must be zero or explicitly bounded.

\end{formaldefinition}

\begin{formaldefinition}[EasyCrypt line 316]
\noindent\textbf{EasyCrypt name.}
\begin{lstlisting}[language=EasyCrypt,basicstyle=\ttfamily\scriptsize]
signing_attempt_state_rejection_aborts
\end{lstlisting}
\noindent\textbf{Checked EasyCrypt statement.}
\begin{lstlisting}[language=EasyCrypt,basicstyle=\ttfamily\scriptsize]
op signing_attempt_state_rejection_aborts
   (md : mode) (pk : pkey) (m : message) (ctx : context)
   (st : signing_attempt_state) : bool =
  ! signing_attempt_state_rejection_accepts md pk m ctx st.
\end{lstlisting}
\noindent\textbf{Formal mathematical reading.}
Mathematical definition. This is a Boolean predicate.  The exact displayed EasyCrypt statement gives its arguments; mathematically it is an event, invariant, support condition, or well-formedness predicate over those arguments.

\noindent\textbf{Role in the proof.}
Use in this chapter. It names a bad event or rejection condition whose probability must be zero or explicitly bounded.

\end{formaldefinition}

\begin{formaldefinition}[EasyCrypt line 321]
\noindent\textbf{EasyCrypt name.}
\begin{lstlisting}[language=EasyCrypt,basicstyle=\ttfamily\scriptsize]
signing_attempt_state_response_norm_aborts
\end{lstlisting}
\noindent\textbf{Checked EasyCrypt statement.}
\begin{lstlisting}[language=EasyCrypt,basicstyle=\ttfamily\scriptsize]
op signing_attempt_state_response_norm_aborts
   (md : mode) (st : signing_attempt_state) : bool =
  ! signing_attempt_state_response_norm_accepts md st.
\end{lstlisting}
\noindent\textbf{Formal mathematical reading.}
Mathematical definition. This is a Boolean predicate.  The exact displayed EasyCrypt statement gives its arguments; mathematically it is an event, invariant, support condition, or well-formedness predicate over those arguments.

\noindent\textbf{Role in the proof.}
Use in this chapter. It is part of the exact formal vocabulary used by subsequent lemmas and games.

\end{formaldefinition}

\begin{formaldefinition}[EasyCrypt line 325]
\noindent\textbf{EasyCrypt name.}
\begin{lstlisting}[language=EasyCrypt,basicstyle=\ttfamily\scriptsize]
signing_attempt_state_hint_norm_aborts
\end{lstlisting}
\noindent\textbf{Checked EasyCrypt statement.}
\begin{lstlisting}[language=EasyCrypt,basicstyle=\ttfamily\scriptsize]
op signing_attempt_state_hint_norm_aborts
   (md : mode) (st : signing_attempt_state) : bool =
  ! signing_attempt_state_hint_norm_accepts md st.
\end{lstlisting}
\noindent\textbf{Formal mathematical reading.}
Mathematical definition. This is a Boolean predicate.  The exact displayed EasyCrypt statement gives its arguments; mathematically it is an event, invariant, support condition, or well-formedness predicate over those arguments.

\noindent\textbf{Role in the proof.}
Use in this chapter. It is part of the exact formal vocabulary used by subsequent lemmas and games.

\end{formaldefinition}

\begin{formaldefinition}[EasyCrypt line 329]
\noindent\textbf{EasyCrypt name.}
\begin{lstlisting}[language=EasyCrypt,basicstyle=\ttfamily\scriptsize]
signing_attempt_state_rejection_sampling_aborts
\end{lstlisting}
\noindent\textbf{Checked EasyCrypt statement.}
\begin{lstlisting}[language=EasyCrypt,basicstyle=\ttfamily\scriptsize]
op signing_attempt_state_rejection_sampling_aborts
   (md : mode) (st : signing_attempt_state) : bool =
  ! signing_attempt_state_valid_signature_accepts md st.
\end{lstlisting}
\noindent\textbf{Formal mathematical reading.}
Mathematical definition. This is a Boolean predicate.  The exact displayed EasyCrypt statement gives its arguments; mathematically it is an event, invariant, support condition, or well-formedness predicate over those arguments.

\noindent\textbf{Role in the proof.}
Use in this chapter. It names a bad event or rejection condition whose probability must be zero or explicitly bounded.

\end{formaldefinition}

\begin{formaldefinition}[EasyCrypt line 333]
\noindent\textbf{EasyCrypt name.}
\begin{lstlisting}[language=EasyCrypt,basicstyle=\ttfamily\scriptsize]
signing_attempt_state_aborts
\end{lstlisting}
\noindent\textbf{Checked EasyCrypt statement.}
\begin{lstlisting}[language=EasyCrypt,basicstyle=\ttfamily\scriptsize]
op signing_attempt_state_aborts
   (md : mode) (st : signing_attempt_state) : bool =
  signing_attempt_state_rejection_sampling_aborts md st \/
  signing_attempt_state_response_norm_aborts md st \/
  signing_attempt_state_hint_norm_aborts md st.
\end{lstlisting}
\noindent\textbf{Formal mathematical reading.}
Mathematical definition. This is a Boolean predicate.  The exact displayed EasyCrypt statement gives its arguments; mathematically it is an event, invariant, support condition, or well-formedness predicate over those arguments.

\noindent\textbf{Role in the proof.}
Use in this chapter. It is part of the exact formal vocabulary used by subsequent lemmas and games.

\end{formaldefinition}

\begin{formaldefinition}[EasyCrypt line 339]
\noindent\textbf{EasyCrypt name.}
\begin{lstlisting}[language=EasyCrypt,basicstyle=\ttfamily\scriptsize]
signing_accepts
\end{lstlisting}
\noindent\textbf{Checked EasyCrypt statement.}
\begin{lstlisting}[language=EasyCrypt,basicstyle=\ttfamily\scriptsize]
op signing_accepts
   (md : mode) (sk : skey) (m : message) (ctx : context)
   (coins : random_coins) : bool =
  signing_attempt_state_accepts md
    (signing_attempt_state_of_coins md sk m ctx coins).
\end{lstlisting}
\noindent\textbf{Formal mathematical reading.}
Mathematical definition. This is a Boolean predicate.  The exact displayed EasyCrypt statement gives its arguments; mathematically it is an event, invariant, support condition, or well-formedness predicate over those arguments.

\noindent\textbf{Role in the proof.}
Use in this chapter. It is part of the exact formal vocabulary used by subsequent lemmas and games.

\end{formaldefinition}

\begin{formaldefinition}[EasyCrypt line 345]
\noindent\textbf{EasyCrypt name.}
\begin{lstlisting}[language=EasyCrypt,basicstyle=\ttfamily\scriptsize]
dsigning_attempt_state
\end{lstlisting}
\noindent\textbf{Checked EasyCrypt statement.}
\begin{lstlisting}[language=EasyCrypt,basicstyle=\ttfamily\scriptsize]
op dsigning_attempt_state
   (md : mode) (sk : skey) (m : message) (ctx : context) :
   signing_attempt_state distr =
  dmap drandom_coins (signing_attempt_state_of_coins md sk m ctx).
\end{lstlisting}
\noindent\textbf{Formal mathematical reading.}
Mathematical definition. This is a total EasyCrypt operation.  The exact displayed statement gives the domain, codomain, and defining equation when present.

\noindent\textbf{Role in the proof.}
Use in this chapter. It contributes a sampler or sampler-derived object to the distributional model.

\end{formaldefinition}

\begin{formaldefinition}[EasyCrypt line 350]
\noindent\textbf{EasyCrypt name.}
\begin{lstlisting}[language=EasyCrypt,basicstyle=\ttfamily\scriptsize]
signing_sample_pair_sampler_ok
\end{lstlisting}
\noindent\textbf{Checked EasyCrypt statement.}
\begin{lstlisting}[language=EasyCrypt,basicstyle=\ttfamily\scriptsize]
op signing_sample_pair_sampler_ok
   (md : mode) (dsample : random_coins -> signing_sample_pair distr) :
   bool =
  forall coins,
    signing_randomness_domain coins =>
    signing_sample_pair_distribution_ok md (dsample coins).
\end{lstlisting}
\noindent\textbf{Formal mathematical reading.}
Mathematical definition. This is a Boolean predicate.  The exact displayed EasyCrypt statement gives its arguments; mathematically it is an event, invariant, support condition, or well-formedness predicate over those arguments.

\noindent\textbf{Role in the proof.}
Use in this chapter. It names a well-formedness, support, or validity predicate needed by later preservation lemmas.

\end{formaldefinition}

\begin{formaldefinition}[EasyCrypt line 357]
\noindent\textbf{EasyCrypt name.}
\begin{lstlisting}[language=EasyCrypt,basicstyle=\ttfamily\scriptsize]
dsigning_attempt_state_from_sample_pair_sampler
\end{lstlisting}
\noindent\textbf{Checked EasyCrypt statement.}
\begin{lstlisting}[language=EasyCrypt,basicstyle=\ttfamily\scriptsize]
op dsigning_attempt_state_from_sample_pair_sampler
   (md : mode) (sk : skey) (m : message) (ctx : context)
   (dsample : random_coins -> signing_sample_pair distr) :
   signing_attempt_state distr =
  dlet drandom_coins
    (fun coins =>
      dmap (dsample coins)
        (signing_attempt_state_of_sample_pair md sk m ctx coins)).
\end{lstlisting}
\noindent\textbf{Formal mathematical reading.}
Mathematical definition. This is a total EasyCrypt operation.  The exact displayed statement gives the domain, codomain, and defining equation when present.

\noindent\textbf{Role in the proof.}
Use in this chapter. It contributes a sampler or sampler-derived object to the distributional model.

\end{formaldefinition}

\begin{formaldefinition}[EasyCrypt line 366]
\noindent\textbf{EasyCrypt name.}
\begin{lstlisting}[language=EasyCrypt,basicstyle=\ttfamily\scriptsize]
signing_attempt_coin_sample_pair
\end{lstlisting}
\noindent\textbf{Checked EasyCrypt statement.}
\begin{lstlisting}[language=EasyCrypt,basicstyle=\ttfamily\scriptsize]
type signing_attempt_coin_sample_pair = random_coins * signing_sample_pair.
\end{lstlisting}
\noindent\textbf{Formal mathematical reading.}
Mathematical definition. This is a definitional type abbreviation.  The exact displayed EasyCrypt statement gives the right-hand side; later proof steps may unfold it without adding an assumption.

\noindent\textbf{Role in the proof.}
Use in this chapter. It contributes a sampler or sampler-derived object to the distributional model.

\end{formaldefinition}

\begin{formaldefinition}[EasyCrypt line 368]
\noindent\textbf{EasyCrypt name.}
\begin{lstlisting}[language=EasyCrypt,basicstyle=\ttfamily\scriptsize]
signing_attempt_state_of_coin_sample_pair
\end{lstlisting}
\noindent\textbf{Checked EasyCrypt statement.}
\begin{lstlisting}[language=EasyCrypt,basicstyle=\ttfamily\scriptsize]
op signing_attempt_state_of_coin_sample_pair
   (md : mode) (sk : skey) (m : message) (ctx : context)
   (cs : signing_attempt_coin_sample_pair) : signing_attempt_state =
  signing_attempt_state_of_sample_pair md sk m ctx cs.`1 cs.`2.
\end{lstlisting}
\noindent\textbf{Formal mathematical reading.}
Mathematical definition. This is a total EasyCrypt operation.  The exact displayed statement gives the domain, codomain, and defining equation when present.

\noindent\textbf{Role in the proof.}
Use in this chapter. It contributes a sampler or sampler-derived object to the distributional model.

\end{formaldefinition}

\begin{formaldefinition}[EasyCrypt line 373]
\noindent\textbf{EasyCrypt name.}
\begin{lstlisting}[language=EasyCrypt,basicstyle=\ttfamily\scriptsize]
dstructural_signing_attempt_coin_sample_pair
\end{lstlisting}
\noindent\textbf{Checked EasyCrypt statement.}
\begin{lstlisting}[language=EasyCrypt,basicstyle=\ttfamily\scriptsize]
op dstructural_signing_attempt_coin_sample_pair
   (md : mode) (sk : skey) (m : message) (ctx : context) :
   signing_attempt_coin_sample_pair distr =
  dlet drandom_coins
    (fun coins =>
      dmap (dstructural_signing_sample_pair md sk m ctx coins)
        (fun sample => (coins, sample))).
\end{lstlisting}
\noindent\textbf{Formal mathematical reading.}
Mathematical definition. This is a total EasyCrypt operation.  The exact displayed statement gives the domain, codomain, and defining equation when present.

\noindent\textbf{Role in the proof.}
Use in this chapter. It contributes a sampler or sampler-derived object to the distributional model.

\end{formaldefinition}

\begin{formaldefinition}[EasyCrypt line 381]
\noindent\textbf{EasyCrypt name.}
\begin{lstlisting}[language=EasyCrypt,basicstyle=\ttfamily\scriptsize]
dexact_hyperball_signing_attempt_coin_sample_pair
\end{lstlisting}
\noindent\textbf{Checked EasyCrypt statement.}
\begin{lstlisting}[language=EasyCrypt,basicstyle=\ttfamily\scriptsize]
op dexact_hyperball_signing_attempt_coin_sample_pair
   (md : mode) (sk : skey) (m : message) (ctx : context) :
   signing_attempt_coin_sample_pair distr =
  dlet drandom_coins
    (fun coins =>
      dmap (dexact_hyperball_signing_sample_pair md)
        (fun sample => (coins, sample))).
\end{lstlisting}
\noindent\textbf{Formal mathematical reading.}
Mathematical definition. This is a total EasyCrypt operation.  The exact displayed statement gives the domain, codomain, and defining equation when present.

\noindent\textbf{Role in the proof.}
Use in this chapter. It contributes a sampler or sampler-derived object to the distributional model.

\end{formaldefinition}

\begin{formaldefinition}[EasyCrypt line 389]
\noindent\textbf{EasyCrypt name.}
\begin{lstlisting}[language=EasyCrypt,basicstyle=\ttfamily\scriptsize]
dstructural_signing_attempt_state
\end{lstlisting}
\noindent\textbf{Checked EasyCrypt statement.}
\begin{lstlisting}[language=EasyCrypt,basicstyle=\ttfamily\scriptsize]
op dstructural_signing_attempt_state
   (md : mode) (sk : skey) (m : message) (ctx : context) :
   signing_attempt_state distr =
  dmap drandom_coins
    (fun coins =>
      signing_attempt_state_of_sample_pair md sk m ctx coins
        (signing_sample_pair_of_coins md sk m ctx coins)).
\end{lstlisting}
\noindent\textbf{Formal mathematical reading.}
Mathematical definition. This is a total EasyCrypt operation.  The exact displayed statement gives the domain, codomain, and defining equation when present.

\noindent\textbf{Role in the proof.}
Use in this chapter. It contributes a sampler or sampler-derived object to the distributional model.

\end{formaldefinition}

\begin{formaldefinition}[EasyCrypt line 397]
\noindent\textbf{EasyCrypt name.}
\begin{lstlisting}[language=EasyCrypt,basicstyle=\ttfamily\scriptsize]
dstructural_signing_attempt_state_from_sampler
\end{lstlisting}
\noindent\textbf{Checked EasyCrypt statement.}
\begin{lstlisting}[language=EasyCrypt,basicstyle=\ttfamily\scriptsize]
op dstructural_signing_attempt_state_from_sampler
   (md : mode) (sk : skey) (m : message) (ctx : context) :
   signing_attempt_state distr =
  dsigning_attempt_state_from_sample_pair_sampler md sk m ctx
    (dstructural_signing_sample_pair md sk m ctx).
\end{lstlisting}
\noindent\textbf{Formal mathematical reading.}
Mathematical definition. This is a total EasyCrypt operation.  The exact displayed statement gives the domain, codomain, and defining equation when present.

\noindent\textbf{Role in the proof.}
Use in this chapter. It contributes a sampler or sampler-derived object to the distributional model.

\end{formaldefinition}

\begin{formaldefinition}[EasyCrypt line 403]
\noindent\textbf{EasyCrypt name.}
\begin{lstlisting}[language=EasyCrypt,basicstyle=\ttfamily\scriptsize]
dbounded_unit_signing_attempt_state
\end{lstlisting}
\noindent\textbf{Checked EasyCrypt statement.}
\begin{lstlisting}[language=EasyCrypt,basicstyle=\ttfamily\scriptsize]
op dbounded_unit_signing_attempt_state
   (md : mode) (sk : skey) (m : message) (ctx : context) :
   signing_attempt_state distr =
  dsigning_attempt_state_from_sample_pair_sampler md sk m ctx
    (fun _ => dbounded_unit_signing_sample_pair md).
\end{lstlisting}
\noindent\textbf{Formal mathematical reading.}
Mathematical definition. This is a total EasyCrypt operation.  The exact displayed statement gives the domain, codomain, and defining equation when present.

\noindent\textbf{Role in the proof.}
Use in this chapter. It names a numerical term that appears in a game-hop or final security inequality.

\end{formaldefinition}

\begin{formaldefinition}[EasyCrypt line 409]
\noindent\textbf{EasyCrypt name.}
\begin{lstlisting}[language=EasyCrypt,basicstyle=\ttfamily\scriptsize]
dchecked_hyperball_signing_attempt_state
\end{lstlisting}
\noindent\textbf{Checked EasyCrypt statement.}
\begin{lstlisting}[language=EasyCrypt,basicstyle=\ttfamily\scriptsize]
op dchecked_hyperball_signing_attempt_state
   (md : mode) (sk : skey) (m : message) (ctx : context) :
   signing_attempt_state distr =
  dsigning_attempt_state_from_sample_pair_sampler md sk m ctx
    (fun _ => dchecked_hyperball_signing_sample_pair md).
\end{lstlisting}
\noindent\textbf{Formal mathematical reading.}
Mathematical definition. This is a total EasyCrypt operation.  The exact displayed statement gives the domain, codomain, and defining equation when present.

\noindent\textbf{Role in the proof.}
Use in this chapter. It contributes a sampler or sampler-derived object to the distributional model.

\end{formaldefinition}

\begin{formaldefinition}[EasyCrypt line 415]
\noindent\textbf{EasyCrypt name.}
\begin{lstlisting}[language=EasyCrypt,basicstyle=\ttfamily\scriptsize]
dexact_hyperball_signing_attempt_state
\end{lstlisting}
\noindent\textbf{Checked EasyCrypt statement.}
\begin{lstlisting}[language=EasyCrypt,basicstyle=\ttfamily\scriptsize]
op dexact_hyperball_signing_attempt_state
   (md : mode) (sk : skey) (m : message) (ctx : context) :
   signing_attempt_state distr =
  dsigning_attempt_state_from_sample_pair_sampler md sk m ctx
    (fun _ => dexact_hyperball_signing_sample_pair md).
\end{lstlisting}
\noindent\textbf{Formal mathematical reading.}
Mathematical definition. This is a total EasyCrypt operation.  The exact displayed statement gives the domain, codomain, and defining equation when present.

\noindent\textbf{Role in the proof.}
Use in this chapter. It contributes a sampler or sampler-derived object to the distributional model.

\end{formaldefinition}

\begin{formaldefinition}[EasyCrypt line 421]
\noindent\textbf{EasyCrypt name.}
\begin{lstlisting}[language=EasyCrypt,basicstyle=\ttfamily\scriptsize]
signing_attempt_state_sample_pair_wf
\end{lstlisting}
\noindent\textbf{Checked EasyCrypt statement.}
\begin{lstlisting}[language=EasyCrypt,basicstyle=\ttfamily\scriptsize]
op signing_attempt_state_sample_pair_wf
   (md : mode) (st : signing_attempt_state) : bool =
  polyvecl_wf md (signing_attempt_state_sampled_y1 st) /\
  polyveck_wf md (signing_attempt_state_sampled_y2 st).
\end{lstlisting}
\noindent\textbf{Formal mathematical reading.}
Mathematical definition. This is a Boolean predicate.  The exact displayed EasyCrypt statement gives its arguments; mathematically it is an event, invariant, support condition, or well-formedness predicate over those arguments.

\noindent\textbf{Role in the proof.}
Use in this chapter. It names a well-formedness, support, or validity predicate needed by later preservation lemmas.

\end{formaldefinition}

\begin{formaldefinition}[EasyCrypt line 426]
\noindent\textbf{EasyCrypt name.}
\begin{lstlisting}[language=EasyCrypt,basicstyle=\ttfamily\scriptsize]
signing_attempt_state_sample_pair_unit
\end{lstlisting}
\noindent\textbf{Checked EasyCrypt statement.}
\begin{lstlisting}[language=EasyCrypt,basicstyle=\ttfamily\scriptsize]
op signing_attempt_state_sample_pair_unit
   (md : mode) (st : signing_attempt_state) : bool =
  polyvecl_unit md (signing_attempt_state_sampled_y1 st) /\
  polyveck_unit md (signing_attempt_state_sampled_y2 st).
\end{lstlisting}
\noindent\textbf{Formal mathematical reading.}
Mathematical definition. This is a Boolean predicate.  The exact displayed EasyCrypt statement gives its arguments; mathematically it is an event, invariant, support condition, or well-formedness predicate over those arguments.

\noindent\textbf{Role in the proof.}
Use in this chapter. It contributes a sampler or sampler-derived object to the distributional model.

\end{formaldefinition}

\begin{formaldefinition}[EasyCrypt line 431]
\noindent\textbf{EasyCrypt name.}
\begin{lstlisting}[language=EasyCrypt,basicstyle=\ttfamily\scriptsize]
signing_attempt_state_sample_pair_bounded
\end{lstlisting}
\noindent\textbf{Checked EasyCrypt statement.}
\begin{lstlisting}[language=EasyCrypt,basicstyle=\ttfamily\scriptsize]
op signing_attempt_state_sample_pair_bounded
   (md : mode) (st : signing_attempt_state) : bool =
  polyvecl_bounded_by md (signing_sample_bound md)
    (signing_attempt_state_sampled_y1 st) /\
  polyveck_bounded_by md (signing_sample_bound md)
    (signing_attempt_state_sampled_y2 st).
\end{lstlisting}
\noindent\textbf{Formal mathematical reading.}
Mathematical definition. This is a Boolean predicate.  The exact displayed EasyCrypt statement gives its arguments; mathematically it is an event, invariant, support condition, or well-formedness predicate over those arguments.

\noindent\textbf{Role in the proof.}
Use in this chapter. It names a numerical term that appears in a game-hop or final security inequality.

\end{formaldefinition}

\begin{formaldefinition}[EasyCrypt line 438]
\noindent\textbf{EasyCrypt name.}
\begin{lstlisting}[language=EasyCrypt,basicstyle=\ttfamily\scriptsize]
signing_attempt_relation
\end{lstlisting}
\noindent\textbf{Checked EasyCrypt statement.}
\begin{lstlisting}[language=EasyCrypt,basicstyle=\ttfamily\scriptsize]
op signing_attempt_relation
   (md : mode) (pk : pkey) (sk : skey) (m : message)
   (ctx : context) (coins : random_coins) (sig : signature) : bool =
  valid_keypair md pk sk /\
  signing_accepts md sk m ctx coins /\
  sig = sign_internal md sk m ctx coins.
\end{lstlisting}
\noindent\textbf{Formal mathematical reading.}
Mathematical definition. This is a Boolean predicate.  The exact displayed EasyCrypt statement gives its arguments; mathematically it is an event, invariant, support condition, or well-formedness predicate over those arguments.

\noindent\textbf{Role in the proof.}
Use in this chapter. It is part of the exact formal vocabulary used by subsequent lemmas and games.

\end{formaldefinition}

\begin{formaldefinition}[EasyCrypt line 445]
\noindent\textbf{EasyCrypt name.}
\begin{lstlisting}[language=EasyCrypt,basicstyle=\ttfamily\scriptsize]
signing_simulation_relation
\end{lstlisting}
\noindent\textbf{Checked EasyCrypt statement.}
\begin{lstlisting}[language=EasyCrypt,basicstyle=\ttfamily\scriptsize]
op signing_simulation_relation
   (md : mode) (pk : pkey) (m : message) (ctx : context)
   (sig : signature) : bool =
  verify_internal md pk m ctx sig /\
  exists (sk : skey) (coins : random_coins),
    signing_attempt_relation md pk sk m ctx coins sig.
\end{lstlisting}
\noindent\textbf{Formal mathematical reading.}
Mathematical definition. This is a Boolean predicate.  The exact displayed EasyCrypt statement gives its arguments; mathematically it is an event, invariant, support condition, or well-formedness predicate over those arguments.

\noindent\textbf{Role in the proof.}
Use in this chapter. It is part of the exact formal vocabulary used by subsequent lemmas and games.

\end{formaldefinition}

\begin{formaldefinition}[EasyCrypt line 452]
\noindent\textbf{EasyCrypt name.}
\begin{lstlisting}[language=EasyCrypt,basicstyle=\ttfamily\scriptsize]
signing_transcript_relation
\end{lstlisting}
\noindent\textbf{Checked EasyCrypt statement.}
\begin{lstlisting}[language=EasyCrypt,basicstyle=\ttfamily\scriptsize]
op signing_transcript_relation
   (md : mode) (pk : pkey) (m : message) (ctx : context)
   (sig : signature) (tr : transcript) : bool =
  signing_simulation_relation md pk m ctx sig /\
  tr = transcript_of_signature md pk m ctx sig /\
  transcript_valid md tr.
\end{lstlisting}
\noindent\textbf{Formal mathematical reading.}
Mathematical definition. This is a Boolean predicate.  The exact displayed EasyCrypt statement gives its arguments; mathematically it is an event, invariant, support condition, or well-formedness predicate over those arguments.

\noindent\textbf{Role in the proof.}
Use in this chapter. It defines transcript/extraction data for the Fiat-Shamir forking and special-soundness argument.

\end{formaldefinition}

\begin{formaldefinition}[EasyCrypt line 459]
\noindent\textbf{EasyCrypt name.}
\begin{lstlisting}[language=EasyCrypt,basicstyle=\ttfamily\scriptsize]
signing_record_message
\end{lstlisting}
\noindent\textbf{Checked EasyCrypt statement.}
\begin{lstlisting}[language=EasyCrypt,basicstyle=\ttfamily\scriptsize]
op signing_record_message (r : signing_transcript_record) : message = r.`1.
\end{lstlisting}
\noindent\textbf{Formal mathematical reading.}
Mathematical definition. This is a total EasyCrypt operation.  The exact displayed statement gives the domain, codomain, and defining equation when present.

\noindent\textbf{Role in the proof.}
Use in this chapter. It is part of the exact formal vocabulary used by subsequent lemmas and games.

\end{formaldefinition}

\begin{formaldefinition}[EasyCrypt line 460]
\noindent\textbf{EasyCrypt name.}
\begin{lstlisting}[language=EasyCrypt,basicstyle=\ttfamily\scriptsize]
signing_record_context
\end{lstlisting}
\noindent\textbf{Checked EasyCrypt statement.}
\begin{lstlisting}[language=EasyCrypt,basicstyle=\ttfamily\scriptsize]
op signing_record_context (r : signing_transcript_record) : context = r.`2.
\end{lstlisting}
\noindent\textbf{Formal mathematical reading.}
Mathematical definition. This is a total EasyCrypt operation.  The exact displayed statement gives the domain, codomain, and defining equation when present.

\noindent\textbf{Role in the proof.}
Use in this chapter. It is part of the exact formal vocabulary used by subsequent lemmas and games.

\end{formaldefinition}

\begin{formaldefinition}[EasyCrypt line 461]
\noindent\textbf{EasyCrypt name.}
\begin{lstlisting}[language=EasyCrypt,basicstyle=\ttfamily\scriptsize]
signing_record_signature
\end{lstlisting}
\noindent\textbf{Checked EasyCrypt statement.}
\begin{lstlisting}[language=EasyCrypt,basicstyle=\ttfamily\scriptsize]
op signing_record_signature (r : signing_transcript_record) : signature = r.`3.
\end{lstlisting}
\noindent\textbf{Formal mathematical reading.}
Mathematical definition. This is a total EasyCrypt operation.  The exact displayed statement gives the domain, codomain, and defining equation when present.

\noindent\textbf{Role in the proof.}
Use in this chapter. It is part of the exact formal vocabulary used by subsequent lemmas and games.

\end{formaldefinition}

\begin{formaldefinition}[EasyCrypt line 462]
\noindent\textbf{EasyCrypt name.}
\begin{lstlisting}[language=EasyCrypt,basicstyle=\ttfamily\scriptsize]
signing_record_transcript
\end{lstlisting}
\noindent\textbf{Checked EasyCrypt statement.}
\begin{lstlisting}[language=EasyCrypt,basicstyle=\ttfamily\scriptsize]
op signing_record_transcript (r : signing_transcript_record) : transcript =
  r.`4.
\end{lstlisting}
\noindent\textbf{Formal mathematical reading.}
Mathematical definition. This is a total EasyCrypt operation.  The exact displayed statement gives the domain, codomain, and defining equation when present.

\noindent\textbf{Role in the proof.}
Use in this chapter. It defines transcript/extraction data for the Fiat-Shamir forking and special-soundness argument.

\end{formaldefinition}

\begin{formaldefinition}[EasyCrypt line 465]
\noindent\textbf{EasyCrypt name.}
\begin{lstlisting}[language=EasyCrypt,basicstyle=\ttfamily\scriptsize]
signing_record_query
\end{lstlisting}
\noindent\textbf{Checked EasyCrypt statement.}
\begin{lstlisting}[language=EasyCrypt,basicstyle=\ttfamily\scriptsize]
op signing_record_query (r : signing_transcript_record) : message * context =
  (signing_record_message r, signing_record_context r).
\end{lstlisting}
\noindent\textbf{Formal mathematical reading.}
Mathematical definition. This is a total EasyCrypt operation.  The exact displayed statement gives the domain, codomain, and defining equation when present.

\noindent\textbf{Role in the proof.}
Use in this chapter. It is part of the exact formal vocabulary used by subsequent lemmas and games.

\end{formaldefinition}

\begin{formaldefinition}[EasyCrypt line 468]
\noindent\textbf{EasyCrypt name.}
\begin{lstlisting}[language=EasyCrypt,basicstyle=\ttfamily\scriptsize]
signing_record_relation
\end{lstlisting}
\noindent\textbf{Checked EasyCrypt statement.}
\begin{lstlisting}[language=EasyCrypt,basicstyle=\ttfamily\scriptsize]
op signing_record_relation
   (md : mode) (pk : pkey) (r : signing_transcript_record) : bool =
  signing_transcript_relation md pk
    (signing_record_message r)
    (signing_record_context r)
    (signing_record_signature r)
    (signing_record_transcript r).
\end{lstlisting}
\noindent\textbf{Formal mathematical reading.}
Mathematical definition. This is a Boolean predicate.  The exact displayed EasyCrypt statement gives its arguments; mathematically it is an event, invariant, support condition, or well-formedness predicate over those arguments.

\noindent\textbf{Role in the proof.}
Use in this chapter. It is part of the exact formal vocabulary used by subsequent lemmas and games.

\end{formaldefinition}

\begin{formaldefinition}[EasyCrypt line 476]
\noindent\textbf{EasyCrypt name.}
\begin{lstlisting}[language=EasyCrypt,basicstyle=\ttfamily\scriptsize]
signing_record_queries
\end{lstlisting}
\noindent\textbf{Checked EasyCrypt statement.}
\begin{lstlisting}[language=EasyCrypt,basicstyle=\ttfamily\scriptsize]
op signing_record_queries (rs : signing_transcript_record list) :
  (message * context) list =
  map signing_record_query rs.
\end{lstlisting}
\noindent\textbf{Formal mathematical reading.}
Mathematical definition. This is a total EasyCrypt operation.  The exact displayed statement gives the domain, codomain, and defining equation when present.

\noindent\textbf{Role in the proof.}
Use in this chapter. It is part of the exact formal vocabulary used by subsequent lemmas and games.

\end{formaldefinition}

\begin{formaldefinition}[EasyCrypt line 480]
\noindent\textbf{EasyCrypt name.}
\begin{lstlisting}[language=EasyCrypt,basicstyle=\ttfamily\scriptsize]
signing_record_transcripts
\end{lstlisting}
\noindent\textbf{Checked EasyCrypt statement.}
\begin{lstlisting}[language=EasyCrypt,basicstyle=\ttfamily\scriptsize]
op signing_record_transcripts (rs : signing_transcript_record list) :
  transcript list =
  map signing_record_transcript rs.
\end{lstlisting}
\noindent\textbf{Formal mathematical reading.}
Mathematical definition. This is a total EasyCrypt operation.  The exact displayed statement gives the domain, codomain, and defining equation when present.

\noindent\textbf{Role in the proof.}
Use in this chapter. It defines transcript/extraction data for the Fiat-Shamir forking and special-soundness argument.

\end{formaldefinition}

\begin{formaldefinition}[EasyCrypt line 484]
\noindent\textbf{EasyCrypt name.}
\begin{lstlisting}[language=EasyCrypt,basicstyle=\ttfamily\scriptsize]
signing_record_log_sound
\end{lstlisting}
\noindent\textbf{Checked EasyCrypt statement.}
\begin{lstlisting}[language=EasyCrypt,basicstyle=\ttfamily\scriptsize]
op signing_record_log_sound
   (md : mode) (pk : pkey) (rs : signing_transcript_record list) : bool =
  all (signing_record_relation md pk) rs.
\end{lstlisting}
\noindent\textbf{Formal mathematical reading.}
Mathematical definition. This is a Boolean predicate.  The exact displayed EasyCrypt statement gives its arguments; mathematically it is an event, invariant, support condition, or well-formedness predicate over those arguments.

\noindent\textbf{Role in the proof.}
Use in this chapter. It is part of the exact formal vocabulary used by subsequent lemmas and games.

\end{formaldefinition}

\begin{formaldefinition}[EasyCrypt line 488]
\noindent\textbf{EasyCrypt name.}
\begin{lstlisting}[language=EasyCrypt,basicstyle=\ttfamily\scriptsize]
transcript_log_valid
\end{lstlisting}
\noindent\textbf{Checked EasyCrypt statement.}
\begin{lstlisting}[language=EasyCrypt,basicstyle=\ttfamily\scriptsize]
op transcript_log_valid (md : mode) (trs : transcript list) : bool =
  all (transcript_valid md) trs.
\end{lstlisting}
\noindent\textbf{Formal mathematical reading.}
Mathematical definition. This is a Boolean predicate.  The exact displayed EasyCrypt statement gives its arguments; mathematically it is an event, invariant, support condition, or well-formedness predicate over those arguments.

\noindent\textbf{Role in the proof.}
Use in this chapter. It names a well-formedness, support, or validity predicate needed by later preservation lemmas.

\end{formaldefinition}

\begin{formaldefinition}[EasyCrypt line 491]
\noindent\textbf{EasyCrypt name.}
\begin{lstlisting}[language=EasyCrypt,basicstyle=\ttfamily\scriptsize]
structural_to_exact_hyperball_sample_pair_point_loss_obligation
\end{lstlisting}
\noindent\textbf{Checked EasyCrypt statement.}
\begin{lstlisting}[language=EasyCrypt,basicstyle=\ttfamily\scriptsize]
op structural_to_exact_hyperball_sample_pair_point_loss_obligation : bool =
  forall (md : mode) (sk : skey) (m : message) (ctx : context)
         (coins : random_coins),
    1%r <=
    mu (dexact_hyperball_signing_sample_pair md)
      (pred1 (signing_sample_pair_of_coins md sk m ctx coins)) +
      rejection_sampling_loss_term.
\end{lstlisting}
\noindent\textbf{Formal mathematical reading.}
Mathematical definition. This is a Boolean predicate.  The exact displayed EasyCrypt statement gives its arguments; mathematically it is an event, invariant, support condition, or well-formedness predicate over those arguments.

\noindent\textbf{Role in the proof.}
Use in this chapter. It names a numerical term that appears in a game-hop or final security inequality.

\end{formaldefinition}

\begin{formaldefinition}[EasyCrypt line 499]
\noindent\textbf{EasyCrypt name.}
\begin{lstlisting}[language=EasyCrypt,basicstyle=\ttfamily\scriptsize]
structural_to_exact_hyperball_sample_pair_loss_obligation
\end{lstlisting}
\noindent\textbf{Checked EasyCrypt statement.}
\begin{lstlisting}[language=EasyCrypt,basicstyle=\ttfamily\scriptsize]
op structural_to_exact_hyperball_sample_pair_loss_obligation : bool =
  forall (md : mode) (sk : skey) (m : message) (ctx : context)
         (coins : random_coins) (p : signing_sample_pair -> bool),
    mu (dstructural_signing_sample_pair md sk m ctx coins) p <=
    mu (dexact_hyperball_signing_sample_pair md) p +
      rejection_sampling_loss_term.
\end{lstlisting}
\noindent\textbf{Formal mathematical reading.}
Mathematical definition. This is a Boolean predicate.  The exact displayed EasyCrypt statement gives its arguments; mathematically it is an event, invariant, support condition, or well-formedness predicate over those arguments.

\noindent\textbf{Role in the proof.}
Use in this chapter. It names a numerical term that appears in a game-hop or final security inequality.

\end{formaldefinition}

\begin{formaldefinition}[EasyCrypt line 506]
\noindent\textbf{EasyCrypt name.}
\begin{lstlisting}[language=EasyCrypt,basicstyle=\ttfamily\scriptsize]
structural_to_exact_hyperball_coin_sample_pair_loss_obligation
\end{lstlisting}
\noindent\textbf{Checked EasyCrypt statement.}
\begin{lstlisting}[language=EasyCrypt,basicstyle=\ttfamily\scriptsize]
op structural_to_exact_hyperball_coin_sample_pair_loss_obligation : bool =
  forall (md : mode) (sk : skey) (m : message) (ctx : context)
         (p : signing_attempt_coin_sample_pair -> bool),
    mu (dstructural_signing_attempt_coin_sample_pair md sk m ctx) p <=
    mu (dexact_hyperball_signing_attempt_coin_sample_pair md sk m ctx) p +
      rejection_sampling_loss_term.
\end{lstlisting}
\noindent\textbf{Formal mathematical reading.}
Mathematical definition. This is a Boolean predicate.  The exact displayed EasyCrypt statement gives its arguments; mathematically it is an event, invariant, support condition, or well-formedness predicate over those arguments.

\noindent\textbf{Role in the proof.}
Use in this chapter. It names a numerical term that appears in a game-hop or final security inequality.

\end{formaldefinition}

\begin{formaldefinition}[EasyCrypt line 513]
\noindent\textbf{EasyCrypt name.}
\begin{lstlisting}[language=EasyCrypt,basicstyle=\ttfamily\scriptsize]
structural_to_exact_hyperball_attempt_loss_obligation
\end{lstlisting}
\noindent\textbf{Checked EasyCrypt statement.}
\begin{lstlisting}[language=EasyCrypt,basicstyle=\ttfamily\scriptsize]
op structural_to_exact_hyperball_attempt_loss_obligation : bool =
  forall (md : mode) (sk : skey) (m : message) (ctx : context)
         (p : signing_attempt_state -> bool),
    mu (dsigning_attempt_state md sk m ctx) p <=
    mu (dexact_hyperball_signing_attempt_state md sk m ctx) p +
      rejection_sampling_loss_term.
\end{lstlisting}
\noindent\textbf{Formal mathematical reading.}
Mathematical definition. This is a Boolean predicate.  The exact displayed EasyCrypt statement gives its arguments; mathematically it is an event, invariant, support condition, or well-formedness predicate over those arguments.

\noindent\textbf{Role in the proof.}
Use in this chapter. It names a numerical term that appears in a game-hop or final security inequality.

\end{formaldefinition}

\begin{formaldefinition}[EasyCrypt line 520]
\noindent\textbf{EasyCrypt name.}
\begin{lstlisting}[language=EasyCrypt,basicstyle=\ttfamily\scriptsize]
rejection_sampling_bound_obligation
\end{lstlisting}
\noindent\textbf{Checked EasyCrypt statement.}
\begin{lstlisting}[language=EasyCrypt,basicstyle=\ttfamily\scriptsize]
op rejection_sampling_bound_obligation : bool =
  0%r <= rejection_sampling_loss_term /\
  0%r <= fs_with_aborts_reprogramming_term /\
  0%r <= fs_with_aborts_min_entropy_term.
\end{lstlisting}
\noindent\textbf{Formal mathematical reading.}
Mathematical definition. This is a Boolean predicate.  The exact displayed EasyCrypt statement gives its arguments; mathematically it is an event, invariant, support condition, or well-formedness predicate over those arguments.

\noindent\textbf{Role in the proof.}
Use in this chapter. It names a numerical term that appears in a game-hop or final security inequality.

\end{formaldefinition}

\subsubsection*{Lemmas, Theorems, Relational Equivalences, and Their Proof Dependencies}
\begin{formaltheorem}[EasyCrypt line 525]
\noindent\textbf{EasyCrypt name.}
\begin{lstlisting}[language=EasyCrypt,basicstyle=\ttfamily\scriptsize]
signing_attempt_verifies
\end{lstlisting}
\noindent\textbf{Checked EasyCrypt statement.}
\begin{lstlisting}[language=EasyCrypt,basicstyle=\ttfamily\scriptsize]
lemma signing_attempt_verifies md pk sk m ctx coins sig :
  signing_attempt_relation md pk sk m ctx coins sig =>
  verify_internal md pk m ctx sig.
\end{lstlisting}
\noindent\textbf{Formal mathematical reading.}
Mathematical theorem. This is an implication, normally turning one invariant, validity predicate, or proof obligation into another.

\noindent\textbf{Role in the proof.}
Use in this chapter. It is a reusable checked theorem cited by later proof scripts.

\noindent\textbf{Proof-script interpretation.}
Proof-script reading. The machine proof decomposes the theorem into rewrites, program-logic obligations, relational equivalences, and arithmetic side conditions.
\emph{Main tactics visible in the proof script:} \ec{rewrite}, \ec{apply}.
\emph{Named identifiers syntactically cited in the proof block:}
\begin{lstlisting}[language=EasyCrypt,basicstyle=\ttfamily\scriptsize]
pkE
signing_attempt_relation
valid_keypair
verify_internal_sign_internal
\end{lstlisting}

\end{formaltheorem}

\begin{formaltheorem}[EasyCrypt line 535]
\noindent\textbf{EasyCrypt name.}
\begin{lstlisting}[language=EasyCrypt,basicstyle=\ttfamily\scriptsize]
signing_simulation_verifies
\end{lstlisting}
\noindent\textbf{Checked EasyCrypt statement.}
\begin{lstlisting}[language=EasyCrypt,basicstyle=\ttfamily\scriptsize]
lemma signing_simulation_verifies md pk m ctx sig :
  signing_simulation_relation md pk m ctx sig =>
  verify_internal md pk m ctx sig.
\end{lstlisting}
\noindent\textbf{Formal mathematical reading.}
Mathematical theorem. This is an implication, normally turning one invariant, validity predicate, or proof obligation into another.

\noindent\textbf{Role in the proof.}
Use in this chapter. It proves an exact game replacement or simulator equivalence.

\noindent\textbf{Proof-script interpretation.}
Proof-script reading. The machine proof decomposes the theorem into rewrites, program-logic obligations, relational equivalences, and arithmetic side conditions.
\emph{Main tactics visible in the proof script:} \ec{rewrite}.
\emph{Named identifiers syntactically cited in the proof block:}
\begin{lstlisting}[language=EasyCrypt,basicstyle=\ttfamily\scriptsize]
signing_simulation_relation
\end{lstlisting}

\end{formaltheorem}

\begin{formaltheorem}[EasyCrypt line 542]
\noindent\textbf{EasyCrypt name.}
\begin{lstlisting}[language=EasyCrypt,basicstyle=\ttfamily\scriptsize]
signing_attempt_simulates
\end{lstlisting}
\noindent\textbf{Checked EasyCrypt statement.}
\begin{lstlisting}[language=EasyCrypt,basicstyle=\ttfamily\scriptsize]
lemma signing_attempt_simulates md pk sk m ctx coins sig :
  signing_attempt_relation md pk sk m ctx coins sig =>
  signing_simulation_relation md pk m ctx sig.
\end{lstlisting}
\noindent\textbf{Formal mathematical reading.}
Mathematical theorem. This is an implication, normally turning one invariant, validity predicate, or proof obligation into another.

\noindent\textbf{Role in the proof.}
Use in this chapter. It proves an exact game replacement or simulator equivalence.

\noindent\textbf{Proof-script interpretation.}
Proof-script reading. The machine proof decomposes the theorem into rewrites, program-logic obligations, relational equivalences, and arithmetic side conditions.
\emph{Main tactics visible in the proof script:} \ec{rewrite}, \ec{split}, \ec{apply}.
\emph{Named identifiers syntactically cited in the proof block:}
\begin{lstlisting}[language=EasyCrypt,basicstyle=\ttfamily\scriptsize]
coins
ctx
md
pk
sig
signing_attempt_verifies
signing_simulation_relation
sk
\end{lstlisting}

\end{formaltheorem}

\begin{formaltheorem}[EasyCrypt line 554]
\noindent\textbf{EasyCrypt name.}
\begin{lstlisting}[language=EasyCrypt,basicstyle=\ttfamily\scriptsize]
sign_internal_attempt_relation
\end{lstlisting}
\noindent\textbf{Checked EasyCrypt statement.}
\begin{lstlisting}[language=EasyCrypt,basicstyle=\ttfamily\scriptsize]
lemma sign_internal_attempt_relation md sk m ctx coins :
  signing_attempt_relation md
    (public_key_of_secret md sk) sk m ctx coins
    (sign_internal md sk m ctx coins).
\end{lstlisting}
\noindent\textbf{Formal mathematical reading.}
Mathematical theorem. This lemma is a universally quantified proposition over the variables shown in the EasyCrypt statement.

\noindent\textbf{Role in the proof.}
Use in this chapter. It is a reusable checked theorem cited by later proof scripts.

\noindent\textbf{Proof-script interpretation.}
Proof-script reading. The machine proof decomposes the theorem into rewrites, program-logic obligations, relational equivalences, and arithmetic side conditions.
\emph{Main tactics visible in the proof script:} \ec{rewrite}, \ec{split}, \ec{smt}.
\emph{Named identifiers syntactically cited in the proof block:}
\begin{lstlisting}[language=EasyCrypt,basicstyle=\ttfamily\scriptsize]
response_norm_ok_current
signing_accepts
signing_attempt_relation
signing_attempt_state_accepts
signing_attempt_state_hint_norm_accepts
signing_attempt_state_response_norm_accepts
signing_attempt_state_valid_signature_accepts
trivial
valid_keypair
valid_signature_sign_internal
verify_norm_ok_current
\end{lstlisting}

\end{formaltheorem}

\begin{formaltheorem}[EasyCrypt line 571]
\noindent\textbf{EasyCrypt name.}
\begin{lstlisting}[language=EasyCrypt,basicstyle=\ttfamily\scriptsize]
signing_attempt_sample_of_pair_currentE
\end{lstlisting}
\noindent\textbf{Checked EasyCrypt statement.}
\begin{lstlisting}[language=EasyCrypt,basicstyle=\ttfamily\scriptsize]
lemma signing_attempt_sample_of_pair_currentE md sk m ctx coins :
  signing_attempt_sample_of_pair md sk m ctx coins
    (signing_sample_pair_of_coins md sk m ctx coins) =
  ( signing_sample_y1 md sk m ctx coins,
    signing_sample_y2 md sk m ctx coins,
    commitment_raw md sk m ctx coins).
\end{lstlisting}
\noindent\textbf{Formal mathematical reading.}
Mathematical theorem. This is an equality or definitional expansion used for rewriting and normalization.

\noindent\textbf{Role in the proof.}
Use in this chapter. It contributes directly to an advantage bound, bad-event bound, or real-arithmetic composition step.

\noindent\textbf{Proof-script interpretation.}
Proof-script reading. The machine proof decomposes the theorem into rewrites, program-logic obligations, relational equivalences, and arithmetic side conditions.
\emph{Main tactics visible in the proof script:} \ec{rewrite}.
\emph{Named identifiers syntactically cited in the proof block:}
\begin{lstlisting}[language=EasyCrypt,basicstyle=\ttfamily\scriptsize]
signing_attempt_sample_of_pair
signing_sample_pair_of_coins
signing_sample_pair_y1
signing_sample_pair_y2
\end{lstlisting}

\end{formaltheorem}

\begin{formaltheorem}[EasyCrypt line 583]
\noindent\textbf{EasyCrypt name.}
\begin{lstlisting}[language=EasyCrypt,basicstyle=\ttfamily\scriptsize]
signing_attempt_state_of_sample_pair_currentE
\end{lstlisting}
\noindent\textbf{Checked EasyCrypt statement.}
\begin{lstlisting}[language=EasyCrypt,basicstyle=\ttfamily\scriptsize]
lemma signing_attempt_state_of_sample_pair_currentE md sk m ctx coins :
  signing_attempt_state_of_sample_pair md sk m ctx coins
    (signing_sample_pair_of_coins md sk m ctx coins) =
  signing_attempt_state_of_coins md sk m ctx coins.
\end{lstlisting}
\noindent\textbf{Formal mathematical reading.}
Mathematical theorem. This is an equality or definitional expansion used for rewriting and normalization.

\noindent\textbf{Role in the proof.}
Use in this chapter. It contributes directly to an advantage bound, bad-event bound, or real-arithmetic composition step.

\noindent\textbf{Proof-script interpretation.}
Proof-script reading. The machine proof decomposes the theorem into rewrites, program-logic obligations, relational equivalences, and arithmetic side conditions.
\emph{Main tactics visible in the proof script:} \ec{rewrite}.
\emph{Named identifiers syntactically cited in the proof block:}
\begin{lstlisting}[language=EasyCrypt,basicstyle=\ttfamily\scriptsize]
signing_attempt_sample_of_pair_currentE
signing_attempt_state_of_coins
signing_attempt_state_of_sample_pair
\end{lstlisting}

\end{formaltheorem}

\begin{formaltheorem}[EasyCrypt line 593]
\noindent\textbf{EasyCrypt name.}
\begin{lstlisting}[language=EasyCrypt,basicstyle=\ttfamily\scriptsize]
structural_signing_sample_pair_sampler_ok
\end{lstlisting}
\noindent\textbf{Checked EasyCrypt statement.}
\begin{lstlisting}[language=EasyCrypt,basicstyle=\ttfamily\scriptsize]
lemma structural_signing_sample_pair_sampler_ok md sk m ctx :
  signing_sample_pair_sampler_ok md
    (dstructural_signing_sample_pair md sk m ctx).
\end{lstlisting}
\noindent\textbf{Formal mathematical reading.}
Mathematical theorem. This lemma is a universally quantified proposition over the variables shown in the EasyCrypt statement.

\noindent\textbf{Role in the proof.}
Use in this chapter. It contributes directly to an advantage bound, bad-event bound, or real-arithmetic composition step.

\noindent\textbf{Proof-script interpretation.}
Proof-script reading. The machine proof decomposes the theorem into rewrites, program-logic obligations, relational equivalences, and arithmetic side conditions.
\emph{Main tactics visible in the proof script:} \ec{rewrite}, \ec{apply}.
\emph{Named identifiers syntactically cited in the proof block:}
\begin{lstlisting}[language=EasyCrypt,basicstyle=\ttfamily\scriptsize]
coins
signing_sample_pair_sampler_ok
structural_signing_sample_pair_distribution_ok
\end{lstlisting}

\end{formaltheorem}

\begin{formaltheorem}[EasyCrypt line 601]
\noindent\textbf{EasyCrypt name.}
\begin{lstlisting}[language=EasyCrypt,basicstyle=\ttfamily\scriptsize]
bounded_unit_signing_sample_pair_sampler_ok
\end{lstlisting}
\noindent\textbf{Checked EasyCrypt statement.}
\begin{lstlisting}[language=EasyCrypt,basicstyle=\ttfamily\scriptsize]
lemma bounded_unit_signing_sample_pair_sampler_ok md :
  signing_sample_pair_sampler_ok md
    (fun _ => dbounded_unit_signing_sample_pair md).
\end{lstlisting}
\noindent\textbf{Formal mathematical reading.}
Mathematical theorem. This is an implication, normally turning one invariant, validity predicate, or proof obligation into another.

\noindent\textbf{Role in the proof.}
Use in this chapter. It contributes directly to an advantage bound, bad-event bound, or real-arithmetic composition step.

\noindent\textbf{Proof-script interpretation.}
Proof-script reading. The machine proof decomposes the theorem into rewrites, program-logic obligations, relational equivalences, and arithmetic side conditions.
\emph{Main tactics visible in the proof script:} \ec{rewrite}, \ec{apply}.
\emph{Named identifiers syntactically cited in the proof block:}
\begin{lstlisting}[language=EasyCrypt,basicstyle=\ttfamily\scriptsize]
bounded_unit_signing_sample_pair_distribution_ok
coins
signing_sample_pair_sampler_ok
\end{lstlisting}

\end{formaltheorem}

\begin{formaltheorem}[EasyCrypt line 609]
\noindent\textbf{EasyCrypt name.}
\begin{lstlisting}[language=EasyCrypt,basicstyle=\ttfamily\scriptsize]
checked_hyperball_signing_sample_pair_sampler_ok
\end{lstlisting}
\noindent\textbf{Checked EasyCrypt statement.}
\begin{lstlisting}[language=EasyCrypt,basicstyle=\ttfamily\scriptsize]
lemma checked_hyperball_signing_sample_pair_sampler_ok md :
  signing_sample_pair_sampler_ok md
    (fun _ => dchecked_hyperball_signing_sample_pair md).
\end{lstlisting}
\noindent\textbf{Formal mathematical reading.}
Mathematical theorem. This is an implication, normally turning one invariant, validity predicate, or proof obligation into another.

\noindent\textbf{Role in the proof.}
Use in this chapter. It contributes directly to an advantage bound, bad-event bound, or real-arithmetic composition step.

\noindent\textbf{Proof-script interpretation.}
Proof-script reading. The machine proof decomposes the theorem into rewrites, program-logic obligations, relational equivalences, and arithmetic side conditions.
\emph{Main tactics visible in the proof script:} \ec{rewrite}, \ec{apply}.
\emph{Named identifiers syntactically cited in the proof block:}
\begin{lstlisting}[language=EasyCrypt,basicstyle=\ttfamily\scriptsize]
checked_hyperball_signing_sample_pair_distribution_ok
coins
signing_sample_pair_sampler_ok
\end{lstlisting}

\end{formaltheorem}

\begin{formaltheorem}[EasyCrypt line 617]
\noindent\textbf{EasyCrypt name.}
\begin{lstlisting}[language=EasyCrypt,basicstyle=\ttfamily\scriptsize]
exact_hyperball_signing_sample_pair_sampler_ok
\end{lstlisting}
\noindent\textbf{Checked EasyCrypt statement.}
\begin{lstlisting}[language=EasyCrypt,basicstyle=\ttfamily\scriptsize]
lemma exact_hyperball_signing_sample_pair_sampler_ok md :
  signing_sample_pair_sampler_ok md
    (fun _ => dexact_hyperball_signing_sample_pair md).
\end{lstlisting}
\noindent\textbf{Formal mathematical reading.}
Mathematical theorem. This is an implication, normally turning one invariant, validity predicate, or proof obligation into another.

\noindent\textbf{Role in the proof.}
Use in this chapter. It contributes directly to an advantage bound, bad-event bound, or real-arithmetic composition step.

\noindent\textbf{Proof-script interpretation.}
Proof-script reading. The machine proof decomposes the theorem into rewrites, program-logic obligations, relational equivalences, and arithmetic side conditions.
\emph{Main tactics visible in the proof script:} \ec{rewrite}, \ec{apply}.
\emph{Named identifiers syntactically cited in the proof block:}
\begin{lstlisting}[language=EasyCrypt,basicstyle=\ttfamily\scriptsize]
coins
exact_hyperball_signing_sample_pair_distribution_ok
signing_sample_pair_sampler_ok
\end{lstlisting}

\end{formaltheorem}

\begin{formaltheorem}[EasyCrypt line 625]
\noindent\textbf{EasyCrypt name.}
\begin{lstlisting}[language=EasyCrypt,basicstyle=\ttfamily\scriptsize]
dsigning_attempt_state_from_sample_pair_sampler_lossless
\end{lstlisting}
\noindent\textbf{Checked EasyCrypt statement.}
\begin{lstlisting}[language=EasyCrypt,basicstyle=\ttfamily\scriptsize]
lemma dsigning_attempt_state_from_sample_pair_sampler_lossless
   md sk m ctx dsample :
  signing_sample_pair_sampler_ok md dsample =>
  is_lossless
    (dsigning_attempt_state_from_sample_pair_sampler
      md sk m ctx dsample).
\end{lstlisting}
\noindent\textbf{Formal mathematical reading.}
Mathematical theorem. This is an implication, normally turning one invariant, validity predicate, or proof obligation into another.

\noindent\textbf{Role in the proof.}
Use in this chapter. It contributes directly to an advantage bound, bad-event bound, or real-arithmetic composition step.

\noindent\textbf{Proof-script interpretation.}
Proof-script reading. The machine proof decomposes the theorem into rewrites, program-logic obligations, relational equivalences, and arithmetic side conditions.
\emph{Main tactics visible in the proof script:} \ec{rewrite}, \ec{apply}.
\emph{Named identifiers syntactically cited in the proof block:}
\begin{lstlisting}[language=EasyCrypt,basicstyle=\ttfamily\scriptsize]
coins
coins_supp
dlet_ll
dmap_ll
dsigning_attempt_state_from_sample_pair_sampler
hll
hsample
signing_coin_distribution_domain
signing_coin_distribution_lossless
signing_sample_pair_distribution_ok
\end{lstlisting}

\end{formaltheorem}

\begin{formaltheorem}[EasyCrypt line 643]
\noindent\textbf{EasyCrypt name.}
\begin{lstlisting}[language=EasyCrypt,basicstyle=\ttfamily\scriptsize]
dsigning_attempt_state_from_sample_pair_sampler_sample_ok
\end{lstlisting}
\noindent\textbf{Checked EasyCrypt statement.}
\begin{lstlisting}[language=EasyCrypt,basicstyle=\ttfamily\scriptsize]
lemma dsigning_attempt_state_from_sample_pair_sampler_sample_ok
   md sk m ctx dsample st :
  signing_sample_pair_sampler_ok md dsample =>
  st \in
    dsigning_attempt_state_from_sample_pair_sampler
      md sk m ctx dsample =>
  signing_attempt_state_sample_pair_wf md st /\
  signing_attempt_state_sample_pair_bounded md st.
\end{lstlisting}
\noindent\textbf{Formal mathematical reading.}
Mathematical theorem. This is an implication, normally turning one invariant, validity predicate, or proof obligation into another.

\noindent\textbf{Role in the proof.}
Use in this chapter. It contributes directly to an advantage bound, bad-event bound, or real-arithmetic composition step.

\noindent\textbf{Proof-script interpretation.}
Proof-script reading. The machine proof decomposes the theorem into rewrites, program-logic obligations, relational equivalences, and arithmetic side conditions.
\emph{Main tactics visible in the proof script:} \ec{rewrite}, \ec{apply}, \ec{split}.
\emph{Named identifiers syntactically cited in the proof block:}
\begin{lstlisting}[language=EasyCrypt,basicstyle=\ttfamily\scriptsize]
coins
coins_supp
dsigning_attempt_state_from_sample_pair_sampler
hsample
hsupp
sample
sample_bounded
sample_supp
sample_wf
signing_attempt_sample_of_pair
signing_attempt_sample_y1
signing_attempt_sample_y2
signing_attempt_state_of_sample_pair
signing_attempt_state_sample
signing_attempt_state_sample_pair_bounded
signing_attempt_state_sample_pair_wf
signing_attempt_state_sampled_y1
signing_attempt_state_sampled_y2
signing_coin_distribution_domain
signing_sample_pair_bounded
signing_sample_pair_distribution_ok
signing_sample_pair_sample_ok
signing_sample_pair_wf
signing_sample_pair_y1
signing_sample_pair_y2
supp_dlet
supp_dmap
\end{lstlisting}

\end{formaltheorem}

\begin{formaltheorem}[EasyCrypt line 678]
\noindent\textbf{EasyCrypt name.}
\begin{lstlisting}[language=EasyCrypt,basicstyle=\ttfamily\scriptsize]
dstructural_signing_attempt_state_lossless
\end{lstlisting}
\noindent\textbf{Checked EasyCrypt statement.}
\begin{lstlisting}[language=EasyCrypt,basicstyle=\ttfamily\scriptsize]
lemma dstructural_signing_attempt_state_lossless md sk m ctx :
  is_lossless (dstructural_signing_attempt_state md sk m ctx).
\end{lstlisting}
\noindent\textbf{Formal mathematical reading.}
Mathematical theorem. This lemma is a universally quantified proposition over the variables shown in the EasyCrypt statement.

\noindent\textbf{Role in the proof.}
Use in this chapter. It contributes directly to an advantage bound, bad-event bound, or real-arithmetic composition step.

\noindent\textbf{Proof-script interpretation.}
Proof-script reading. The machine proof decomposes the theorem into rewrites, program-logic obligations, relational equivalences, and arithmetic side conditions.
\emph{Main tactics visible in the proof script:} \ec{rewrite}, \ec{apply}.
\emph{Named identifiers syntactically cited in the proof block:}
\begin{lstlisting}[language=EasyCrypt,basicstyle=\ttfamily\scriptsize]
dmap_ll
dstructural_signing_attempt_state
signing_coin_distribution_lossless
\end{lstlisting}

\end{formaltheorem}

\begin{formaltheorem}[EasyCrypt line 685]
\noindent\textbf{EasyCrypt name.}
\begin{lstlisting}[language=EasyCrypt,basicstyle=\ttfamily\scriptsize]
dstructural_signing_attempt_state_from_sampler_lossless
\end{lstlisting}
\noindent\textbf{Checked EasyCrypt statement.}
\begin{lstlisting}[language=EasyCrypt,basicstyle=\ttfamily\scriptsize]
lemma dstructural_signing_attempt_state_from_sampler_lossless
   md sk m ctx :
  is_lossless
    (dstructural_signing_attempt_state_from_sampler md sk m ctx).
\end{lstlisting}
\noindent\textbf{Formal mathematical reading.}
Mathematical theorem. This lemma is a universally quantified proposition over the variables shown in the EasyCrypt statement.

\noindent\textbf{Role in the proof.}
Use in this chapter. It contributes directly to an advantage bound, bad-event bound, or real-arithmetic composition step.

\noindent\textbf{Proof-script interpretation.}
Proof-script reading. The machine proof decomposes the theorem into rewrites, program-logic obligations, relational equivalences, and arithmetic side conditions.
\emph{Main tactics visible in the proof script:} \ec{rewrite}, \ec{apply}.
\emph{Named identifiers syntactically cited in the proof block:}
\begin{lstlisting}[language=EasyCrypt,basicstyle=\ttfamily\scriptsize]
dsigning_attempt_state_from_sample_pair_sampler_lossless
dstructural_signing_attempt_state_from_sampler
structural_signing_sample_pair_sampler_ok
\end{lstlisting}

\end{formaltheorem}

\begin{formaltheorem}[EasyCrypt line 695]
\noindent\textbf{EasyCrypt name.}
\begin{lstlisting}[language=EasyCrypt,basicstyle=\ttfamily\scriptsize]
dbounded_unit_signing_attempt_state_lossless
\end{lstlisting}
\noindent\textbf{Checked EasyCrypt statement.}
\begin{lstlisting}[language=EasyCrypt,basicstyle=\ttfamily\scriptsize]
lemma dbounded_unit_signing_attempt_state_lossless md sk m ctx :
  is_lossless (dbounded_unit_signing_attempt_state md sk m ctx).
\end{lstlisting}
\noindent\textbf{Formal mathematical reading.}
Mathematical theorem. This lemma is a universally quantified proposition over the variables shown in the EasyCrypt statement.

\noindent\textbf{Role in the proof.}
Use in this chapter. It contributes directly to an advantage bound, bad-event bound, or real-arithmetic composition step.

\noindent\textbf{Proof-script interpretation.}
Proof-script reading. The machine proof decomposes the theorem into rewrites, program-logic obligations, relational equivalences, and arithmetic side conditions.
\emph{Main tactics visible in the proof script:} \ec{rewrite}, \ec{apply}.
\emph{Named identifiers syntactically cited in the proof block:}
\begin{lstlisting}[language=EasyCrypt,basicstyle=\ttfamily\scriptsize]
bounded_unit_signing_sample_pair_sampler_ok
dbounded_unit_signing_attempt_state
dsigning_attempt_state_from_sample_pair_sampler_lossless
\end{lstlisting}

\end{formaltheorem}

\begin{formaltheorem}[EasyCrypt line 703]
\noindent\textbf{EasyCrypt name.}
\begin{lstlisting}[language=EasyCrypt,basicstyle=\ttfamily\scriptsize]
dbounded_unit_signing_attempt_state_sample_ok
\end{lstlisting}
\noindent\textbf{Checked EasyCrypt statement.}
\begin{lstlisting}[language=EasyCrypt,basicstyle=\ttfamily\scriptsize]
lemma dbounded_unit_signing_attempt_state_sample_ok md sk m ctx st :
  st \in dbounded_unit_signing_attempt_state md sk m ctx =>
  signing_attempt_state_sample_pair_wf md st /\
  signing_attempt_state_sample_pair_bounded md st.
\end{lstlisting}
\noindent\textbf{Formal mathematical reading.}
Mathematical theorem. This is an implication, normally turning one invariant, validity predicate, or proof obligation into another.

\noindent\textbf{Role in the proof.}
Use in this chapter. It contributes directly to an advantage bound, bad-event bound, or real-arithmetic composition step.

\noindent\textbf{Proof-script interpretation.}
Proof-script reading. The machine proof decomposes the theorem into rewrites, program-logic obligations, relational equivalences, and arithmetic side conditions.
\emph{Main tactics visible in the proof script:} \ec{rewrite}, \ec{apply}.
\emph{Named identifiers syntactically cited in the proof block:}
\begin{lstlisting}[language=EasyCrypt,basicstyle=\ttfamily\scriptsize]
bounded_unit_signing_sample_pair_sampler_ok
dbounded_unit_signing_attempt_state
dsigning_attempt_state_from_sample_pair_sampler_sample_ok
\end{lstlisting}

\end{formaltheorem}

\begin{formaltheorem}[EasyCrypt line 713]
\noindent\textbf{EasyCrypt name.}
\begin{lstlisting}[language=EasyCrypt,basicstyle=\ttfamily\scriptsize]
dchecked_hyperball_signing_attempt_state_lossless
\end{lstlisting}
\noindent\textbf{Checked EasyCrypt statement.}
\begin{lstlisting}[language=EasyCrypt,basicstyle=\ttfamily\scriptsize]
lemma dchecked_hyperball_signing_attempt_state_lossless md sk m ctx :
  is_lossless (dchecked_hyperball_signing_attempt_state md sk m ctx).
\end{lstlisting}
\noindent\textbf{Formal mathematical reading.}
Mathematical theorem. This lemma is a universally quantified proposition over the variables shown in the EasyCrypt statement.

\noindent\textbf{Role in the proof.}
Use in this chapter. It contributes directly to an advantage bound, bad-event bound, or real-arithmetic composition step.

\noindent\textbf{Proof-script interpretation.}
Proof-script reading. The machine proof decomposes the theorem into rewrites, program-logic obligations, relational equivalences, and arithmetic side conditions.
\emph{Main tactics visible in the proof script:} \ec{rewrite}, \ec{apply}.
\emph{Named identifiers syntactically cited in the proof block:}
\begin{lstlisting}[language=EasyCrypt,basicstyle=\ttfamily\scriptsize]
checked_hyperball_signing_sample_pair_sampler_ok
dchecked_hyperball_signing_attempt_state
dsigning_attempt_state_from_sample_pair_sampler_lossless
\end{lstlisting}

\end{formaltheorem}

\begin{formaltheorem}[EasyCrypt line 721]
\noindent\textbf{EasyCrypt name.}
\begin{lstlisting}[language=EasyCrypt,basicstyle=\ttfamily\scriptsize]
dchecked_hyperball_signing_attempt_state_sample_ok
\end{lstlisting}
\noindent\textbf{Checked EasyCrypt statement.}
\begin{lstlisting}[language=EasyCrypt,basicstyle=\ttfamily\scriptsize]
lemma dchecked_hyperball_signing_attempt_state_sample_ok md sk m ctx st :
  st \in dchecked_hyperball_signing_attempt_state md sk m ctx =>
  signing_attempt_state_sample_pair_wf md st /\
  signing_attempt_state_sample_pair_bounded md st.
\end{lstlisting}
\noindent\textbf{Formal mathematical reading.}
Mathematical theorem. This is an implication, normally turning one invariant, validity predicate, or proof obligation into another.

\noindent\textbf{Role in the proof.}
Use in this chapter. It contributes directly to an advantage bound, bad-event bound, or real-arithmetic composition step.

\noindent\textbf{Proof-script interpretation.}
Proof-script reading. The machine proof decomposes the theorem into rewrites, program-logic obligations, relational equivalences, and arithmetic side conditions.
\emph{Main tactics visible in the proof script:} \ec{rewrite}, \ec{apply}.
\emph{Named identifiers syntactically cited in the proof block:}
\begin{lstlisting}[language=EasyCrypt,basicstyle=\ttfamily\scriptsize]
checked_hyperball_signing_sample_pair_sampler_ok
dchecked_hyperball_signing_attempt_state
dsigning_attempt_state_from_sample_pair_sampler_sample_ok
\end{lstlisting}

\end{formaltheorem}

\begin{formaltheorem}[EasyCrypt line 731]
\noindent\textbf{EasyCrypt name.}
\begin{lstlisting}[language=EasyCrypt,basicstyle=\ttfamily\scriptsize]
dexact_hyperball_signing_attempt_state_lossless
\end{lstlisting}
\noindent\textbf{Checked EasyCrypt statement.}
\begin{lstlisting}[language=EasyCrypt,basicstyle=\ttfamily\scriptsize]
lemma dexact_hyperball_signing_attempt_state_lossless md sk m ctx :
  is_lossless (dexact_hyperball_signing_attempt_state md sk m ctx).
\end{lstlisting}
\noindent\textbf{Formal mathematical reading.}
Mathematical theorem. This lemma is a universally quantified proposition over the variables shown in the EasyCrypt statement.

\noindent\textbf{Role in the proof.}
Use in this chapter. It contributes directly to an advantage bound, bad-event bound, or real-arithmetic composition step.

\noindent\textbf{Proof-script interpretation.}
Proof-script reading. The machine proof decomposes the theorem into rewrites, program-logic obligations, relational equivalences, and arithmetic side conditions.
\emph{Main tactics visible in the proof script:} \ec{rewrite}, \ec{apply}.
\emph{Named identifiers syntactically cited in the proof block:}
\begin{lstlisting}[language=EasyCrypt,basicstyle=\ttfamily\scriptsize]
dexact_hyperball_signing_attempt_state
dsigning_attempt_state_from_sample_pair_sampler_lossless
exact_hyperball_signing_sample_pair_sampler_ok
\end{lstlisting}

\end{formaltheorem}

\begin{formaltheorem}[EasyCrypt line 739]
\noindent\textbf{EasyCrypt name.}
\begin{lstlisting}[language=EasyCrypt,basicstyle=\ttfamily\scriptsize]
dexact_hyperball_signing_attempt_state_sample_ok
\end{lstlisting}
\noindent\textbf{Checked EasyCrypt statement.}
\begin{lstlisting}[language=EasyCrypt,basicstyle=\ttfamily\scriptsize]
lemma dexact_hyperball_signing_attempt_state_sample_ok md sk m ctx st :
  st \in dexact_hyperball_signing_attempt_state md sk m ctx =>
  signing_attempt_state_sample_pair_wf md st /\
  signing_attempt_state_sample_pair_bounded md st.
\end{lstlisting}
\noindent\textbf{Formal mathematical reading.}
Mathematical theorem. This is an implication, normally turning one invariant, validity predicate, or proof obligation into another.

\noindent\textbf{Role in the proof.}
Use in this chapter. It contributes directly to an advantage bound, bad-event bound, or real-arithmetic composition step.

\noindent\textbf{Proof-script interpretation.}
Proof-script reading. The machine proof decomposes the theorem into rewrites, program-logic obligations, relational equivalences, and arithmetic side conditions.
\emph{Main tactics visible in the proof script:} \ec{rewrite}, \ec{apply}.
\emph{Named identifiers syntactically cited in the proof block:}
\begin{lstlisting}[language=EasyCrypt,basicstyle=\ttfamily\scriptsize]
dexact_hyperball_signing_attempt_state
dsigning_attempt_state_from_sample_pair_sampler_sample_ok
exact_hyperball_signing_sample_pair_sampler_ok
\end{lstlisting}

\end{formaltheorem}

\begin{formaltheorem}[EasyCrypt line 749]
\noindent\textbf{EasyCrypt name.}
\begin{lstlisting}[language=EasyCrypt,basicstyle=\ttfamily\scriptsize]
dexact_hyperball_signing_attempt_state_supportP
\end{lstlisting}
\noindent\textbf{Checked EasyCrypt statement.}
\begin{lstlisting}[language=EasyCrypt,basicstyle=\ttfamily\scriptsize]
lemma dexact_hyperball_signing_attempt_state_supportP md sk m ctx st :
  st \in dexact_hyperball_signing_attempt_state md sk m ctx <=>
  exists (coins : random_coins) (sample : signing_sample_pair),
    coins \in drandom_coins /\
    hyperball_sample_ok md sample /\
    st = signing_attempt_state_of_sample_pair md sk m ctx coins sample.
\end{lstlisting}
\noindent\textbf{Formal mathematical reading.}
Mathematical theorem. This is an inequality, usually an arithmetic loss bound, event inclusion bound, or monotonicity fact used to compose probabilities.

\noindent\textbf{Role in the proof.}
Use in this chapter. It proves an exact game replacement or simulator equivalence.

\noindent\textbf{Proof-script interpretation.}
Proof-script reading. The machine proof decomposes the theorem into rewrites, program-logic obligations, relational equivalences, and arithmetic side conditions.
\emph{Main tactics visible in the proof script:} \ec{rewrite}, \ec{split}.
\emph{Named identifiers syntactically cited in the proof block:}
\begin{lstlisting}[language=EasyCrypt,basicstyle=\ttfamily\scriptsize]
coins
coins_supp
dexact_hyperball_signing_attempt_state
dsigning_attempt_state_from_sample_pair_sampler
exact_hyperball_signing_sample_pair_supportP
sample
sample_ok
sample_supp
stE
supp_dlet
supp_dmap
\end{lstlisting}

\end{formaltheorem}

\begin{formaltheorem}[EasyCrypt line 778]
\noindent\textbf{EasyCrypt name.}
\begin{lstlisting}[language=EasyCrypt,basicstyle=\ttfamily\scriptsize]
dexact_hyperball_signing_attempt_state_reference_rejection_abort_bound1
\end{lstlisting}
\noindent\textbf{Checked EasyCrypt statement.}
\begin{lstlisting}[language=EasyCrypt,basicstyle=\ttfamily\scriptsize]
lemma dexact_hyperball_signing_attempt_state_reference_rejection_abort_bound1
  md sk m ctx :
  mu (dexact_hyperball_signing_attempt_state md sk m ctx)
    (signing_attempt_state_reference_rejection_aborts md) <= 1%r.
\end{lstlisting}
\noindent\textbf{Formal mathematical reading.}
Mathematical theorem. This is an inequality, usually an arithmetic loss bound, event inclusion bound, or monotonicity fact used to compose probabilities.

\noindent\textbf{Role in the proof.}
Use in this chapter. It contributes directly to an advantage bound, bad-event bound, or real-arithmetic composition step.

\noindent\textbf{Proof-script interpretation.}
Proof-script reading. The machine proof decomposes the theorem into rewrites, program-logic obligations, relational equivalences, and arithmetic side conditions.
\emph{Main tactics visible in the proof script:} \ec{smt}.
\emph{Named identifiers syntactically cited in the proof block:}
\begin{lstlisting}[language=EasyCrypt,basicstyle=\ttfamily\scriptsize]
mu_bounded
\end{lstlisting}

\end{formaltheorem}

\begin{formaltheorem}[EasyCrypt line 786]
\noindent\textbf{EasyCrypt name.}
\begin{lstlisting}[language=EasyCrypt,basicstyle=\ttfamily\scriptsize]
dexact_hyperball_signing_attempt_state_mu_pullback
\end{lstlisting}
\noindent\textbf{Checked EasyCrypt statement.}
\begin{lstlisting}[language=EasyCrypt,basicstyle=\ttfamily\scriptsize]
lemma dexact_hyperball_signing_attempt_state_mu_pullback
  md sk m ctx (P : signing_attempt_state -> bool) :
  mu (dexact_hyperball_signing_attempt_state md sk m ctx) P =
  mu (dlet drandom_coins
        (fun coins =>
           dmap (dexact_hyperball_signing_sample_pair md)
             (fun sample =>
                signing_attempt_state_of_sample_pair
                  md sk m ctx coins sample)))
    P.
\end{lstlisting}
\noindent\textbf{Formal mathematical reading.}
Mathematical theorem. This is an implication, normally turning one invariant, validity predicate, or proof obligation into another.

\noindent\textbf{Role in the proof.}
Use in this chapter. It proves an exact game replacement or simulator equivalence.

\noindent\textbf{Proof-script interpretation.}
Proof-script reading. The machine proof decomposes the theorem into rewrites, program-logic obligations, relational equivalences, and arithmetic side conditions.
\emph{Main tactics visible in the proof script:} \ec{rewrite}.
\emph{Named identifiers syntactically cited in the proof block:}
\begin{lstlisting}[language=EasyCrypt,basicstyle=\ttfamily\scriptsize]
dexact_hyperball_signing_attempt_state
dsigning_attempt_state_from_sample_pair_sampler
\end{lstlisting}

\end{formaltheorem}

\begin{formaltheorem}[EasyCrypt line 801]
\noindent\textbf{EasyCrypt name.}
\begin{lstlisting}[language=EasyCrypt,basicstyle=\ttfamily\scriptsize]
signing_attempt_state_reference_rejection_aborts_componentE
\end{lstlisting}
\noindent\textbf{Checked EasyCrypt statement.}
\begin{lstlisting}[language=EasyCrypt,basicstyle=\ttfamily\scriptsize]
lemma signing_attempt_state_reference_rejection_aborts_componentE md st :
  signing_attempt_state_reference_rejection_aborts md st =
  (signing_attempt_state_reject1_aborts md st \/
   signing_attempt_state_reject2_aborts md st \/
   signing_attempt_state_pack_aborts md st \/
   signing_attempt_state_hint_norm_aborts md st).
\end{lstlisting}
\noindent\textbf{Formal mathematical reading.}
Mathematical theorem. This is an equality or definitional expansion used for rewriting and normalization.

\noindent\textbf{Role in the proof.}
Use in this chapter. It preserves an invariant across an oracle call, adversary call, or game procedure.

\noindent\textbf{Proof-script interpretation.}
Proof-script reading. The machine proof decomposes the theorem into rewrites, program-logic obligations, relational equivalences, and arithmetic side conditions.
\emph{Main tactics visible in the proof script:} \ec{rewrite}, \ec{smt}.
\emph{Named identifiers syntactically cited in the proof block:}
\begin{lstlisting}[language=EasyCrypt,basicstyle=\ttfamily\scriptsize]
signing_attempt_state_hint_norm_aborts
signing_attempt_state_pack_aborts
signing_attempt_state_reference_rejection_aborts
signing_attempt_state_reference_rejection_accepts
signing_attempt_state_reject1_aborts
signing_attempt_state_reject2_accepts
\end{lstlisting}

\end{formaltheorem}

\begin{formaltheorem}[EasyCrypt line 817]
\noindent\textbf{EasyCrypt name.}
\begin{lstlisting}[language=EasyCrypt,basicstyle=\ttfamily\scriptsize]
dexact_hyperball_signing_attempt_state_reference_rejection_abort_split
\end{lstlisting}
\noindent\textbf{Checked EasyCrypt statement.}
\begin{lstlisting}[language=EasyCrypt,basicstyle=\ttfamily\scriptsize]
lemma dexact_hyperball_signing_attempt_state_reference_rejection_abort_split
  md sk m ctx :
  mu (dexact_hyperball_signing_attempt_state md sk m ctx)
    (signing_attempt_state_reference_rejection_aborts md) =
  mu (dexact_hyperball_signing_attempt_state md sk m ctx)
    (predU
      (signing_attempt_state_reject1_aborts md)
      (predU
        (signing_attempt_state_reject2_aborts md)
        (predU
          (signing_attempt_state_pack_aborts md)
          (signing_attempt_state_hint_norm_aborts md)))).
\end{lstlisting}
\noindent\textbf{Formal mathematical reading.}
Mathematical theorem. This is an equality or definitional expansion used for rewriting and normalization.

\noindent\textbf{Role in the proof.}
Use in this chapter. It proves an exact game replacement or simulator equivalence.

\noindent\textbf{Proof-script interpretation.}
Proof-script reading. The machine proof decomposes the theorem into rewrites, program-logic obligations, relational equivalences, and arithmetic side conditions.
\emph{Main tactics visible in the proof script:} \ec{apply}, \ec{rewrite}.
\emph{Named identifiers syntactically cited in the proof block:}
\begin{lstlisting}[language=EasyCrypt,basicstyle=\ttfamily\scriptsize]
mu_eq
predU
signing_attempt_state_reference_rejection_aborts_componentE
st
\end{lstlisting}

\end{formaltheorem}

\begin{formaltheorem}[EasyCrypt line 834]
\noindent\textbf{EasyCrypt name.}
\begin{lstlisting}[language=EasyCrypt,basicstyle=\ttfamily\scriptsize]
dexact_hyperball_signing_attempt_state_reference_rejection_abort_union_bound
\end{lstlisting}
\noindent\textbf{Checked EasyCrypt statement.}
\begin{lstlisting}[language=EasyCrypt,basicstyle=\ttfamily\scriptsize]
lemma dexact_hyperball_signing_attempt_state_reference_rejection_abort_union_bound
  md sk m ctx :
  mu (dexact_hyperball_signing_attempt_state md sk m ctx)
    (signing_attempt_state_reference_rejection_aborts md) <=
  mu (dexact_hyperball_signing_attempt_state md sk m ctx)
    (signing_attempt_state_reject1_aborts md) +
  mu (dexact_hyperball_signing_attempt_state md sk m ctx)
    (signing_attempt_state_reject2_aborts md) +
  mu (dexact_hyperball_signing_attempt_state md sk m ctx)
    (signing_attempt_state_pack_aborts md) +
  mu (dexact_hyperball_signing_attempt_state md sk m ctx)
    (signing_attempt_state_hint_norm_aborts md).
\end{lstlisting}
\noindent\textbf{Formal mathematical reading.}
Mathematical theorem. This is an inequality, usually an arithmetic loss bound, event inclusion bound, or monotonicity fact used to compose probabilities.

\noindent\textbf{Role in the proof.}
Use in this chapter. It contributes directly to an advantage bound, bad-event bound, or real-arithmetic composition step.

\noindent\textbf{Proof-script interpretation.}
Proof-script reading. The machine proof decomposes the theorem into rewrites, program-logic obligations, relational equivalences, and arithmetic side conditions.
\emph{Main tactics visible in the proof script:} \ec{rewrite}, \ec{smt}.
\emph{Named identifiers syntactically cited in the proof block:}
\begin{lstlisting}[language=EasyCrypt,basicstyle=\ttfamily\scriptsize]
dexact_hyperball_signing_attempt_state_reference_rejection_abort_split
mu_bounded
mu_or
\end{lstlisting}

\end{formaltheorem}

\begin{formaltheorem}[EasyCrypt line 852]
\noindent\textbf{EasyCrypt name.}
\begin{lstlisting}[language=EasyCrypt,basicstyle=\ttfamily\scriptsize]
signing_attempt_state_of_sample_pair_signatureE
\end{lstlisting}
\noindent\textbf{Checked EasyCrypt statement.}
\begin{lstlisting}[language=EasyCrypt,basicstyle=\ttfamily\scriptsize]
lemma signing_attempt_state_of_sample_pair_signatureE
  md sk m ctx coins sample :
  signing_attempt_state_signature
    (signing_attempt_state_of_sample_pair md sk m ctx coins sample) =
  sign_internal md sk m ctx coins.
\end{lstlisting}
\noindent\textbf{Formal mathematical reading.}
Mathematical theorem. This is an equality or definitional expansion used for rewriting and normalization.

\noindent\textbf{Role in the proof.}
Use in this chapter. It contributes directly to an advantage bound, bad-event bound, or real-arithmetic composition step.

\noindent\textbf{Proof-script interpretation.}
Proof-script reading. The machine proof decomposes the theorem into rewrites, program-logic obligations, relational equivalences, and arithmetic side conditions.
\emph{Main tactics visible in the proof script:} \ec{rewrite}.
\emph{Named identifiers syntactically cited in the proof block:}
\begin{lstlisting}[language=EasyCrypt,basicstyle=\ttfamily\scriptsize]
signing_attempt_state_of_sample_pair
signing_attempt_state_signature
\end{lstlisting}

\end{formaltheorem}

\begin{formaltheorem}[EasyCrypt line 862]
\noindent\textbf{EasyCrypt name.}
\begin{lstlisting}[language=EasyCrypt,basicstyle=\ttfamily\scriptsize]
signing_attempt_state_of_sample_pair_pack_accepts
\end{lstlisting}
\noindent\textbf{Checked EasyCrypt statement.}
\begin{lstlisting}[language=EasyCrypt,basicstyle=\ttfamily\scriptsize]
lemma signing_attempt_state_of_sample_pair_pack_accepts
  md sk m ctx coins sample :
  signing_attempt_state_pack_accepts md
    (signing_attempt_state_of_sample_pair md sk m ctx coins sample).
\end{lstlisting}
\noindent\textbf{Formal mathematical reading.}
Mathematical theorem. This lemma is a universally quantified proposition over the variables shown in the EasyCrypt statement.

\noindent\textbf{Role in the proof.}
Use in this chapter. It contributes directly to an advantage bound, bad-event bound, or real-arithmetic composition step.

\noindent\textbf{Proof-script interpretation.}
Proof-script reading. The machine proof decomposes the theorem into rewrites, program-logic obligations, relational equivalences, and arithmetic side conditions.
\emph{Main tactics visible in the proof script:} \ec{rewrite}, \ec{split}, \ec{apply}.
\emph{Named identifiers syntactically cited in the proof block:}
\begin{lstlisting}[language=EasyCrypt,basicstyle=\ttfamily\scriptsize]
mode_crypto_bytes_gt0
signing_attempt_state_of_sample_pair_signatureE
signing_attempt_state_pack_accepts
signing_attempt_state_valid_signature_accepts
valid_signature_sign_internal
\end{lstlisting}

\end{formaltheorem}

\begin{formaltheorem}[EasyCrypt line 875]
\noindent\textbf{EasyCrypt name.}
\begin{lstlisting}[language=EasyCrypt,basicstyle=\ttfamily\scriptsize]
signing_attempt_state_of_sample_pair_pack_no_abort
\end{lstlisting}
\noindent\textbf{Checked EasyCrypt statement.}
\begin{lstlisting}[language=EasyCrypt,basicstyle=\ttfamily\scriptsize]
lemma signing_attempt_state_of_sample_pair_pack_no_abort
  md sk m ctx coins sample :
  ! signing_attempt_state_pack_aborts md
      (signing_attempt_state_of_sample_pair md sk m ctx coins sample).
\end{lstlisting}
\noindent\textbf{Formal mathematical reading.}
Mathematical theorem. This lemma is a universally quantified proposition over the variables shown in the EasyCrypt statement.

\noindent\textbf{Role in the proof.}
Use in this chapter. It contributes directly to an advantage bound, bad-event bound, or real-arithmetic composition step.

\noindent\textbf{Proof-script interpretation.}
Proof-script reading. The machine proof decomposes the theorem into rewrites, program-logic obligations, relational equivalences, and arithmetic side conditions.
\emph{Main tactics visible in the proof script:} \ec{rewrite}.
\emph{Named identifiers syntactically cited in the proof block:}
\begin{lstlisting}[language=EasyCrypt,basicstyle=\ttfamily\scriptsize]
signing_attempt_state_of_sample_pair_pack_accepts
signing_attempt_state_pack_aborts
\end{lstlisting}

\end{formaltheorem}

\begin{formaltheorem}[EasyCrypt line 884]
\noindent\textbf{EasyCrypt name.}
\begin{lstlisting}[language=EasyCrypt,basicstyle=\ttfamily\scriptsize]
dexact_hyperball_signing_attempt_state_pack_abort_zero
\end{lstlisting}
\noindent\textbf{Checked EasyCrypt statement.}
\begin{lstlisting}[language=EasyCrypt,basicstyle=\ttfamily\scriptsize]
lemma dexact_hyperball_signing_attempt_state_pack_abort_zero
  md sk m ctx :
  mu (dexact_hyperball_signing_attempt_state md sk m ctx)
    (signing_attempt_state_pack_aborts md) =
  0%r.
\end{lstlisting}
\noindent\textbf{Formal mathematical reading.}
Mathematical theorem. This is an equality or definitional expansion used for rewriting and normalization.

\noindent\textbf{Role in the proof.}
Use in this chapter. It proves an exact game replacement or simulator equivalence.

\noindent\textbf{Proof-script interpretation.}
Proof-script reading. The machine proof decomposes the theorem into rewrites, program-logic obligations, relational equivalences, and arithmetic side conditions.
\emph{Main tactics visible in the proof script:} \ec{have}, \ec{apply}, \ec{rewrite}, \ec{smt}.
\emph{Named identifiers syntactically cited in the proof block:}
\begin{lstlisting}[language=EasyCrypt,basicstyle=\ttfamily\scriptsize]
coins
ctx
dexact_hyperball_signing_attempt_state
dexact_hyperball_signing_attempt_state_supportP
le0
md
mu
mu0
mu_bounded
mu_le
pred0
sample
signing_attempt_state_of_sample_pair_pack_no_abort
signing_attempt_state_pack_aborts
sk
st
stE
st_supp
\end{lstlisting}

\end{formaltheorem}

\begin{formaltheorem}[EasyCrypt line 904]
\noindent\textbf{EasyCrypt name.}
\begin{lstlisting}[language=EasyCrypt,basicstyle=\ttfamily\scriptsize]
dstructural_signing_attempt_state_from_samplerE
\end{lstlisting}
\noindent\textbf{Checked EasyCrypt statement.}
\begin{lstlisting}[language=EasyCrypt,basicstyle=\ttfamily\scriptsize]
lemma dstructural_signing_attempt_state_from_samplerE md sk m ctx :
  dstructural_signing_attempt_state_from_sampler md sk m ctx =
  dstructural_signing_attempt_state md sk m ctx.
\end{lstlisting}
\noindent\textbf{Formal mathematical reading.}
Mathematical theorem. This is an equality or definitional expansion used for rewriting and normalization.

\noindent\textbf{Role in the proof.}
Use in this chapter. It contributes directly to an advantage bound, bad-event bound, or real-arithmetic composition step.

\noindent\textbf{Proof-script interpretation.}
Proof-script reading. The machine proof decomposes the theorem into rewrites, program-logic obligations, relational equivalences, and arithmetic side conditions.
\emph{Main tactics visible in the proof script:} \ec{rewrite}, \ec{have}, \ec{apply}.
\emph{Named identifiers syntactically cited in the proof block:}
\begin{lstlisting}[language=EasyCrypt,basicstyle=\ttfamily\scriptsize]
coins
ctx
dlet_dunit
dmap
dmap_dunit
dsigning_attempt_state_from_sample_pair_sampler
dstructural_signing_attempt_state
dstructural_signing_attempt_state_from_sampler
dstructural_signing_sample_pair
dunit
fun_ext
md
signing_attempt_state_of_sample_pair
signing_sample_pair_of_coins
sk
\end{lstlisting}

\end{formaltheorem}

\begin{formaltheorem}[EasyCrypt line 925]
\noindent\textbf{EasyCrypt name.}
\begin{lstlisting}[language=EasyCrypt,basicstyle=\ttfamily\scriptsize]
dstructural_signing_attempt_stateE
\end{lstlisting}
\noindent\textbf{Checked EasyCrypt statement.}
\begin{lstlisting}[language=EasyCrypt,basicstyle=\ttfamily\scriptsize]
lemma dstructural_signing_attempt_stateE md sk m ctx :
  dstructural_signing_attempt_state md sk m ctx =
  dsigning_attempt_state md sk m ctx.
\end{lstlisting}
\noindent\textbf{Formal mathematical reading.}
Mathematical theorem. This is an equality or definitional expansion used for rewriting and normalization.

\noindent\textbf{Role in the proof.}
Use in this chapter. It preserves an invariant across an oracle call, adversary call, or game procedure.

\noindent\textbf{Proof-script interpretation.}
Proof-script reading. The machine proof decomposes the theorem into rewrites, program-logic obligations, relational equivalences, and arithmetic side conditions.
\emph{Main tactics visible in the proof script:} \ec{rewrite}, \ec{apply}.
\emph{Named identifiers syntactically cited in the proof block:}
\begin{lstlisting}[language=EasyCrypt,basicstyle=\ttfamily\scriptsize]
coins
dsigning_attempt_state
dstructural_signing_attempt_state
eq_dmap
signing_attempt_state_of_sample_pair_currentE
\end{lstlisting}

\end{formaltheorem}

\begin{formaltheorem}[EasyCrypt line 934]
\noindent\textbf{EasyCrypt name.}
\begin{lstlisting}[language=EasyCrypt,basicstyle=\ttfamily\scriptsize]
dstructural_signing_attempt_state_from_sampler_coin_sample_pairE
\end{lstlisting}
\noindent\textbf{Checked EasyCrypt statement.}
\begin{lstlisting}[language=EasyCrypt,basicstyle=\ttfamily\scriptsize]
lemma dstructural_signing_attempt_state_from_sampler_coin_sample_pairE
   md sk m ctx :
  dmap (dstructural_signing_attempt_coin_sample_pair md sk m ctx)
    (signing_attempt_state_of_coin_sample_pair md sk m ctx) =
  dstructural_signing_attempt_state_from_sampler md sk m ctx.
\end{lstlisting}
\noindent\textbf{Formal mathematical reading.}
Mathematical theorem. This is an equality or definitional expansion used for rewriting and normalization.

\noindent\textbf{Role in the proof.}
Use in this chapter. It contributes directly to an advantage bound, bad-event bound, or real-arithmetic composition step.

\noindent\textbf{Proof-script interpretation.}
Proof-script reading. The machine proof decomposes the theorem into rewrites, program-logic obligations, relational equivalences, and arithmetic side conditions.
\emph{Main tactics visible in the proof script:} \ec{rewrite}, \ec{apply}.
\emph{Named identifiers syntactically cited in the proof block:}
\begin{lstlisting}[language=EasyCrypt,basicstyle=\ttfamily\scriptsize]
dmap_comp
dmap_dlet
dsigning_attempt_state_from_sample_pair_sampler
dstructural_signing_attempt_coin_sample_pair
dstructural_signing_attempt_state_from_sampler
eq_dlet
eq_dmap
sample
signing_attempt_state_of_coin_sample_pair
\end{lstlisting}

\end{formaltheorem}

\begin{formaltheorem}[EasyCrypt line 950]
\noindent\textbf{EasyCrypt name.}
\begin{lstlisting}[language=EasyCrypt,basicstyle=\ttfamily\scriptsize]
dexact_hyperball_signing_attempt_state_coin_sample_pairE
\end{lstlisting}
\noindent\textbf{Checked EasyCrypt statement.}
\begin{lstlisting}[language=EasyCrypt,basicstyle=\ttfamily\scriptsize]
lemma dexact_hyperball_signing_attempt_state_coin_sample_pairE
   md sk m ctx :
  dmap (dexact_hyperball_signing_attempt_coin_sample_pair md sk m ctx)
    (signing_attempt_state_of_coin_sample_pair md sk m ctx) =
  dexact_hyperball_signing_attempt_state md sk m ctx.
\end{lstlisting}
\noindent\textbf{Formal mathematical reading.}
Mathematical theorem. This is an equality or definitional expansion used for rewriting and normalization.

\noindent\textbf{Role in the proof.}
Use in this chapter. It contributes directly to an advantage bound, bad-event bound, or real-arithmetic composition step.

\noindent\textbf{Proof-script interpretation.}
Proof-script reading. The machine proof decomposes the theorem into rewrites, program-logic obligations, relational equivalences, and arithmetic side conditions.
\emph{Main tactics visible in the proof script:} \ec{rewrite}, \ec{apply}.
\emph{Named identifiers syntactically cited in the proof block:}
\begin{lstlisting}[language=EasyCrypt,basicstyle=\ttfamily\scriptsize]
dexact_hyperball_signing_attempt_coin_sample_pair
dexact_hyperball_signing_attempt_state
dmap_comp
dmap_dlet
dsigning_attempt_state_from_sample_pair_sampler
eq_dlet
eq_dmap
sample
signing_attempt_state_of_coin_sample_pair
\end{lstlisting}

\end{formaltheorem}

\begin{formaltheorem}[EasyCrypt line 966]
\noindent\textbf{EasyCrypt name.}
\begin{lstlisting}[language=EasyCrypt,basicstyle=\ttfamily\scriptsize]
mu_dlet_le_add_all
\end{lstlisting}
\noindent\textbf{Checked EasyCrypt statement.}
\begin{lstlisting}[language=EasyCrypt,basicstyle=\ttfamily\scriptsize]
lemma mu_dlet_le_add_all ['a 'b 'c]
   (d : 'a distr) (F1 : 'a -> 'b distr) (F2 : 'a -> 'c distr)
   (P1 : 'b -> bool) (P2 : 'c -> bool) (eps : real) :
  0%r <= eps =>
  (forall x, mu (F1 x) P1 <= mu (F2 x) P2 + eps) =>
  mu (dlet d F1) P1 <= mu (dlet d F2) P2 + eps.
\end{lstlisting}
\noindent\textbf{Formal mathematical reading.}
Mathematical theorem. This is an inequality, usually an arithmetic loss bound, event inclusion bound, or monotonicity fact used to compose probabilities.

\noindent\textbf{Role in the proof.}
Use in this chapter. It contributes directly to an advantage bound, bad-event bound, or real-arithmetic composition step.

\noindent\textbf{Proof-script interpretation.}
Proof-script reading. The machine proof decomposes the theorem into rewrites, program-logic obligations, relational equivalences, and arithmetic side conditions.
\emph{Main tactics visible in the proof script:} \ec{rewrite}, \ec{have}, \ec{apply}, \ec{smt}.
\emph{Named identifiers syntactically cited in the proof block:}
\begin{lstlisting}[language=EasyCrypt,basicstyle=\ttfamily\scriptsize]
F1
F2
P1
P2
RealSeries
dletE
eps
eps_ge
eq_sum
exact
fun_ext
ge0_mu1
hsum
ler_add2l
ler_sum
ler_trans
ler_wpmul2l
mu
mu1
mu_bounded
s1_sbl
s2_sbl
s2eps_sbl
seps_sbl
sum
sumD
summable
summableD
summableZ
summable_mu1
summable_mu1_wght
\end{lstlisting}

\end{formaltheorem}

\begin{formaltheorem}[EasyCrypt line 1004]
\noindent\textbf{EasyCrypt name.}
\begin{lstlisting}[language=EasyCrypt,basicstyle=\ttfamily\scriptsize]
structural_to_exact_hyperball_sample_pair_loss_from_point_loss
\end{lstlisting}
\noindent\textbf{Checked EasyCrypt statement.}
\begin{lstlisting}[language=EasyCrypt,basicstyle=\ttfamily\scriptsize]
lemma structural_to_exact_hyperball_sample_pair_loss_from_point_loss :
  structural_to_exact_hyperball_sample_pair_point_loss_obligation =>
  structural_to_exact_hyperball_sample_pair_loss_obligation.
\end{lstlisting}
\noindent\textbf{Formal mathematical reading.}
Mathematical theorem. This is an implication, normally turning one invariant, validity predicate, or proof obligation into another.

\noindent\textbf{Role in the proof.}
Use in this chapter. It contributes directly to an advantage bound, bad-event bound, or real-arithmetic composition step.

\noindent\textbf{Proof-script interpretation.}
Proof-script reading. The machine proof decomposes the theorem into rewrites, program-logic obligations, relational equivalences, and arithmetic side conditions.
\emph{Main tactics visible in the proof script:} \ec{rewrite}, \ec{case}, \ec{have}, \ec{apply}, \ec{smt}.
\emph{Named identifiers syntactically cited in the proof block:}
\begin{lstlisting}[language=EasyCrypt,basicstyle=\ttfamily\scriptsize]
coins
ctx
dexact_hyperball_signing_sample_pair
dstructural_signing_sample_pair
dunitE
hop
hpoint
hs
hsub
md
mu
mu_bounded
mu_sub
pred1
rejection_sampling_loss_term_nonnegative
signing_sample_pair_of_coins
sk
structural_to_exact_hyperball_sample_pair_loss_obligation
structural_to_exact_hyperball_sample_pair_point_loss_obligation
\end{lstlisting}

\end{formaltheorem}

\begin{formaltheorem}[EasyCrypt line 1023]
\noindent\textbf{EasyCrypt name.}
\begin{lstlisting}[language=EasyCrypt,basicstyle=\ttfamily\scriptsize]
structural_to_exact_hyperball_coin_sample_pair_loss_from_sample_pair_loss
\end{lstlisting}
\noindent\textbf{Checked EasyCrypt statement.}
\begin{lstlisting}[language=EasyCrypt,basicstyle=\ttfamily\scriptsize]
lemma structural_to_exact_hyperball_coin_sample_pair_loss_from_sample_pair_loss :
  structural_to_exact_hyperball_sample_pair_loss_obligation =>
  structural_to_exact_hyperball_coin_sample_pair_loss_obligation.
\end{lstlisting}
\noindent\textbf{Formal mathematical reading.}
Mathematical theorem. This is an implication, normally turning one invariant, validity predicate, or proof obligation into another.

\noindent\textbf{Role in the proof.}
Use in this chapter. It contributes directly to an advantage bound, bad-event bound, or real-arithmetic composition step.

\noindent\textbf{Proof-script interpretation.}
Proof-script reading. The machine proof decomposes the theorem into rewrites, program-logic obligations, relational equivalences, and arithmetic side conditions.
\emph{Main tactics visible in the proof script:} \ec{rewrite}, \ec{apply}.
\emph{Named identifiers syntactically cited in the proof block:}
\begin{lstlisting}[language=EasyCrypt,basicstyle=\ttfamily\scriptsize]
coins
ctx
dexact_hyperball_signing_attempt_coin_sample_pair
dmapE
dstructural_signing_attempt_coin_sample_pair
exact
hop
md
mu_dlet_le_add_all
rejection_sampling_loss_term_nonnegative
sample
sk
structural_to_exact_hyperball_coin_sample_pair_loss_obligation
structural_to_exact_hyperball_sample_pair_loss_obligation
\end{lstlisting}

\end{formaltheorem}

\begin{formaltheorem}[EasyCrypt line 1039]
\noindent\textbf{EasyCrypt name.}
\begin{lstlisting}[language=EasyCrypt,basicstyle=\ttfamily\scriptsize]
structural_to_exact_hyperball_attempt_loss_from_coin_sample_pair_loss
\end{lstlisting}
\noindent\textbf{Checked EasyCrypt statement.}
\begin{lstlisting}[language=EasyCrypt,basicstyle=\ttfamily\scriptsize]
lemma structural_to_exact_hyperball_attempt_loss_from_coin_sample_pair_loss :
  structural_to_exact_hyperball_coin_sample_pair_loss_obligation =>
  structural_to_exact_hyperball_attempt_loss_obligation.
\end{lstlisting}
\noindent\textbf{Formal mathematical reading.}
Mathematical theorem. This is an implication, normally turning one invariant, validity predicate, or proof obligation into another.

\noindent\textbf{Role in the proof.}
Use in this chapter. It contributes directly to an advantage bound, bad-event bound, or real-arithmetic composition step.

\noindent\textbf{Proof-script interpretation.}
Proof-script reading. The machine proof decomposes the theorem into rewrites, program-logic obligations, relational equivalences, and arithmetic side conditions.
\emph{Main tactics visible in the proof script:} \ec{rewrite}.
\emph{Named identifiers syntactically cited in the proof block:}
\begin{lstlisting}[language=EasyCrypt,basicstyle=\ttfamily\scriptsize]
cs
ctx
dexact_hyperball_signing_attempt_state_coin_sample_pairE
dmapE
dstructural_signing_attempt_stateE
dstructural_signing_attempt_state_from_samplerE
dstructural_signing_attempt_state_from_sampler_coin_sample_pairE
exact
hop
md
signing_attempt_state_of_coin_sample_pair
sk
structural_to_exact_hyperball_attempt_loss_obligation
structural_to_exact_hyperball_coin_sample_pair_loss_obligation
\end{lstlisting}

\end{formaltheorem}

\begin{formaltheorem}[EasyCrypt line 1055]
\noindent\textbf{EasyCrypt name.}
\begin{lstlisting}[language=EasyCrypt,basicstyle=\ttfamily\scriptsize]
signing_attempt_state_of_coins_signatureE
\end{lstlisting}
\noindent\textbf{Checked EasyCrypt statement.}
\begin{lstlisting}[language=EasyCrypt,basicstyle=\ttfamily\scriptsize]
lemma signing_attempt_state_of_coins_signatureE md sk m ctx coins :
  signing_attempt_state_signature
    (signing_attempt_state_of_coins md sk m ctx coins) =
  sign_internal md sk m ctx coins.
\end{lstlisting}
\noindent\textbf{Formal mathematical reading.}
Mathematical theorem. This is an equality or definitional expansion used for rewriting and normalization.

\noindent\textbf{Role in the proof.}
Use in this chapter. It preserves an invariant across an oracle call, adversary call, or game procedure.

\noindent\textbf{Proof-script interpretation.}
No proof block was collected for this item.

\end{formaltheorem}

\begin{formaltheorem}[EasyCrypt line 1062]
\noindent\textbf{EasyCrypt name.}
\begin{lstlisting}[language=EasyCrypt,basicstyle=\ttfamily\scriptsize]
signing_attempt_state_of_coins_sampled_y1E
\end{lstlisting}
\noindent\textbf{Checked EasyCrypt statement.}
\begin{lstlisting}[language=EasyCrypt,basicstyle=\ttfamily\scriptsize]
lemma signing_attempt_state_of_coins_sampled_y1E md sk m ctx coins :
  signing_attempt_state_sampled_y1
    (signing_attempt_state_of_coins md sk m ctx coins) =
  signing_sample_y1 md sk m ctx coins.
\end{lstlisting}
\noindent\textbf{Formal mathematical reading.}
Mathematical theorem. This is an equality or definitional expansion used for rewriting and normalization.

\noindent\textbf{Role in the proof.}
Use in this chapter. It contributes directly to an advantage bound, bad-event bound, or real-arithmetic composition step.

\noindent\textbf{Proof-script interpretation.}
Proof-script reading. The machine proof decomposes the theorem into rewrites, program-logic obligations, relational equivalences, and arithmetic side conditions.
\emph{Main tactics visible in the proof script:} \ec{rewrite}.
\emph{Named identifiers syntactically cited in the proof block:}
\begin{lstlisting}[language=EasyCrypt,basicstyle=\ttfamily\scriptsize]
signing_attempt_sample_y1
signing_attempt_state_of_coins
signing_attempt_state_sample
signing_attempt_state_sampled_y1
\end{lstlisting}

\end{formaltheorem}

\begin{formaltheorem}[EasyCrypt line 1073]
\noindent\textbf{EasyCrypt name.}
\begin{lstlisting}[language=EasyCrypt,basicstyle=\ttfamily\scriptsize]
signing_attempt_state_of_coins_sampled_y2E
\end{lstlisting}
\noindent\textbf{Checked EasyCrypt statement.}
\begin{lstlisting}[language=EasyCrypt,basicstyle=\ttfamily\scriptsize]
lemma signing_attempt_state_of_coins_sampled_y2E md sk m ctx coins :
  signing_attempt_state_sampled_y2
    (signing_attempt_state_of_coins md sk m ctx coins) =
  signing_sample_y2 md sk m ctx coins.
\end{lstlisting}
\noindent\textbf{Formal mathematical reading.}
Mathematical theorem. This is an equality or definitional expansion used for rewriting and normalization.

\noindent\textbf{Role in the proof.}
Use in this chapter. It contributes directly to an advantage bound, bad-event bound, or real-arithmetic composition step.

\noindent\textbf{Proof-script interpretation.}
Proof-script reading. The machine proof decomposes the theorem into rewrites, program-logic obligations, relational equivalences, and arithmetic side conditions.
\emph{Main tactics visible in the proof script:} \ec{rewrite}.
\emph{Named identifiers syntactically cited in the proof block:}
\begin{lstlisting}[language=EasyCrypt,basicstyle=\ttfamily\scriptsize]
signing_attempt_sample_y2
signing_attempt_state_of_coins
signing_attempt_state_sample
signing_attempt_state_sampled_y2
\end{lstlisting}

\end{formaltheorem}

\begin{formaltheorem}[EasyCrypt line 1084]
\noindent\textbf{EasyCrypt name.}
\begin{lstlisting}[language=EasyCrypt,basicstyle=\ttfamily\scriptsize]
signing_attempt_state_of_coins_sampled_yE
\end{lstlisting}
\noindent\textbf{Checked EasyCrypt statement.}
\begin{lstlisting}[language=EasyCrypt,basicstyle=\ttfamily\scriptsize]
lemma signing_attempt_state_of_coins_sampled_yE md sk m ctx coins :
  signing_attempt_state_sampled_y
    (signing_attempt_state_of_coins md sk m ctx coins) =
  commitment_raw md sk m ctx coins.
\end{lstlisting}
\noindent\textbf{Formal mathematical reading.}
Mathematical theorem. This is an equality or definitional expansion used for rewriting and normalization.

\noindent\textbf{Role in the proof.}
Use in this chapter. It contributes directly to an advantage bound, bad-event bound, or real-arithmetic composition step.

\noindent\textbf{Proof-script interpretation.}
Proof-script reading. The machine proof decomposes the theorem into rewrites, program-logic obligations, relational equivalences, and arithmetic side conditions.
\emph{Main tactics visible in the proof script:} \ec{rewrite}.
\emph{Named identifiers syntactically cited in the proof block:}
\begin{lstlisting}[language=EasyCrypt,basicstyle=\ttfamily\scriptsize]
signing_attempt_sample_commitment_raw
signing_attempt_state_of_coins
signing_attempt_state_sample
signing_attempt_state_sampled_y
\end{lstlisting}

\end{formaltheorem}

\begin{formaltheorem}[EasyCrypt line 1095]
\noindent\textbf{EasyCrypt name.}
\begin{lstlisting}[language=EasyCrypt,basicstyle=\ttfamily\scriptsize]
signing_attempt_state_of_coins_highbitsE
\end{lstlisting}
\noindent\textbf{Checked EasyCrypt statement.}
\begin{lstlisting}[language=EasyCrypt,basicstyle=\ttfamily\scriptsize]
lemma signing_attempt_state_of_coins_highbitsE md sk m ctx coins :
  signing_attempt_state_highbits
    (signing_attempt_state_of_coins md sk m ctx coins) =
  commitment_highbits md sk m ctx coins.
\end{lstlisting}
\noindent\textbf{Formal mathematical reading.}
Mathematical theorem. This is an equality or definitional expansion used for rewriting and normalization.

\noindent\textbf{Role in the proof.}
Use in this chapter. It preserves an invariant across an oracle call, adversary call, or game procedure.

\noindent\textbf{Proof-script interpretation.}
No proof block was collected for this item.

\end{formaltheorem}

\begin{formaltheorem}[EasyCrypt line 1102]
\noindent\textbf{EasyCrypt name.}
\begin{lstlisting}[language=EasyCrypt,basicstyle=\ttfamily\scriptsize]
signing_attempt_state_of_coins_lowbitsE
\end{lstlisting}
\noindent\textbf{Checked EasyCrypt statement.}
\begin{lstlisting}[language=EasyCrypt,basicstyle=\ttfamily\scriptsize]
lemma signing_attempt_state_of_coins_lowbitsE md sk m ctx coins :
  signing_attempt_state_lowbits
    (signing_attempt_state_of_coins md sk m ctx coins) =
  commitment_lowbits md sk m ctx coins.
\end{lstlisting}
\noindent\textbf{Formal mathematical reading.}
Mathematical theorem. This is an equality or definitional expansion used for rewriting and normalization.

\noindent\textbf{Role in the proof.}
Use in this chapter. It preserves an invariant across an oracle call, adversary call, or game procedure.

\noindent\textbf{Proof-script interpretation.}
No proof block was collected for this item.

\end{formaltheorem}

\begin{formaltheorem}[EasyCrypt line 1109]
\noindent\textbf{EasyCrypt name.}
\begin{lstlisting}[language=EasyCrypt,basicstyle=\ttfamily\scriptsize]
signing_attempt_state_of_coins_challengeE
\end{lstlisting}
\noindent\textbf{Checked EasyCrypt statement.}
\begin{lstlisting}[language=EasyCrypt,basicstyle=\ttfamily\scriptsize]
lemma signing_attempt_state_of_coins_challengeE md sk m ctx coins :
  signing_attempt_state_challenge
    (signing_attempt_state_of_coins md sk m ctx coins) =
  commitment_challenge md sk m ctx coins.
\end{lstlisting}
\noindent\textbf{Formal mathematical reading.}
Mathematical theorem. This is an equality or definitional expansion used for rewriting and normalization.

\noindent\textbf{Role in the proof.}
Use in this chapter. It contributes directly to an advantage bound, bad-event bound, or real-arithmetic composition step.

\noindent\textbf{Proof-script interpretation.}
No proof block was collected for this item.

\end{formaltheorem}

\begin{formaltheorem}[EasyCrypt line 1116]
\noindent\textbf{EasyCrypt name.}
\begin{lstlisting}[language=EasyCrypt,basicstyle=\ttfamily\scriptsize]
signing_attempt_state_of_coins_responseE
\end{lstlisting}
\noindent\textbf{Checked EasyCrypt statement.}
\begin{lstlisting}[language=EasyCrypt,basicstyle=\ttfamily\scriptsize]
lemma signing_attempt_state_of_coins_responseE md sk m ctx coins :
  signing_attempt_state_response
    (signing_attempt_state_of_coins md sk m ctx coins) =
  response_vector md sk m ctx coins.
\end{lstlisting}
\noindent\textbf{Formal mathematical reading.}
Mathematical theorem. This is an equality or definitional expansion used for rewriting and normalization.

\noindent\textbf{Role in the proof.}
Use in this chapter. It preserves an invariant across an oracle call, adversary call, or game procedure.

\noindent\textbf{Proof-script interpretation.}
No proof block was collected for this item.

\end{formaltheorem}

\begin{formaltheorem}[EasyCrypt line 1123]
\noindent\textbf{EasyCrypt name.}
\begin{lstlisting}[language=EasyCrypt,basicstyle=\ttfamily\scriptsize]
signing_attempt_state_of_coins_response_auxE
\end{lstlisting}
\noindent\textbf{Checked EasyCrypt statement.}
\begin{lstlisting}[language=EasyCrypt,basicstyle=\ttfamily\scriptsize]
lemma signing_attempt_state_of_coins_response_auxE md sk m ctx coins :
  signing_attempt_state_response_aux
    (signing_attempt_state_of_coins md sk m ctx coins) =
  response_aux_vector md sk m ctx coins.
\end{lstlisting}
\noindent\textbf{Formal mathematical reading.}
Mathematical theorem. This is an equality or definitional expansion used for rewriting and normalization.

\noindent\textbf{Role in the proof.}
Use in this chapter. It preserves an invariant across an oracle call, adversary call, or game procedure.

\noindent\textbf{Proof-script interpretation.}
No proof block was collected for this item.

\end{formaltheorem}

\begin{formaltheorem}[EasyCrypt line 1130]
\noindent\textbf{EasyCrypt name.}
\begin{lstlisting}[language=EasyCrypt,basicstyle=\ttfamily\scriptsize]
signing_attempt_state_of_coins_hintE
\end{lstlisting}
\noindent\textbf{Checked EasyCrypt statement.}
\begin{lstlisting}[language=EasyCrypt,basicstyle=\ttfamily\scriptsize]
lemma signing_attempt_state_of_coins_hintE md sk m ctx coins :
  signing_attempt_state_hint
    (signing_attempt_state_of_coins md sk m ctx coins) =
  hint_vector md sk m ctx coins.
\end{lstlisting}
\noindent\textbf{Formal mathematical reading.}
Mathematical theorem. This is an equality or definitional expansion used for rewriting and normalization.

\noindent\textbf{Role in the proof.}
Use in this chapter. It preserves an invariant across an oracle call, adversary call, or game procedure.

\noindent\textbf{Proof-script interpretation.}
No proof block was collected for this item.

\end{formaltheorem}

\begin{formaltheorem}[EasyCrypt line 1137]
\noindent\textbf{EasyCrypt name.}
\begin{lstlisting}[language=EasyCrypt,basicstyle=\ttfamily\scriptsize]
signing_attempt_state_of_coins_lowbits_carrierE
\end{lstlisting}
\noindent\textbf{Checked EasyCrypt statement.}
\begin{lstlisting}[language=EasyCrypt,basicstyle=\ttfamily\scriptsize]
lemma signing_attempt_state_of_coins_lowbits_carrierE md sk m ctx coins :
  signing_attempt_state_lowbits_carrier
    (signing_attempt_state_of_coins md sk m ctx coins) =
  nth poly_zero (commitment_raw md sk m ctx coins) 0.
\end{lstlisting}
\noindent\textbf{Formal mathematical reading.}
Mathematical theorem. This is an equality or definitional expansion used for rewriting and normalization.

\noindent\textbf{Role in the proof.}
Use in this chapter. It preserves an invariant across an oracle call, adversary call, or game procedure.

\noindent\textbf{Proof-script interpretation.}
Proof-script reading. The machine proof decomposes the theorem into rewrites, program-logic obligations, relational equivalences, and arithmetic side conditions.
\emph{Main tactics visible in the proof script:} \ec{rewrite}.
\emph{Named identifiers syntactically cited in the proof block:}
\begin{lstlisting}[language=EasyCrypt,basicstyle=\ttfamily\scriptsize]
signing_attempt_state_lowbits_carrier
signing_attempt_state_of_coins_sampled_yE
\end{lstlisting}

\end{formaltheorem}

\begin{formaltheorem}[EasyCrypt line 1146]
\noindent\textbf{EasyCrypt name.}
\begin{lstlisting}[language=EasyCrypt,basicstyle=\ttfamily\scriptsize]
signing_attempt_state_of_coins_tokenE
\end{lstlisting}
\noindent\textbf{Checked EasyCrypt statement.}
\begin{lstlisting}[language=EasyCrypt,basicstyle=\ttfamily\scriptsize]
lemma signing_attempt_state_of_coins_tokenE md sk m ctx coins :
  signing_attempt_state_token
    (signing_attempt_state_of_coins md sk m ctx coins) =
  signature_token md sk m ctx coins.
\end{lstlisting}
\noindent\textbf{Formal mathematical reading.}
Mathematical theorem. This is an equality or definitional expansion used for rewriting and normalization.

\noindent\textbf{Role in the proof.}
Use in this chapter. It preserves an invariant across an oracle call, adversary call, or game procedure.

\noindent\textbf{Proof-script interpretation.}
Proof-script reading. The machine proof decomposes the theorem into rewrites, program-logic obligations, relational equivalences, and arithmetic side conditions.
\emph{Main tactics visible in the proof script:} \ec{rewrite}.
\emph{Named identifiers syntactically cited in the proof block:}
\begin{lstlisting}[language=EasyCrypt,basicstyle=\ttfamily\scriptsize]
signature_token
signing_attempt_state_of_coins_hintE
signing_attempt_state_of_coins_responseE
signing_attempt_state_of_coins_response_auxE
signing_attempt_state_token
\end{lstlisting}

\end{formaltheorem}

\begin{formaltheorem}[EasyCrypt line 1158]
\noindent\textbf{EasyCrypt name.}
\begin{lstlisting}[language=EasyCrypt,basicstyle=\ttfamily\scriptsize]
signing_attempt_state_of_coins_programmed_lowbitsE
\end{lstlisting}
\noindent\textbf{Checked EasyCrypt statement.}
\begin{lstlisting}[language=EasyCrypt,basicstyle=\ttfamily\scriptsize]
lemma signing_attempt_state_of_coins_programmed_lowbitsE md sk m ctx coins :
  signing_attempt_state_programmed_lowbits
    (signing_attempt_state_of_coins md sk m ctx coins) =
  commitment_lowbits md sk m ctx coins.
\end{lstlisting}
\noindent\textbf{Formal mathematical reading.}
Mathematical theorem. This is an equality or definitional expansion used for rewriting and normalization.

\noindent\textbf{Role in the proof.}
Use in this chapter. It preserves an invariant across an oracle call, adversary call, or game procedure.

\noindent\textbf{Proof-script interpretation.}
Proof-script reading. The machine proof decomposes the theorem into rewrites, program-logic obligations, relational equivalences, and arithmetic side conditions.
\emph{Main tactics visible in the proof script:} \ec{rewrite}.
\emph{Named identifiers syntactically cited in the proof block:}
\begin{lstlisting}[language=EasyCrypt,basicstyle=\ttfamily\scriptsize]
commitment_lowbits
signing_attempt_state_coins
signing_attempt_state_of_coins
signing_attempt_state_of_coins_lowbits_carrierE
signing_attempt_state_programmed_lowbits
\end{lstlisting}

\end{formaltheorem}

\begin{formaltheorem}[EasyCrypt line 1169]
\noindent\textbf{EasyCrypt name.}
\begin{lstlisting}[language=EasyCrypt,basicstyle=\ttfamily\scriptsize]
signing_attempt_state_of_coins_lowbits_consistent
\end{lstlisting}
\noindent\textbf{Checked EasyCrypt statement.}
\begin{lstlisting}[language=EasyCrypt,basicstyle=\ttfamily\scriptsize]
lemma signing_attempt_state_of_coins_lowbits_consistent md sk m ctx coins :
  signing_attempt_state_lowbits_consistent
    (signing_attempt_state_of_coins md sk m ctx coins).
\end{lstlisting}
\noindent\textbf{Formal mathematical reading.}
Mathematical theorem. This lemma is a universally quantified proposition over the variables shown in the EasyCrypt statement.

\noindent\textbf{Role in the proof.}
Use in this chapter. It preserves an invariant across an oracle call, adversary call, or game procedure.

\noindent\textbf{Proof-script interpretation.}
Proof-script reading. The machine proof decomposes the theorem into rewrites, program-logic obligations, relational equivalences, and arithmetic side conditions.
\emph{Main tactics visible in the proof script:} \ec{rewrite}.
\emph{Named identifiers syntactically cited in the proof block:}
\begin{lstlisting}[language=EasyCrypt,basicstyle=\ttfamily\scriptsize]
signing_attempt_state_lowbits_consistent
signing_attempt_state_of_coins_lowbitsE
signing_attempt_state_of_coins_programmed_lowbitsE
\end{lstlisting}

\end{formaltheorem}

\begin{formaltheorem}[EasyCrypt line 1178]
\noindent\textbf{EasyCrypt name.}
\begin{lstlisting}[language=EasyCrypt,basicstyle=\ttfamily\scriptsize]
signing_attempt_state_of_coins_programmed_challengeE
\end{lstlisting}
\noindent\textbf{Checked EasyCrypt statement.}
\begin{lstlisting}[language=EasyCrypt,basicstyle=\ttfamily\scriptsize]
lemma signing_attempt_state_of_coins_programmed_challengeE md sk m ctx coins :
  signing_attempt_state_programmed_challenge md (public_key_of_secret md sk)
    m ctx (signing_attempt_state_of_coins md sk m ctx coins) =
  commitment_challenge md sk m ctx coins.
\end{lstlisting}
\noindent\textbf{Formal mathematical reading.}
Mathematical theorem. This is an equality or definitional expansion used for rewriting and normalization.

\noindent\textbf{Role in the proof.}
Use in this chapter. It contributes directly to an advantage bound, bad-event bound, or real-arithmetic composition step.

\noindent\textbf{Proof-script interpretation.}
Proof-script reading. The machine proof decomposes the theorem into rewrites, program-logic obligations, relational equivalences, and arithmetic side conditions.
\emph{Main tactics visible in the proof script:} \ec{rewrite}.
\emph{Named identifiers syntactically cited in the proof block:}
\begin{lstlisting}[language=EasyCrypt,basicstyle=\ttfamily\scriptsize]
commitment_challenge
signing_attempt_state_of_coins_lowbitsE
signing_attempt_state_of_coins_response_auxE
signing_attempt_state_programmed_challenge
\end{lstlisting}

\end{formaltheorem}

\begin{formaltheorem}[EasyCrypt line 1189]
\noindent\textbf{EasyCrypt name.}
\begin{lstlisting}[language=EasyCrypt,basicstyle=\ttfamily\scriptsize]
signing_attempt_state_of_coins_challenge_consistent
\end{lstlisting}
\noindent\textbf{Checked EasyCrypt statement.}
\begin{lstlisting}[language=EasyCrypt,basicstyle=\ttfamily\scriptsize]
lemma signing_attempt_state_of_coins_challenge_consistent md sk m ctx coins :
  signing_attempt_state_challenge_consistent md (public_key_of_secret md sk)
    m ctx (signing_attempt_state_of_coins md sk m ctx coins).
\end{lstlisting}
\noindent\textbf{Formal mathematical reading.}
Mathematical theorem. This lemma is a universally quantified proposition over the variables shown in the EasyCrypt statement.

\noindent\textbf{Role in the proof.}
Use in this chapter. It contributes directly to an advantage bound, bad-event bound, or real-arithmetic composition step.

\noindent\textbf{Proof-script interpretation.}
Proof-script reading. The machine proof decomposes the theorem into rewrites, program-logic obligations, relational equivalences, and arithmetic side conditions.
\emph{Main tactics visible in the proof script:} \ec{rewrite}.
\emph{Named identifiers syntactically cited in the proof block:}
\begin{lstlisting}[language=EasyCrypt,basicstyle=\ttfamily\scriptsize]
signing_attempt_state_challenge_consistent
signing_attempt_state_of_coins_challengeE
signing_attempt_state_of_coins_programmed_challengeE
\end{lstlisting}

\end{formaltheorem}

\begin{formaltheorem}[EasyCrypt line 1198]
\noindent\textbf{EasyCrypt name.}
\begin{lstlisting}[language=EasyCrypt,basicstyle=\ttfamily\scriptsize]
signing_attempt_state_of_coins_reconstructed_highbitsE
\end{lstlisting}
\noindent\textbf{Checked EasyCrypt statement.}
\begin{lstlisting}[language=EasyCrypt,basicstyle=\ttfamily\scriptsize]
lemma signing_attempt_state_of_coins_reconstructed_highbitsE
   md sk m ctx coins :
  signing_attempt_state_reconstructed_highbits md (public_key_of_secret md sk)
    (signing_attempt_state_of_coins md sk m ctx coins) =
  commitment_highbits md sk m ctx coins.
\end{lstlisting}
\noindent\textbf{Formal mathematical reading.}
Mathematical theorem. This is an equality or definitional expansion used for rewriting and normalization.

\noindent\textbf{Role in the proof.}
Use in this chapter. It preserves an invariant across an oracle call, adversary call, or game procedure.

\noindent\textbf{Proof-script interpretation.}
Proof-script reading. The machine proof decomposes the theorem into rewrites, program-logic obligations, relational equivalences, and arithmetic side conditions.
\emph{Main tactics visible in the proof script:} \ec{rewrite}.
\emph{Named identifiers syntactically cited in the proof block:}
\begin{lstlisting}[language=EasyCrypt,basicstyle=\ttfamily\scriptsize]
commitment_highbits
signing_attempt_state_of_coins_challengeE
signing_attempt_state_of_coins_response_auxE
signing_attempt_state_reconstructed_highbits
\end{lstlisting}

\end{formaltheorem}

\begin{formaltheorem}[EasyCrypt line 1210]
\noindent\textbf{EasyCrypt name.}
\begin{lstlisting}[language=EasyCrypt,basicstyle=\ttfamily\scriptsize]
signing_attempt_state_of_coins_public_equation_consistent
\end{lstlisting}
\noindent\textbf{Checked EasyCrypt statement.}
\begin{lstlisting}[language=EasyCrypt,basicstyle=\ttfamily\scriptsize]
lemma signing_attempt_state_of_coins_public_equation_consistent
   md sk m ctx coins :
  signing_attempt_state_public_equation_consistent md
    (public_key_of_secret md sk)
    (signing_attempt_state_of_coins md sk m ctx coins).
\end{lstlisting}
\noindent\textbf{Formal mathematical reading.}
Mathematical theorem. This lemma is a universally quantified proposition over the variables shown in the EasyCrypt statement.

\noindent\textbf{Role in the proof.}
Use in this chapter. It preserves an invariant across an oracle call, adversary call, or game procedure.

\noindent\textbf{Proof-script interpretation.}
Proof-script reading. The machine proof decomposes the theorem into rewrites, program-logic obligations, relational equivalences, and arithmetic side conditions.
\emph{Main tactics visible in the proof script:} \ec{rewrite}.
\emph{Named identifiers syntactically cited in the proof block:}
\begin{lstlisting}[language=EasyCrypt,basicstyle=\ttfamily\scriptsize]
signing_attempt_state_of_coins_highbitsE
signing_attempt_state_of_coins_reconstructed_highbitsE
signing_attempt_state_public_equation_consistent
\end{lstlisting}

\end{formaltheorem}

\begin{formaltheorem}[EasyCrypt line 1221]
\noindent\textbf{EasyCrypt name.}
\begin{lstlisting}[language=EasyCrypt,basicstyle=\ttfamily\scriptsize]
signing_attempt_state_of_coins_signature_of_fieldsE
\end{lstlisting}
\noindent\textbf{Checked EasyCrypt statement.}
\begin{lstlisting}[language=EasyCrypt,basicstyle=\ttfamily\scriptsize]
lemma signing_attempt_state_of_coins_signature_of_fieldsE md sk m ctx coins :
  signing_attempt_state_signature_of_fields
    (public_key_of_secret md sk) m ctx
    (signing_attempt_state_of_coins md sk m ctx coins) =
  sign_internal md sk m ctx coins.
\end{lstlisting}
\noindent\textbf{Formal mathematical reading.}
Mathematical theorem. This is an equality or definitional expansion used for rewriting and normalization.

\noindent\textbf{Role in the proof.}
Use in this chapter. It preserves an invariant across an oracle call, adversary call, or game procedure.

\noindent\textbf{Proof-script interpretation.}
Proof-script reading. The machine proof decomposes the theorem into rewrites, program-logic obligations, relational equivalences, and arithmetic side conditions.
\emph{Main tactics visible in the proof script:} \ec{rewrite}.
\emph{Named identifiers syntactically cited in the proof block:}
\begin{lstlisting}[language=EasyCrypt,basicstyle=\ttfamily\scriptsize]
sign_internal
signing_attempt_state_of_coins_challengeE
signing_attempt_state_of_coins_highbitsE
signing_attempt_state_of_coins_lowbitsE
signing_attempt_state_of_coins_tokenE
signing_attempt_state_signature_of_fields
\end{lstlisting}

\end{formaltheorem}

\begin{formaltheorem}[EasyCrypt line 1234]
\noindent\textbf{EasyCrypt name.}
\begin{lstlisting}[language=EasyCrypt,basicstyle=\ttfamily\scriptsize]
signing_attempt_state_of_coins_fields_match_signature
\end{lstlisting}
\noindent\textbf{Checked EasyCrypt statement.}
\begin{lstlisting}[language=EasyCrypt,basicstyle=\ttfamily\scriptsize]
lemma signing_attempt_state_of_coins_fields_match_signature md sk m ctx coins :
  signing_attempt_state_fields_match_signature
    (public_key_of_secret md sk) m ctx
    (signing_attempt_state_of_coins md sk m ctx coins).
\end{lstlisting}
\noindent\textbf{Formal mathematical reading.}
Mathematical theorem. This lemma is a universally quantified proposition over the variables shown in the EasyCrypt statement.

\noindent\textbf{Role in the proof.}
Use in this chapter. It preserves an invariant across an oracle call, adversary call, or game procedure.

\noindent\textbf{Proof-script interpretation.}
Proof-script reading. The machine proof decomposes the theorem into rewrites, program-logic obligations, relational equivalences, and arithmetic side conditions.
\emph{Main tactics visible in the proof script:} \ec{rewrite}.
\emph{Named identifiers syntactically cited in the proof block:}
\begin{lstlisting}[language=EasyCrypt,basicstyle=\ttfamily\scriptsize]
signing_attempt_state_fields_match_signature
signing_attempt_state_of_coins_signatureE
signing_attempt_state_of_coins_signature_of_fieldsE
\end{lstlisting}

\end{formaltheorem}

\begin{formaltheorem}[EasyCrypt line 1244]
\noindent\textbf{EasyCrypt name.}
\begin{lstlisting}[language=EasyCrypt,basicstyle=\ttfamily\scriptsize]
signing_attempt_state_of_coins_field_accepts
\end{lstlisting}
\noindent\textbf{Checked EasyCrypt statement.}
\begin{lstlisting}[language=EasyCrypt,basicstyle=\ttfamily\scriptsize]
lemma signing_attempt_state_of_coins_field_accepts md sk m ctx coins :
  signing_attempt_state_field_accepts md (public_key_of_secret md sk)
    m ctx (signing_attempt_state_of_coins md sk m ctx coins).
\end{lstlisting}
\noindent\textbf{Formal mathematical reading.}
Mathematical theorem. This lemma is a universally quantified proposition over the variables shown in the EasyCrypt statement.

\noindent\textbf{Role in the proof.}
Use in this chapter. It preserves an invariant across an oracle call, adversary call, or game procedure.

\noindent\textbf{Proof-script interpretation.}
Proof-script reading. The machine proof decomposes the theorem into rewrites, program-logic obligations, relational equivalences, and arithmetic side conditions.
\emph{Main tactics visible in the proof script:} \ec{rewrite}, \ec{split}, \ec{apply}.
\emph{Named identifiers syntactically cited in the proof block:}
\begin{lstlisting}[language=EasyCrypt,basicstyle=\ttfamily\scriptsize]
signing_attempt_state_field_accepts
signing_attempt_state_of_coins_challenge_consistent
signing_attempt_state_of_coins_fields_match_signature
signing_attempt_state_of_coins_lowbits_consistent
signing_attempt_state_of_coins_public_equation_consistent
\end{lstlisting}

\end{formaltheorem}

\begin{formaltheorem}[EasyCrypt line 1258]
\noindent\textbf{EasyCrypt name.}
\begin{lstlisting}[language=EasyCrypt,basicstyle=\ttfamily\scriptsize]
signing_attempt_state_of_coins_accepts_current
\end{lstlisting}
\noindent\textbf{Checked EasyCrypt statement.}
\begin{lstlisting}[language=EasyCrypt,basicstyle=\ttfamily\scriptsize]
lemma signing_attempt_state_of_coins_accepts_current md sk m ctx coins :
  signing_attempt_state_accepts md
    (signing_attempt_state_of_coins md sk m ctx coins).
\end{lstlisting}
\noindent\textbf{Formal mathematical reading.}
Mathematical theorem. This lemma is a universally quantified proposition over the variables shown in the EasyCrypt statement.

\noindent\textbf{Role in the proof.}
Use in this chapter. It preserves an invariant across an oracle call, adversary call, or game procedure.

\noindent\textbf{Proof-script interpretation.}
Proof-script reading. The machine proof decomposes the theorem into rewrites, program-logic obligations, relational equivalences, and arithmetic side conditions.
\emph{Main tactics visible in the proof script:} \ec{rewrite}.
\emph{Named identifiers syntactically cited in the proof block:}
\begin{lstlisting}[language=EasyCrypt,basicstyle=\ttfamily\scriptsize]
response_norm_ok_current
signing_attempt_state_accepts
signing_attempt_state_hint_norm_accepts
signing_attempt_state_of_coins_signatureE
signing_attempt_state_response_norm_accepts
signing_attempt_state_valid_signature_accepts
valid_signature_sign_internal
verify_norm_ok_current
\end{lstlisting}

\end{formaltheorem}

\begin{formaltheorem}[EasyCrypt line 1272]
\noindent\textbf{EasyCrypt name.}
\begin{lstlisting}[language=EasyCrypt,basicstyle=\ttfamily\scriptsize]
signing_attempt_reject1_bound_ge_response_bound
\end{lstlisting}
\noindent\textbf{Checked EasyCrypt statement.}
\begin{lstlisting}[language=EasyCrypt,basicstyle=\ttfamily\scriptsize]
lemma signing_attempt_reject1_bound_ge_response_bound md :
  mode_b1sq md <= signing_attempt_reject1_bound md.
\end{lstlisting}
\noindent\textbf{Formal mathematical reading.}
Mathematical theorem. This is an inequality, usually an arithmetic loss bound, event inclusion bound, or monotonicity fact used to compose probabilities.

\noindent\textbf{Role in the proof.}
Use in this chapter. It contributes directly to an advantage bound, bad-event bound, or real-arithmetic composition step.

\noindent\textbf{Proof-script interpretation.}
No proof block was collected for this item.

\end{formaltheorem}

\begin{formaltheorem}[EasyCrypt line 1276]
\noindent\textbf{EasyCrypt name.}
\begin{lstlisting}[language=EasyCrypt,basicstyle=\ttfamily\scriptsize]
signing_attempt_state_reject1_accepts_from_response_norm_ok
\end{lstlisting}
\noindent\textbf{Checked EasyCrypt statement.}
\begin{lstlisting}[language=EasyCrypt,basicstyle=\ttfamily\scriptsize]
lemma signing_attempt_state_reject1_accepts_from_response_norm_ok md st :
  signing_attempt_state_response_norm_accepts md st =>
  signing_attempt_state_reject1_accepts md st.
\end{lstlisting}
\noindent\textbf{Formal mathematical reading.}
Mathematical theorem. This is an implication, normally turning one invariant, validity predicate, or proof obligation into another.

\noindent\textbf{Role in the proof.}
Use in this chapter. It preserves an invariant across an oracle call, adversary call, or game procedure.

\noindent\textbf{Proof-script interpretation.}
Proof-script reading. The machine proof decomposes the theorem into rewrites, program-logic obligations, relational equivalences, and arithmetic side conditions.
\emph{Main tactics visible in the proof script:} \ec{rewrite}, \ec{split}, \ec{smt}.
\emph{Named identifiers syntactically cited in the proof block:}
\begin{lstlisting}[language=EasyCrypt,basicstyle=\ttfamily\scriptsize]
norm_ge0
norm_le
response_norm_ok
response_norm_sq
signing_attempt_reject1_bound_ge_response_bound
signing_attempt_state_reject1_accepts
signing_attempt_state_reject1_norm_sq
signing_attempt_state_response_norm_accepts
\end{lstlisting}

\end{formaltheorem}

\begin{formaltheorem}[EasyCrypt line 1289]
\noindent\textbf{EasyCrypt name.}
\begin{lstlisting}[language=EasyCrypt,basicstyle=\ttfamily\scriptsize]
signing_attempt_state_balance_norm_sq_concreteE
\end{lstlisting}
\noindent\textbf{Checked EasyCrypt statement.}
\begin{lstlisting}[language=EasyCrypt,basicstyle=\ttfamily\scriptsize]
lemma signing_attempt_state_balance_norm_sq_concreteE md st :
  signing_attempt_state_balance_norm_sq md st =
  signing_attempt_reject2_balance_norm_sq md
    (signing_attempt_state_response st)
    (signing_attempt_state_response_aux st)
    (signing_attempt_state_reject2_sampled_y1 st)
    (signing_attempt_state_reject2_sampled_y2 st).
\end{lstlisting}
\noindent\textbf{Formal mathematical reading.}
Mathematical theorem. This is an equality or definitional expansion used for rewriting and normalization.

\noindent\textbf{Role in the proof.}
Use in this chapter. It preserves an invariant across an oracle call, adversary call, or game procedure.

\noindent\textbf{Proof-script interpretation.}
No proof block was collected for this item.

\end{formaltheorem}

\begin{formaltheorem}[EasyCrypt line 1298]
\noindent\textbf{EasyCrypt name.}
\begin{lstlisting}[language=EasyCrypt,basicstyle=\ttfamily\scriptsize]
signing_attempt_state_reject2_under_bound_concreteE
\end{lstlisting}
\noindent\textbf{Checked EasyCrypt statement.}
\begin{lstlisting}[language=EasyCrypt,basicstyle=\ttfamily\scriptsize]
lemma signing_attempt_state_reject2_under_bound_concreteE md st :
  signing_attempt_state_reject2_under_bound md st =
  (signing_attempt_reject2_balance_norm_sq md
    (signing_attempt_state_response st)
    (signing_attempt_state_response_aux st)
    (signing_attempt_state_reject2_sampled_y1 st)
    (signing_attempt_state_reject2_sampled_y2 st) <
   signing_attempt_reject2_bound md).
\end{lstlisting}
\noindent\textbf{Formal mathematical reading.}
Mathematical theorem. This is an equality or definitional expansion used for rewriting and normalization.

\noindent\textbf{Role in the proof.}
Use in this chapter. It contributes directly to an advantage bound, bad-event bound, or real-arithmetic composition step.

\noindent\textbf{Proof-script interpretation.}
Proof-script reading. The machine proof decomposes the theorem into rewrites, program-logic obligations, relational equivalences, and arithmetic side conditions.
\emph{Main tactics visible in the proof script:} \ec{rewrite}.
\emph{Named identifiers syntactically cited in the proof block:}
\begin{lstlisting}[language=EasyCrypt,basicstyle=\ttfamily\scriptsize]
signing_attempt_state_balance_norm_sq_concreteE
signing_attempt_state_reject2_under_bound
\end{lstlisting}

\end{formaltheorem}

\begin{formaltheorem}[EasyCrypt line 1311]
\noindent\textbf{EasyCrypt name.}
\begin{lstlisting}[language=EasyCrypt,basicstyle=\ttfamily\scriptsize]
signing_attempt_state_reject2_aborts_concreteE
\end{lstlisting}
\noindent\textbf{Checked EasyCrypt statement.}
\begin{lstlisting}[language=EasyCrypt,basicstyle=\ttfamily\scriptsize]
lemma signing_attempt_state_reject2_aborts_concreteE md st :
  signing_attempt_state_reject2_aborts md st =
  signing_attempt_reject2_concrete_aborts md
    (signing_attempt_state_bimodal_reject2_branch st)
    (signing_attempt_state_response st)
    (signing_attempt_state_response_aux st)
    (signing_attempt_state_reject2_sampled_y1 st)
    (signing_attempt_state_reject2_sampled_y2 st).
\end{lstlisting}
\noindent\textbf{Formal mathematical reading.}
Mathematical theorem. This is an equality or definitional expansion used for rewriting and normalization.

\noindent\textbf{Role in the proof.}
Use in this chapter. It preserves an invariant across an oracle call, adversary call, or game procedure.

\noindent\textbf{Proof-script interpretation.}
No proof block was collected for this item.

\end{formaltheorem}

\begin{formaltheorem}[EasyCrypt line 1321]
\noindent\textbf{EasyCrypt name.}
\begin{lstlisting}[language=EasyCrypt,basicstyle=\ttfamily\scriptsize]
signing_attempt_state_reject2_accepts_from_branch_clear
\end{lstlisting}
\noindent\textbf{Checked EasyCrypt statement.}
\begin{lstlisting}[language=EasyCrypt,basicstyle=\ttfamily\scriptsize]
lemma signing_attempt_state_reject2_accepts_from_branch_clear md st :
  ! signing_attempt_state_bimodal_reject2_branch st =>
  signing_attempt_state_reject2_accepts md st.
\end{lstlisting}
\noindent\textbf{Formal mathematical reading.}
Mathematical theorem. This is an implication, normally turning one invariant, validity predicate, or proof obligation into another.

\noindent\textbf{Role in the proof.}
Use in this chapter. It contributes directly to an advantage bound, bad-event bound, or real-arithmetic composition step.

\noindent\textbf{Proof-script interpretation.}
Proof-script reading. The machine proof decomposes the theorem into rewrites, program-logic obligations, relational equivalences, and arithmetic side conditions.
\emph{Main tactics visible in the proof script:} \ec{rewrite}.
\emph{Named identifiers syntactically cited in the proof block:}
\begin{lstlisting}[language=EasyCrypt,basicstyle=\ttfamily\scriptsize]
branch_clear
signing_attempt_reject2_concrete_aborts
signing_attempt_state_reject2_aborts
signing_attempt_state_reject2_accepts
\end{lstlisting}

\end{formaltheorem}

\begin{formaltheorem}[EasyCrypt line 1332]
\noindent\textbf{EasyCrypt name.}
\begin{lstlisting}[language=EasyCrypt,basicstyle=\ttfamily\scriptsize]
signing_attempt_state_of_sample_pair_reject2_branch_clear
\end{lstlisting}
\noindent\textbf{Checked EasyCrypt statement.}
\begin{lstlisting}[language=EasyCrypt,basicstyle=\ttfamily\scriptsize]
lemma signing_attempt_state_of_sample_pair_reject2_branch_clear
   md sk m ctx coins sample :
  ! signing_attempt_state_bimodal_reject2_branch
      (signing_attempt_state_of_sample_pair md sk m ctx coins sample).
\end{lstlisting}
\noindent\textbf{Formal mathematical reading.}
Mathematical theorem. This lemma is a universally quantified proposition over the variables shown in the EasyCrypt statement.

\noindent\textbf{Role in the proof.}
Use in this chapter. It contributes directly to an advantage bound, bad-event bound, or real-arithmetic composition step.

\noindent\textbf{Proof-script interpretation.}
Proof-script reading. The machine proof decomposes the theorem into rewrites, program-logic obligations, relational equivalences, and arithmetic side conditions.
\emph{Main tactics visible in the proof script:} \ec{rewrite}.
\emph{Named identifiers syntactically cited in the proof block:}
\begin{lstlisting}[language=EasyCrypt,basicstyle=\ttfamily\scriptsize]
sign_internal
signature_rejection_branch_byte
signing_attempt_state_bimodal_reject2_branch
signing_attempt_state_of_sample_pair_signatureE
signing_attempt_state_reject2_branch_bit
signing_attempt_state_reject2_branch_byte
\end{lstlisting}

\end{formaltheorem}

\begin{formaltheorem}[EasyCrypt line 1345]
\noindent\textbf{EasyCrypt name.}
\begin{lstlisting}[language=EasyCrypt,basicstyle=\ttfamily\scriptsize]
signing_attempt_state_of_sample_pair_reject2_accepts
\end{lstlisting}
\noindent\textbf{Checked EasyCrypt statement.}
\begin{lstlisting}[language=EasyCrypt,basicstyle=\ttfamily\scriptsize]
lemma signing_attempt_state_of_sample_pair_reject2_accepts
   md sk m ctx coins sample :
  signing_attempt_state_reject2_accepts md
    (signing_attempt_state_of_sample_pair md sk m ctx coins sample).
\end{lstlisting}
\noindent\textbf{Formal mathematical reading.}
Mathematical theorem. This lemma is a universally quantified proposition over the variables shown in the EasyCrypt statement.

\noindent\textbf{Role in the proof.}
Use in this chapter. It contributes directly to an advantage bound, bad-event bound, or real-arithmetic composition step.

\noindent\textbf{Proof-script interpretation.}
Proof-script reading. The machine proof decomposes the theorem into rewrites, program-logic obligations, relational equivalences, and arithmetic side conditions.
\emph{Main tactics visible in the proof script:} \ec{apply}.
\emph{Named identifiers syntactically cited in the proof block:}
\begin{lstlisting}[language=EasyCrypt,basicstyle=\ttfamily\scriptsize]
signing_attempt_state_of_sample_pair_reject2_branch_clear
signing_attempt_state_reject2_accepts_from_branch_clear
\end{lstlisting}

\end{formaltheorem}

\begin{formaltheorem}[EasyCrypt line 1354]
\noindent\textbf{EasyCrypt name.}
\begin{lstlisting}[language=EasyCrypt,basicstyle=\ttfamily\scriptsize]
signing_attempt_state_of_sample_pair_reject2_no_abort
\end{lstlisting}
\noindent\textbf{Checked EasyCrypt statement.}
\begin{lstlisting}[language=EasyCrypt,basicstyle=\ttfamily\scriptsize]
lemma signing_attempt_state_of_sample_pair_reject2_no_abort
   md sk m ctx coins sample :
  ! signing_attempt_state_reject2_aborts md
      (signing_attempt_state_of_sample_pair md sk m ctx coins sample).
\end{lstlisting}
\noindent\textbf{Formal mathematical reading.}
Mathematical theorem. This lemma is a universally quantified proposition over the variables shown in the EasyCrypt statement.

\noindent\textbf{Role in the proof.}
Use in this chapter. It contributes directly to an advantage bound, bad-event bound, or real-arithmetic composition step.

\noindent\textbf{Proof-script interpretation.}
Proof-script reading. The machine proof decomposes the theorem into rewrites, program-logic obligations, relational equivalences, and arithmetic side conditions.
\emph{Main tactics visible in the proof script:} \ec{rewrite}.
\emph{Named identifiers syntactically cited in the proof block:}
\begin{lstlisting}[language=EasyCrypt,basicstyle=\ttfamily\scriptsize]
signing_attempt_reject2_concrete_aborts
signing_attempt_state_of_sample_pair_reject2_branch_clear
signing_attempt_state_reject2_aborts
\end{lstlisting}

\end{formaltheorem}

\begin{formaltheorem}[EasyCrypt line 1364]
\noindent\textbf{EasyCrypt name.}
\begin{lstlisting}[language=EasyCrypt,basicstyle=\ttfamily\scriptsize]
dexact_hyperball_signing_attempt_state_reject2_abort_zero
\end{lstlisting}
\noindent\textbf{Checked EasyCrypt statement.}
\begin{lstlisting}[language=EasyCrypt,basicstyle=\ttfamily\scriptsize]
lemma dexact_hyperball_signing_attempt_state_reject2_abort_zero
  md sk m ctx :
  mu (dexact_hyperball_signing_attempt_state md sk m ctx)
    (signing_attempt_state_reject2_aborts md) =
  0%r.
\end{lstlisting}
\noindent\textbf{Formal mathematical reading.}
Mathematical theorem. This is an equality or definitional expansion used for rewriting and normalization.

\noindent\textbf{Role in the proof.}
Use in this chapter. It proves an exact game replacement or simulator equivalence.

\noindent\textbf{Proof-script interpretation.}
Proof-script reading. The machine proof decomposes the theorem into rewrites, program-logic obligations, relational equivalences, and arithmetic side conditions.
\emph{Main tactics visible in the proof script:} \ec{have}, \ec{apply}, \ec{rewrite}, \ec{smt}.
\emph{Named identifiers syntactically cited in the proof block:}
\begin{lstlisting}[language=EasyCrypt,basicstyle=\ttfamily\scriptsize]
coins
ctx
dexact_hyperball_signing_attempt_state
dexact_hyperball_signing_attempt_state_supportP
le0
md
mu
mu0
mu_bounded
mu_le
pred0
sample
signing_attempt_state_of_sample_pair_reject2_no_abort
signing_attempt_state_reject2_aborts
sk
st
stE
st_supp
\end{lstlisting}

\end{formaltheorem}

\begin{formaltheorem}[EasyCrypt line 1384]
\noindent\textbf{EasyCrypt name.}
\begin{lstlisting}[language=EasyCrypt,basicstyle=\ttfamily\scriptsize]
dexact_hyperball_signing_attempt_state_reference_rejection_abort_reduced_bound
\end{lstlisting}
\noindent\textbf{Checked EasyCrypt statement.}
\begin{lstlisting}[language=EasyCrypt,basicstyle=\ttfamily\scriptsize]
lemma dexact_hyperball_signing_attempt_state_reference_rejection_abort_reduced_bound
  md sk m ctx :
  mu (dexact_hyperball_signing_attempt_state md sk m ctx)
    (signing_attempt_state_reference_rejection_aborts md) <=
  mu (dexact_hyperball_signing_attempt_state md sk m ctx)
    (signing_attempt_state_reject1_aborts md) +
  mu (dexact_hyperball_signing_attempt_state md sk m ctx)
    (signing_attempt_state_hint_norm_aborts md).
\end{lstlisting}
\noindent\textbf{Formal mathematical reading.}
Mathematical theorem. This is an inequality, usually an arithmetic loss bound, event inclusion bound, or monotonicity fact used to compose probabilities.

\noindent\textbf{Role in the proof.}
Use in this chapter. It contributes directly to an advantage bound, bad-event bound, or real-arithmetic composition step.

\noindent\textbf{Proof-script interpretation.}
Proof-script reading. The machine proof decomposes the theorem into rewrites, program-logic obligations, relational equivalences, and arithmetic side conditions.
\emph{Main tactics visible in the proof script:} \ec{have}, \ec{rewrite}, \ec{smt}.
\emph{Named identifiers syntactically cited in the proof block:}
\begin{lstlisting}[language=EasyCrypt,basicstyle=\ttfamily\scriptsize]
ctx
dexact_hyperball_signing_attempt_state_pack_abort_zero
dexact_hyperball_signing_attempt_state_reference_rejection_abort_union_bound
dexact_hyperball_signing_attempt_state_reject2_abort_zero
md
mu_bounded
sk
ub
\end{lstlisting}

\end{formaltheorem}

\begin{formaltheorem}[EasyCrypt line 1401]
\noindent\textbf{EasyCrypt name.}
\begin{lstlisting}[language=EasyCrypt,basicstyle=\ttfamily\scriptsize]
signing_attempt_state_of_sample_pair_response_norm_accepts
\end{lstlisting}
\noindent\textbf{Checked EasyCrypt statement.}
\begin{lstlisting}[language=EasyCrypt,basicstyle=\ttfamily\scriptsize]
lemma signing_attempt_state_of_sample_pair_response_norm_accepts
   md sk m ctx coins sample :
  signing_attempt_state_response_norm_accepts md
    (signing_attempt_state_of_sample_pair md sk m ctx coins sample).
\end{lstlisting}
\noindent\textbf{Formal mathematical reading.}
Mathematical theorem. This lemma is a universally quantified proposition over the variables shown in the EasyCrypt statement.

\noindent\textbf{Role in the proof.}
Use in this chapter. It contributes directly to an advantage bound, bad-event bound, or real-arithmetic composition step.

\noindent\textbf{Proof-script interpretation.}
Proof-script reading. The machine proof decomposes the theorem into rewrites, program-logic obligations, relational equivalences, and arithmetic side conditions.
\emph{Main tactics visible in the proof script:} \ec{rewrite}.
\emph{Named identifiers syntactically cited in the proof block:}
\begin{lstlisting}[language=EasyCrypt,basicstyle=\ttfamily\scriptsize]
response_norm_ok_current
signing_attempt_state_of_sample_pair_signatureE
signing_attempt_state_response_norm_accepts
\end{lstlisting}

\end{formaltheorem}

\begin{formaltheorem}[EasyCrypt line 1411]
\noindent\textbf{EasyCrypt name.}
\begin{lstlisting}[language=EasyCrypt,basicstyle=\ttfamily\scriptsize]
signing_attempt_state_of_sample_pair_reject1_accepts
\end{lstlisting}
\noindent\textbf{Checked EasyCrypt statement.}
\begin{lstlisting}[language=EasyCrypt,basicstyle=\ttfamily\scriptsize]
lemma signing_attempt_state_of_sample_pair_reject1_accepts
   md sk m ctx coins sample :
  signing_attempt_state_reject1_accepts md
    (signing_attempt_state_of_sample_pair md sk m ctx coins sample).
\end{lstlisting}
\noindent\textbf{Formal mathematical reading.}
Mathematical theorem. This lemma is a universally quantified proposition over the variables shown in the EasyCrypt statement.

\noindent\textbf{Role in the proof.}
Use in this chapter. It contributes directly to an advantage bound, bad-event bound, or real-arithmetic composition step.

\noindent\textbf{Proof-script interpretation.}
Proof-script reading. The machine proof decomposes the theorem into rewrites, program-logic obligations, relational equivalences, and arithmetic side conditions.
\emph{Main tactics visible in the proof script:} \ec{apply}.
\emph{Named identifiers syntactically cited in the proof block:}
\begin{lstlisting}[language=EasyCrypt,basicstyle=\ttfamily\scriptsize]
signing_attempt_state_of_sample_pair_response_norm_accepts
signing_attempt_state_reject1_accepts_from_response_norm_ok
\end{lstlisting}

\end{formaltheorem}

\begin{formaltheorem}[EasyCrypt line 1420]
\noindent\textbf{EasyCrypt name.}
\begin{lstlisting}[language=EasyCrypt,basicstyle=\ttfamily\scriptsize]
signing_attempt_state_of_sample_pair_reject1_no_abort
\end{lstlisting}
\noindent\textbf{Checked EasyCrypt statement.}
\begin{lstlisting}[language=EasyCrypt,basicstyle=\ttfamily\scriptsize]
lemma signing_attempt_state_of_sample_pair_reject1_no_abort
   md sk m ctx coins sample :
  ! signing_attempt_state_reject1_aborts md
      (signing_attempt_state_of_sample_pair md sk m ctx coins sample).
\end{lstlisting}
\noindent\textbf{Formal mathematical reading.}
Mathematical theorem. This lemma is a universally quantified proposition over the variables shown in the EasyCrypt statement.

\noindent\textbf{Role in the proof.}
Use in this chapter. It contributes directly to an advantage bound, bad-event bound, or real-arithmetic composition step.

\noindent\textbf{Proof-script interpretation.}
Proof-script reading. The machine proof decomposes the theorem into rewrites, program-logic obligations, relational equivalences, and arithmetic side conditions.
\emph{Main tactics visible in the proof script:} \ec{rewrite}.
\emph{Named identifiers syntactically cited in the proof block:}
\begin{lstlisting}[language=EasyCrypt,basicstyle=\ttfamily\scriptsize]
signing_attempt_state_of_sample_pair_reject1_accepts
signing_attempt_state_reject1_aborts
\end{lstlisting}

\end{formaltheorem}

\begin{formaltheorem}[EasyCrypt line 1429]
\noindent\textbf{EasyCrypt name.}
\begin{lstlisting}[language=EasyCrypt,basicstyle=\ttfamily\scriptsize]
dexact_hyperball_signing_attempt_state_reject1_abort_zero
\end{lstlisting}
\noindent\textbf{Checked EasyCrypt statement.}
\begin{lstlisting}[language=EasyCrypt,basicstyle=\ttfamily\scriptsize]
lemma dexact_hyperball_signing_attempt_state_reject1_abort_zero
  md sk m ctx :
  mu (dexact_hyperball_signing_attempt_state md sk m ctx)
    (signing_attempt_state_reject1_aborts md) =
  0%r.
\end{lstlisting}
\noindent\textbf{Formal mathematical reading.}
Mathematical theorem. This is an equality or definitional expansion used for rewriting and normalization.

\noindent\textbf{Role in the proof.}
Use in this chapter. It proves an exact game replacement or simulator equivalence.

\noindent\textbf{Proof-script interpretation.}
Proof-script reading. The machine proof decomposes the theorem into rewrites, program-logic obligations, relational equivalences, and arithmetic side conditions.
\emph{Main tactics visible in the proof script:} \ec{have}, \ec{apply}, \ec{rewrite}, \ec{smt}.
\emph{Named identifiers syntactically cited in the proof block:}
\begin{lstlisting}[language=EasyCrypt,basicstyle=\ttfamily\scriptsize]
coins
ctx
dexact_hyperball_signing_attempt_state
dexact_hyperball_signing_attempt_state_supportP
le0
md
mu
mu0
mu_bounded
mu_le
pred0
sample
signing_attempt_state_of_sample_pair_reject1_no_abort
signing_attempt_state_reject1_aborts
sk
st
stE
st_supp
\end{lstlisting}

\end{formaltheorem}

\begin{formaltheorem}[EasyCrypt line 1449]
\noindent\textbf{EasyCrypt name.}
\begin{lstlisting}[language=EasyCrypt,basicstyle=\ttfamily\scriptsize]
dexact_hyperball_signing_attempt_state_reference_rejection_abort_hint_bound
\end{lstlisting}
\noindent\textbf{Checked EasyCrypt statement.}
\begin{lstlisting}[language=EasyCrypt,basicstyle=\ttfamily\scriptsize]
lemma dexact_hyperball_signing_attempt_state_reference_rejection_abort_hint_bound
  md sk m ctx :
  mu (dexact_hyperball_signing_attempt_state md sk m ctx)
    (signing_attempt_state_reference_rejection_aborts md) <=
  mu (dexact_hyperball_signing_attempt_state md sk m ctx)
    (signing_attempt_state_hint_norm_aborts md).
\end{lstlisting}
\noindent\textbf{Formal mathematical reading.}
Mathematical theorem. This is an inequality, usually an arithmetic loss bound, event inclusion bound, or monotonicity fact used to compose probabilities.

\noindent\textbf{Role in the proof.}
Use in this chapter. It contributes directly to an advantage bound, bad-event bound, or real-arithmetic composition step.

\noindent\textbf{Proof-script interpretation.}
Proof-script reading. The machine proof decomposes the theorem into rewrites, program-logic obligations, relational equivalences, and arithmetic side conditions.
\emph{Main tactics visible in the proof script:} \ec{have}, \ec{rewrite}, \ec{smt}.
\emph{Named identifiers syntactically cited in the proof block:}
\begin{lstlisting}[language=EasyCrypt,basicstyle=\ttfamily\scriptsize]
ctx
dexact_hyperball_signing_attempt_state_reference_rejection_abort_reduced_bound
dexact_hyperball_signing_attempt_state_reject1_abort_zero
md
mu_bounded
sk
ub
\end{lstlisting}

\end{formaltheorem}

\begin{formaltheorem}[EasyCrypt line 1463]
\noindent\textbf{EasyCrypt name.}
\begin{lstlisting}[language=EasyCrypt,basicstyle=\ttfamily\scriptsize]
signing_attempt_state_of_sample_pair_hint_norm_accepts
\end{lstlisting}
\noindent\textbf{Checked EasyCrypt statement.}
\begin{lstlisting}[language=EasyCrypt,basicstyle=\ttfamily\scriptsize]
lemma signing_attempt_state_of_sample_pair_hint_norm_accepts
   md sk m ctx coins sample :
  signing_attempt_state_hint_norm_accepts md
    (signing_attempt_state_of_sample_pair md sk m ctx coins sample).
\end{lstlisting}
\noindent\textbf{Formal mathematical reading.}
Mathematical theorem. This lemma is a universally quantified proposition over the variables shown in the EasyCrypt statement.

\noindent\textbf{Role in the proof.}
Use in this chapter. It contributes directly to an advantage bound, bad-event bound, or real-arithmetic composition step.

\noindent\textbf{Proof-script interpretation.}
Proof-script reading. The machine proof decomposes the theorem into rewrites, program-logic obligations, relational equivalences, and arithmetic side conditions.
\emph{Main tactics visible in the proof script:} \ec{rewrite}.
\emph{Named identifiers syntactically cited in the proof block:}
\begin{lstlisting}[language=EasyCrypt,basicstyle=\ttfamily\scriptsize]
signing_attempt_state_hint_norm_accepts
signing_attempt_state_of_sample_pair_signatureE
verify_norm_ok_current
\end{lstlisting}

\end{formaltheorem}

\begin{formaltheorem}[EasyCrypt line 1473]
\noindent\textbf{EasyCrypt name.}
\begin{lstlisting}[language=EasyCrypt,basicstyle=\ttfamily\scriptsize]
signing_attempt_state_of_sample_pair_hint_norm_no_abort
\end{lstlisting}
\noindent\textbf{Checked EasyCrypt statement.}
\begin{lstlisting}[language=EasyCrypt,basicstyle=\ttfamily\scriptsize]
lemma signing_attempt_state_of_sample_pair_hint_norm_no_abort
   md sk m ctx coins sample :
  ! signing_attempt_state_hint_norm_aborts md
      (signing_attempt_state_of_sample_pair md sk m ctx coins sample).
\end{lstlisting}
\noindent\textbf{Formal mathematical reading.}
Mathematical theorem. This lemma is a universally quantified proposition over the variables shown in the EasyCrypt statement.

\noindent\textbf{Role in the proof.}
Use in this chapter. It contributes directly to an advantage bound, bad-event bound, or real-arithmetic composition step.

\noindent\textbf{Proof-script interpretation.}
Proof-script reading. The machine proof decomposes the theorem into rewrites, program-logic obligations, relational equivalences, and arithmetic side conditions.
\emph{Main tactics visible in the proof script:} \ec{rewrite}.
\emph{Named identifiers syntactically cited in the proof block:}
\begin{lstlisting}[language=EasyCrypt,basicstyle=\ttfamily\scriptsize]
signing_attempt_state_hint_norm_aborts
signing_attempt_state_of_sample_pair_hint_norm_accepts
\end{lstlisting}

\end{formaltheorem}

\begin{formaltheorem}[EasyCrypt line 1482]
\noindent\textbf{EasyCrypt name.}
\begin{lstlisting}[language=EasyCrypt,basicstyle=\ttfamily\scriptsize]
dexact_hyperball_signing_attempt_state_hint_norm_abort_zero
\end{lstlisting}
\noindent\textbf{Checked EasyCrypt statement.}
\begin{lstlisting}[language=EasyCrypt,basicstyle=\ttfamily\scriptsize]
lemma dexact_hyperball_signing_attempt_state_hint_norm_abort_zero
  md sk m ctx :
  mu (dexact_hyperball_signing_attempt_state md sk m ctx)
    (signing_attempt_state_hint_norm_aborts md) =
  0%r.
\end{lstlisting}
\noindent\textbf{Formal mathematical reading.}
Mathematical theorem. This is an equality or definitional expansion used for rewriting and normalization.

\noindent\textbf{Role in the proof.}
Use in this chapter. It proves an exact game replacement or simulator equivalence.

\noindent\textbf{Proof-script interpretation.}
Proof-script reading. The machine proof decomposes the theorem into rewrites, program-logic obligations, relational equivalences, and arithmetic side conditions.
\emph{Main tactics visible in the proof script:} \ec{have}, \ec{apply}, \ec{rewrite}, \ec{smt}.
\emph{Named identifiers syntactically cited in the proof block:}
\begin{lstlisting}[language=EasyCrypt,basicstyle=\ttfamily\scriptsize]
coins
ctx
dexact_hyperball_signing_attempt_state
dexact_hyperball_signing_attempt_state_supportP
le0
md
mu
mu0
mu_bounded
mu_le
pred0
sample
signing_attempt_state_hint_norm_aborts
signing_attempt_state_of_sample_pair_hint_norm_no_abort
sk
st
stE
st_supp
\end{lstlisting}

\end{formaltheorem}

\begin{formaltheorem}[EasyCrypt line 1502]
\noindent\textbf{EasyCrypt name.}
\begin{lstlisting}[language=EasyCrypt,basicstyle=\ttfamily\scriptsize]
dexact_hyperball_signing_attempt_state_reference_rejection_abort_zero
\end{lstlisting}
\noindent\textbf{Checked EasyCrypt statement.}
\begin{lstlisting}[language=EasyCrypt,basicstyle=\ttfamily\scriptsize]
lemma dexact_hyperball_signing_attempt_state_reference_rejection_abort_zero
  md sk m ctx :
  mu (dexact_hyperball_signing_attempt_state md sk m ctx)
    (signing_attempt_state_reference_rejection_aborts md) =
  0%r.
\end{lstlisting}
\noindent\textbf{Formal mathematical reading.}
Mathematical theorem. This is an equality or definitional expansion used for rewriting and normalization.

\noindent\textbf{Role in the proof.}
Use in this chapter. It proves an exact game replacement or simulator equivalence.

\noindent\textbf{Proof-script interpretation.}
Proof-script reading. The machine proof decomposes the theorem into rewrites, program-logic obligations, relational equivalences, and arithmetic side conditions.
\emph{Main tactics visible in the proof script:} \ec{have}, \ec{rewrite}, \ec{smt}.
\emph{Named identifiers syntactically cited in the proof block:}
\begin{lstlisting}[language=EasyCrypt,basicstyle=\ttfamily\scriptsize]
ctx
dexact_hyperball_signing_attempt_state_hint_norm_abort_zero
dexact_hyperball_signing_attempt_state_reference_rejection_abort_hint_bound
md
mu_bounded
sk
ub
\end{lstlisting}

\end{formaltheorem}

\begin{formaltheorem}[EasyCrypt line 1515]
\noindent\textbf{EasyCrypt name.}
\begin{lstlisting}[language=EasyCrypt,basicstyle=\ttfamily\scriptsize]
signing_attempt_state_of_coins_reject2_branch_clear_current
\end{lstlisting}
\noindent\textbf{Checked EasyCrypt statement.}
\begin{lstlisting}[language=EasyCrypt,basicstyle=\ttfamily\scriptsize]
lemma signing_attempt_state_of_coins_reject2_branch_clear_current
   md sk m ctx coins :
  ! signing_attempt_state_bimodal_reject2_branch
      (signing_attempt_state_of_coins md sk m ctx coins).
\end{lstlisting}
\noindent\textbf{Formal mathematical reading.}
Mathematical theorem. This lemma is a universally quantified proposition over the variables shown in the EasyCrypt statement.

\noindent\textbf{Role in the proof.}
Use in this chapter. It contributes directly to an advantage bound, bad-event bound, or real-arithmetic composition step.

\noindent\textbf{Proof-script interpretation.}
Proof-script reading. The machine proof decomposes the theorem into rewrites, program-logic obligations, relational equivalences, and arithmetic side conditions.
\emph{Main tactics visible in the proof script:} \ec{rewrite}.
\emph{Named identifiers syntactically cited in the proof block:}
\begin{lstlisting}[language=EasyCrypt,basicstyle=\ttfamily\scriptsize]
sign_internal
signature_rejection_branch_byte
signing_attempt_state_bimodal_reject2_branch
signing_attempt_state_of_coins_signatureE
signing_attempt_state_reject2_branch_bit
signing_attempt_state_reject2_branch_byte
\end{lstlisting}

\end{formaltheorem}

\begin{formaltheorem}[EasyCrypt line 1528]
\noindent\textbf{EasyCrypt name.}
\begin{lstlisting}[language=EasyCrypt,basicstyle=\ttfamily\scriptsize]
signing_attempt_state_of_coins_reject2_sampled_y1E
\end{lstlisting}
\noindent\textbf{Checked EasyCrypt statement.}
\begin{lstlisting}[language=EasyCrypt,basicstyle=\ttfamily\scriptsize]
lemma signing_attempt_state_of_coins_reject2_sampled_y1E md sk m ctx coins :
  signing_attempt_state_reject2_sampled_y1
    (signing_attempt_state_of_coins md sk m ctx coins) =
  signing_sample_y1 md sk m ctx coins.
\end{lstlisting}
\noindent\textbf{Formal mathematical reading.}
Mathematical theorem. This is an equality or definitional expansion used for rewriting and normalization.

\noindent\textbf{Role in the proof.}
Use in this chapter. It contributes directly to an advantage bound, bad-event bound, or real-arithmetic composition step.

\noindent\textbf{Proof-script interpretation.}
Proof-script reading. The machine proof decomposes the theorem into rewrites, program-logic obligations, relational equivalences, and arithmetic side conditions.
\emph{Main tactics visible in the proof script:} \ec{rewrite}.
\emph{Named identifiers syntactically cited in the proof block:}
\begin{lstlisting}[language=EasyCrypt,basicstyle=\ttfamily\scriptsize]
signing_attempt_state_of_coins_sampled_y1E
signing_attempt_state_reject2_sampled_y1
\end{lstlisting}

\end{formaltheorem}

\begin{formaltheorem}[EasyCrypt line 1537]
\noindent\textbf{EasyCrypt name.}
\begin{lstlisting}[language=EasyCrypt,basicstyle=\ttfamily\scriptsize]
signing_attempt_state_of_coins_reject2_sampled_y2E
\end{lstlisting}
\noindent\textbf{Checked EasyCrypt statement.}
\begin{lstlisting}[language=EasyCrypt,basicstyle=\ttfamily\scriptsize]
lemma signing_attempt_state_of_coins_reject2_sampled_y2E md sk m ctx coins :
  signing_attempt_state_reject2_sampled_y2
    (signing_attempt_state_of_coins md sk m ctx coins) =
  signing_sample_y2 md sk m ctx coins.
\end{lstlisting}
\noindent\textbf{Formal mathematical reading.}
Mathematical theorem. This is an equality or definitional expansion used for rewriting and normalization.

\noindent\textbf{Role in the proof.}
Use in this chapter. It contributes directly to an advantage bound, bad-event bound, or real-arithmetic composition step.

\noindent\textbf{Proof-script interpretation.}
Proof-script reading. The machine proof decomposes the theorem into rewrites, program-logic obligations, relational equivalences, and arithmetic side conditions.
\emph{Main tactics visible in the proof script:} \ec{rewrite}.
\emph{Named identifiers syntactically cited in the proof block:}
\begin{lstlisting}[language=EasyCrypt,basicstyle=\ttfamily\scriptsize]
signing_attempt_state_of_coins_sampled_y2E
signing_attempt_state_reject2_sampled_y2
\end{lstlisting}

\end{formaltheorem}

\begin{formaltheorem}[EasyCrypt line 1546]
\noindent\textbf{EasyCrypt name.}
\begin{lstlisting}[language=EasyCrypt,basicstyle=\ttfamily\scriptsize]
signing_attempt_state_of_coins_reject2_balance_norm_sqE
\end{lstlisting}
\noindent\textbf{Checked EasyCrypt statement.}
\begin{lstlisting}[language=EasyCrypt,basicstyle=\ttfamily\scriptsize]
lemma signing_attempt_state_of_coins_reject2_balance_norm_sqE
   md sk m ctx coins :
  signing_attempt_state_balance_norm_sq md
    (signing_attempt_state_of_coins md sk m ctx coins) =
  signing_attempt_reject2_balance_norm_sq md
    (response_vector md sk m ctx coins)
    (response_aux_vector md sk m ctx coins)
    (signing_sample_y1 md sk m ctx coins)
    (signing_sample_y2 md sk m ctx coins).
\end{lstlisting}
\noindent\textbf{Formal mathematical reading.}
Mathematical theorem. This is an equality or definitional expansion used for rewriting and normalization.

\noindent\textbf{Role in the proof.}
Use in this chapter. It preserves an invariant across an oracle call, adversary call, or game procedure.

\noindent\textbf{Proof-script interpretation.}
Proof-script reading. The machine proof decomposes the theorem into rewrites, program-logic obligations, relational equivalences, and arithmetic side conditions.
\emph{Main tactics visible in the proof script:} \ec{rewrite}.
\emph{Named identifiers syntactically cited in the proof block:}
\begin{lstlisting}[language=EasyCrypt,basicstyle=\ttfamily\scriptsize]
signing_attempt_state_balance_norm_sq_concreteE
signing_attempt_state_of_coins_reject2_sampled_y1E
signing_attempt_state_of_coins_reject2_sampled_y2E
signing_attempt_state_of_coins_responseE
signing_attempt_state_of_coins_response_auxE
\end{lstlisting}

\end{formaltheorem}

\begin{formaltheorem}[EasyCrypt line 1563]
\noindent\textbf{EasyCrypt name.}
\begin{lstlisting}[language=EasyCrypt,basicstyle=\ttfamily\scriptsize]
signing_attempt_state_of_coins_balance_norm_sq_current
\end{lstlisting}
\noindent\textbf{Checked EasyCrypt statement.}
\begin{lstlisting}[language=EasyCrypt,basicstyle=\ttfamily\scriptsize]
lemma signing_attempt_state_of_coins_balance_norm_sq_current
   md sk m ctx coins :
  signing_attempt_state_balance_norm_sq md
    (signing_attempt_state_of_coins md sk m ctx coins) =
  polyvecl_norm_sq md
    (signing_attempt_state_reject2_balance_response md
      (signing_attempt_state_of_coins md sk m ctx coins)) +
  polyveck_norm_sq md
    (signing_attempt_state_reject2_balance_aux md
      (signing_attempt_state_of_coins md sk m ctx coins)).
\end{lstlisting}
\noindent\textbf{Formal mathematical reading.}
Mathematical theorem. This is an equality or definitional expansion used for rewriting and normalization.

\noindent\textbf{Role in the proof.}
Use in this chapter. It preserves an invariant across an oracle call, adversary call, or game procedure.

\noindent\textbf{Proof-script interpretation.}
Proof-script reading. The machine proof decomposes the theorem into rewrites, program-logic obligations, relational equivalences, and arithmetic side conditions.
\emph{Main tactics visible in the proof script:} \ec{rewrite}.
\emph{Named identifiers syntactically cited in the proof block:}
\begin{lstlisting}[language=EasyCrypt,basicstyle=\ttfamily\scriptsize]
signing_attempt_reject2_balance_norm_sq
signing_attempt_state_balance_norm_sq
signing_attempt_state_reject2_balance_aux
signing_attempt_state_reject2_balance_response
\end{lstlisting}

\end{formaltheorem}

\begin{formaltheorem}[EasyCrypt line 1580]
\noindent\textbf{EasyCrypt name.}
\begin{lstlisting}[language=EasyCrypt,basicstyle=\ttfamily\scriptsize]
signing_attempt_state_of_coins_reject2_accepts_current
\end{lstlisting}
\noindent\textbf{Checked EasyCrypt statement.}
\begin{lstlisting}[language=EasyCrypt,basicstyle=\ttfamily\scriptsize]
lemma signing_attempt_state_of_coins_reject2_accepts_current
   md sk m ctx coins :
  signing_attempt_state_reject2_accepts md
    (signing_attempt_state_of_coins md sk m ctx coins).
\end{lstlisting}
\noindent\textbf{Formal mathematical reading.}
Mathematical theorem. This lemma is a universally quantified proposition over the variables shown in the EasyCrypt statement.

\noindent\textbf{Role in the proof.}
Use in this chapter. It preserves an invariant across an oracle call, adversary call, or game procedure.

\noindent\textbf{Proof-script interpretation.}
Proof-script reading. The machine proof decomposes the theorem into rewrites, program-logic obligations, relational equivalences, and arithmetic side conditions.
\emph{Main tactics visible in the proof script:} \ec{apply}.
\emph{Named identifiers syntactically cited in the proof block:}
\begin{lstlisting}[language=EasyCrypt,basicstyle=\ttfamily\scriptsize]
signing_attempt_state_of_coins_reject2_branch_clear_current
signing_attempt_state_reject2_accepts_from_branch_clear
\end{lstlisting}

\end{formaltheorem}

\begin{formaltheorem}[EasyCrypt line 1589]
\noindent\textbf{EasyCrypt name.}
\begin{lstlisting}[language=EasyCrypt,basicstyle=\ttfamily\scriptsize]
signing_attempt_state_pack_accepts_from_valid
\end{lstlisting}
\noindent\textbf{Checked EasyCrypt statement.}
\begin{lstlisting}[language=EasyCrypt,basicstyle=\ttfamily\scriptsize]
lemma signing_attempt_state_pack_accepts_from_valid md st :
  signing_attempt_state_valid_signature_accepts md st =>
  signing_attempt_state_pack_accepts md st.
\end{lstlisting}
\noindent\textbf{Formal mathematical reading.}
Mathematical theorem. This is an implication, normally turning one invariant, validity predicate, or proof obligation into another.

\noindent\textbf{Role in the proof.}
Use in this chapter. It preserves an invariant across an oracle call, adversary call, or game procedure.

\noindent\textbf{Proof-script interpretation.}
Proof-script reading. The machine proof decomposes the theorem into rewrites, program-logic obligations, relational equivalences, and arithmetic side conditions.
\emph{Main tactics visible in the proof script:} \ec{rewrite}, \ec{split}, \ec{apply}.
\emph{Named identifiers syntactically cited in the proof block:}
\begin{lstlisting}[language=EasyCrypt,basicstyle=\ttfamily\scriptsize]
mode_crypto_bytes_gt0
signing_attempt_state_pack_accepts
valid_ok
\end{lstlisting}

\end{formaltheorem}

\begin{formaltheorem}[EasyCrypt line 1599]
\noindent\textbf{EasyCrypt name.}
\begin{lstlisting}[language=EasyCrypt,basicstyle=\ttfamily\scriptsize]
signing_attempt_state_reference_rejection_accepts_from_accepts
\end{lstlisting}
\noindent\textbf{Checked EasyCrypt statement.}
\begin{lstlisting}[language=EasyCrypt,basicstyle=\ttfamily\scriptsize]
lemma signing_attempt_state_reference_rejection_accepts_from_accepts md st :
  signing_attempt_state_accepts md st =>
  signing_attempt_state_reject2_accepts md st =>
  signing_attempt_state_reference_rejection_accepts md st.
\end{lstlisting}
\noindent\textbf{Formal mathematical reading.}
Mathematical theorem. This is an implication, normally turning one invariant, validity predicate, or proof obligation into another.

\noindent\textbf{Role in the proof.}
Use in this chapter. It preserves an invariant across an oracle call, adversary call, or game procedure.

\noindent\textbf{Proof-script interpretation.}
Proof-script reading. The machine proof decomposes the theorem into rewrites, program-logic obligations, relational equivalences, and arithmetic side conditions.
\emph{Main tactics visible in the proof script:} \ec{rewrite}, \ec{split}, \ec{apply}.
\emph{Named identifiers syntactically cited in the proof block:}
\begin{lstlisting}[language=EasyCrypt,basicstyle=\ttfamily\scriptsize]
hint_ok
reject2_ok
response_ok
signing_attempt_state_accepts
signing_attempt_state_pack_accepts_from_valid
signing_attempt_state_reference_rejection_accepts
signing_attempt_state_reject1_accepts_from_response_norm_ok
valid_ok
\end{lstlisting}

\end{formaltheorem}

\begin{formaltheorem}[EasyCrypt line 1616]
\noindent\textbf{EasyCrypt name.}
\begin{lstlisting}[language=EasyCrypt,basicstyle=\ttfamily\scriptsize]
signing_attempt_state_of_coins_reference_rejection_accepts_current
\end{lstlisting}
\noindent\textbf{Checked EasyCrypt statement.}
\begin{lstlisting}[language=EasyCrypt,basicstyle=\ttfamily\scriptsize]
lemma signing_attempt_state_of_coins_reference_rejection_accepts_current
   md sk m ctx coins :
  signing_attempt_state_reference_rejection_accepts md
    (signing_attempt_state_of_coins md sk m ctx coins).
\end{lstlisting}
\noindent\textbf{Formal mathematical reading.}
Mathematical theorem. This lemma is a universally quantified proposition over the variables shown in the EasyCrypt statement.

\noindent\textbf{Role in the proof.}
Use in this chapter. It preserves an invariant across an oracle call, adversary call, or game procedure.

\noindent\textbf{Proof-script interpretation.}
Proof-script reading. The machine proof decomposes the theorem into rewrites, program-logic obligations, relational equivalences, and arithmetic side conditions.
\emph{Main tactics visible in the proof script:} \ec{apply}.
\emph{Named identifiers syntactically cited in the proof block:}
\begin{lstlisting}[language=EasyCrypt,basicstyle=\ttfamily\scriptsize]
signing_attempt_state_of_coins_accepts_current
signing_attempt_state_of_coins_reject2_accepts_current
signing_attempt_state_reference_rejection_accepts_from_accepts
\end{lstlisting}

\end{formaltheorem}

\begin{formaltheorem}[EasyCrypt line 1626]
\noindent\textbf{EasyCrypt name.}
\begin{lstlisting}[language=EasyCrypt,basicstyle=\ttfamily\scriptsize]
signing_attempt_state_of_coins_verification_accepts_current
\end{lstlisting}
\noindent\textbf{Checked EasyCrypt statement.}
\begin{lstlisting}[language=EasyCrypt,basicstyle=\ttfamily\scriptsize]
lemma signing_attempt_state_of_coins_verification_accepts_current
   md sk m ctx coins :
  signing_attempt_state_verification_accepts md (public_key_of_secret md sk)
    m ctx (signing_attempt_state_of_coins md sk m ctx coins).
\end{lstlisting}
\noindent\textbf{Formal mathematical reading.}
Mathematical theorem. This lemma is a universally quantified proposition over the variables shown in the EasyCrypt statement.

\noindent\textbf{Role in the proof.}
Use in this chapter. It preserves an invariant across an oracle call, adversary call, or game procedure.

\noindent\textbf{Proof-script interpretation.}
Proof-script reading. The machine proof decomposes the theorem into rewrites, program-logic obligations, relational equivalences, and arithmetic side conditions.
\emph{Main tactics visible in the proof script:} \ec{rewrite}, \ec{split}, \ec{apply}.
\emph{Named identifiers syntactically cited in the proof block:}
\begin{lstlisting}[language=EasyCrypt,basicstyle=\ttfamily\scriptsize]
signing_attempt_state_of_coins_accepts_current
signing_attempt_state_of_coins_field_accepts
signing_attempt_state_verification_accepts
\end{lstlisting}

\end{formaltheorem}

\begin{formaltheorem}[EasyCrypt line 1637]
\noindent\textbf{EasyCrypt name.}
\begin{lstlisting}[language=EasyCrypt,basicstyle=\ttfamily\scriptsize]
signing_attempt_state_of_coins_rejection_accepts_current
\end{lstlisting}
\noindent\textbf{Checked EasyCrypt statement.}
\begin{lstlisting}[language=EasyCrypt,basicstyle=\ttfamily\scriptsize]
lemma signing_attempt_state_of_coins_rejection_accepts_current
   md sk m ctx coins :
  signing_attempt_state_rejection_accepts md (public_key_of_secret md sk)
    m ctx (signing_attempt_state_of_coins md sk m ctx coins).
\end{lstlisting}
\noindent\textbf{Formal mathematical reading.}
Mathematical theorem. This lemma is a universally quantified proposition over the variables shown in the EasyCrypt statement.

\noindent\textbf{Role in the proof.}
Use in this chapter. It preserves an invariant across an oracle call, adversary call, or game procedure.

\noindent\textbf{Proof-script interpretation.}
Proof-script reading. The machine proof decomposes the theorem into rewrites, program-logic obligations, relational equivalences, and arithmetic side conditions.
\emph{Main tactics visible in the proof script:} \ec{rewrite}, \ec{split}, \ec{apply}.
\emph{Named identifiers syntactically cited in the proof block:}
\begin{lstlisting}[language=EasyCrypt,basicstyle=\ttfamily\scriptsize]
signing_attempt_state_of_coins_reference_rejection_accepts_current
signing_attempt_state_of_coins_verification_accepts_current
signing_attempt_state_rejection_accepts
\end{lstlisting}

\end{formaltheorem}

\begin{formaltheorem}[EasyCrypt line 1648]
\noindent\textbf{EasyCrypt name.}
\begin{lstlisting}[language=EasyCrypt,basicstyle=\ttfamily\scriptsize]
signing_attempt_state_of_coins_rejection_no_abort_current
\end{lstlisting}
\noindent\textbf{Checked EasyCrypt statement.}
\begin{lstlisting}[language=EasyCrypt,basicstyle=\ttfamily\scriptsize]
lemma signing_attempt_state_of_coins_rejection_no_abort_current
   md sk m ctx coins :
  ! signing_attempt_state_rejection_aborts md (public_key_of_secret md sk)
      m ctx (signing_attempt_state_of_coins md sk m ctx coins).
\end{lstlisting}
\noindent\textbf{Formal mathematical reading.}
Mathematical theorem. This lemma is a universally quantified proposition over the variables shown in the EasyCrypt statement.

\noindent\textbf{Role in the proof.}
Use in this chapter. It preserves an invariant across an oracle call, adversary call, or game procedure.

\noindent\textbf{Proof-script interpretation.}
Proof-script reading. The machine proof decomposes the theorem into rewrites, program-logic obligations, relational equivalences, and arithmetic side conditions.
\emph{Main tactics visible in the proof script:} \ec{rewrite}.
\emph{Named identifiers syntactically cited in the proof block:}
\begin{lstlisting}[language=EasyCrypt,basicstyle=\ttfamily\scriptsize]
signing_attempt_state_of_coins_rejection_accepts_current
signing_attempt_state_rejection_aborts
\end{lstlisting}

\end{formaltheorem}

\begin{formaltheorem}[EasyCrypt line 1657]
\noindent\textbf{EasyCrypt name.}
\begin{lstlisting}[language=EasyCrypt,basicstyle=\ttfamily\scriptsize]
signing_attempt_state_of_coins_no_abort_current
\end{lstlisting}
\noindent\textbf{Checked EasyCrypt statement.}
\begin{lstlisting}[language=EasyCrypt,basicstyle=\ttfamily\scriptsize]
lemma signing_attempt_state_of_coins_no_abort_current md sk m ctx coins :
  ! signing_attempt_state_aborts md
      (signing_attempt_state_of_coins md sk m ctx coins).
\end{lstlisting}
\noindent\textbf{Formal mathematical reading.}
Mathematical theorem. This lemma is a universally quantified proposition over the variables shown in the EasyCrypt statement.

\noindent\textbf{Role in the proof.}
Use in this chapter. It preserves an invariant across an oracle call, adversary call, or game procedure.

\noindent\textbf{Proof-script interpretation.}
Proof-script reading. The machine proof decomposes the theorem into rewrites, program-logic obligations, relational equivalences, and arithmetic side conditions.
\emph{Main tactics visible in the proof script:} \ec{rewrite}.
\emph{Named identifiers syntactically cited in the proof block:}
\begin{lstlisting}[language=EasyCrypt,basicstyle=\ttfamily\scriptsize]
response_norm_ok_current
signing_attempt_state_aborts
signing_attempt_state_hint_norm_aborts
signing_attempt_state_hint_norm_accepts
signing_attempt_state_of_coins_signatureE
signing_attempt_state_rejection_sampling_aborts
signing_attempt_state_response_norm_aborts
signing_attempt_state_response_norm_accepts
signing_attempt_state_valid_signature_accepts
valid_signature_sign_internal
verify_norm_ok_current
\end{lstlisting}

\end{formaltheorem}

\begin{formaltheorem}[EasyCrypt line 1674]
\noindent\textbf{EasyCrypt name.}
\begin{lstlisting}[language=EasyCrypt,basicstyle=\ttfamily\scriptsize]
signing_attempt_state_distribution_lossless
\end{lstlisting}
\noindent\textbf{Checked EasyCrypt statement.}
\begin{lstlisting}[language=EasyCrypt,basicstyle=\ttfamily\scriptsize]
lemma signing_attempt_state_distribution_lossless md sk m ctx :
  is_lossless (dsigning_attempt_state md sk m ctx).
\end{lstlisting}
\noindent\textbf{Formal mathematical reading.}
Mathematical theorem. This lemma is a universally quantified proposition over the variables shown in the EasyCrypt statement.

\noindent\textbf{Role in the proof.}
Use in this chapter. It contributes directly to an advantage bound, bad-event bound, or real-arithmetic composition step.

\noindent\textbf{Proof-script interpretation.}
Proof-script reading. The machine proof decomposes the theorem into rewrites, program-logic obligations, relational equivalences, and arithmetic side conditions.
\emph{Main tactics visible in the proof script:} \ec{rewrite}, \ec{apply}.
\emph{Named identifiers syntactically cited in the proof block:}
\begin{lstlisting}[language=EasyCrypt,basicstyle=\ttfamily\scriptsize]
dmap_ll
dsigning_attempt_state
signing_coin_distribution_lossless
\end{lstlisting}

\end{formaltheorem}

\begin{formaltheorem}[EasyCrypt line 1681]
\noindent\textbf{EasyCrypt name.}
\begin{lstlisting}[language=EasyCrypt,basicstyle=\ttfamily\scriptsize]
signing_attempt_state_distribution_token_spread
\end{lstlisting}
\noindent\textbf{Checked EasyCrypt statement.}
\begin{lstlisting}[language=EasyCrypt,basicstyle=\ttfamily\scriptsize]
lemma signing_attempt_state_distribution_token_spread
  md sk m ctx token :
  mu (dsigning_attempt_state md sk m ctx)
    (fun st =>
      signing_entropy_token_of_coins
        (signing_attempt_state_coins st) = token) <=
  signing_randomness_point_bound.
\end{lstlisting}
\noindent\textbf{Formal mathematical reading.}
Mathematical theorem. This is an inequality, usually an arithmetic loss bound, event inclusion bound, or monotonicity fact used to compose probabilities.

\noindent\textbf{Role in the proof.}
Use in this chapter. It preserves an invariant across an oracle call, adversary call, or game procedure.

\noindent\textbf{Proof-script interpretation.}
Proof-script reading. The machine proof decomposes the theorem into rewrites, program-logic obligations, relational equivalences, and arithmetic side conditions.
\emph{Main tactics visible in the proof script:} \ec{rewrite}, \ec{apply}.
\emph{Named identifiers syntactically cited in the proof block:}
\begin{lstlisting}[language=EasyCrypt,basicstyle=\ttfamily\scriptsize]
coins
dmapE
drandom_coins
dsigning_attempt_state
ler_trans
mu
mu_le
signing_attempt_state_coins
signing_attempt_state_of_coins
signing_coin_distribution_token_point_bound
signing_entropy_token_of_coins
token
\end{lstlisting}

\end{formaltheorem}

\begin{formaltheorem}[EasyCrypt line 1698]
\noindent\textbf{EasyCrypt name.}
\begin{lstlisting}[language=EasyCrypt,basicstyle=\ttfamily\scriptsize]
signing_attempt_state_distribution_reference_rejection_abort_zero_current
\end{lstlisting}
\noindent\textbf{Checked EasyCrypt statement.}
\begin{lstlisting}[language=EasyCrypt,basicstyle=\ttfamily\scriptsize]
lemma signing_attempt_state_distribution_reference_rejection_abort_zero_current
  md sk m ctx :
  mu (dsigning_attempt_state md sk m ctx)
    (signing_attempt_state_reference_rejection_aborts md) =
  0%r.
\end{lstlisting}
\noindent\textbf{Formal mathematical reading.}
Mathematical theorem. This is an equality or definitional expansion used for rewriting and normalization.

\noindent\textbf{Role in the proof.}
Use in this chapter. It preserves an invariant across an oracle call, adversary call, or game procedure.

\noindent\textbf{Proof-script interpretation.}
Proof-script reading. The machine proof decomposes the theorem into rewrites, program-logic obligations, relational equivalences, and arithmetic side conditions.
\emph{Main tactics visible in the proof script:} \ec{rewrite}, \ec{smt}.
\emph{Named identifiers syntactically cited in the proof block:}
\begin{lstlisting}[language=EasyCrypt,basicstyle=\ttfamily\scriptsize]
coins
dmapE
dsigning_attempt_state
mu0
mu_eq
random_coins
signing_attempt_state_of_coins_reference_rejection_accepts_current
signing_attempt_state_reference_rejection_aborts
\end{lstlisting}

\end{formaltheorem}

\begin{formaltheorem}[EasyCrypt line 1712]
\noindent\textbf{EasyCrypt name.}
\begin{lstlisting}[language=EasyCrypt,basicstyle=\ttfamily\scriptsize]
signing_attempt_state_distribution_rejection_abort_zero_current
\end{lstlisting}
\noindent\textbf{Checked EasyCrypt statement.}
\begin{lstlisting}[language=EasyCrypt,basicstyle=\ttfamily\scriptsize]
lemma signing_attempt_state_distribution_rejection_abort_zero_current
  md sk m ctx :
  mu (dsigning_attempt_state md sk m ctx)
    (signing_attempt_state_rejection_aborts md
       (public_key_of_secret md sk) m ctx) =
  0%r.
\end{lstlisting}
\noindent\textbf{Formal mathematical reading.}
Mathematical theorem. This is an equality or definitional expansion used for rewriting and normalization.

\noindent\textbf{Role in the proof.}
Use in this chapter. It preserves an invariant across an oracle call, adversary call, or game procedure.

\noindent\textbf{Proof-script interpretation.}
Proof-script reading. The machine proof decomposes the theorem into rewrites, program-logic obligations, relational equivalences, and arithmetic side conditions.
\emph{Main tactics visible in the proof script:} \ec{rewrite}, \ec{smt}.
\emph{Named identifiers syntactically cited in the proof block:}
\begin{lstlisting}[language=EasyCrypt,basicstyle=\ttfamily\scriptsize]
coins
dmapE
dsigning_attempt_state
mu0
mu_eq
random_coins
signing_attempt_state_of_coins_rejection_no_abort_current
\end{lstlisting}

\end{formaltheorem}

\begin{formaltheorem}[EasyCrypt line 1726]
\noindent\textbf{EasyCrypt name.}
\begin{lstlisting}[language=EasyCrypt,basicstyle=\ttfamily\scriptsize]
sign_internal_simulation_relation
\end{lstlisting}
\noindent\textbf{Checked EasyCrypt statement.}
\begin{lstlisting}[language=EasyCrypt,basicstyle=\ttfamily\scriptsize]
lemma sign_internal_simulation_relation md sk m ctx coins :
  signing_simulation_relation md
    (public_key_of_secret md sk) m ctx
    (sign_internal md sk m ctx coins).
\end{lstlisting}
\noindent\textbf{Formal mathematical reading.}
Mathematical theorem. This lemma is a universally quantified proposition over the variables shown in the EasyCrypt statement.

\noindent\textbf{Role in the proof.}
Use in this chapter. It proves an exact game replacement or simulator equivalence.

\noindent\textbf{Proof-script interpretation.}
Proof-script reading. The machine proof decomposes the theorem into rewrites, program-logic obligations, relational equivalences, and arithmetic side conditions.
\emph{Main tactics visible in the proof script:} \ec{apply}.
\emph{Named identifiers syntactically cited in the proof block:}
\begin{lstlisting}[language=EasyCrypt,basicstyle=\ttfamily\scriptsize]
coins
ctx
md
public_key_of_secret
sign_internal
sign_internal_attempt_relation
signing_attempt_simulates
sk
\end{lstlisting}

\end{formaltheorem}

\begin{formaltheorem}[EasyCrypt line 1736]
\noindent\textbf{EasyCrypt name.}
\begin{lstlisting}[language=EasyCrypt,basicstyle=\ttfamily\scriptsize]
signing_attempt_transcript_valid
\end{lstlisting}
\noindent\textbf{Checked EasyCrypt statement.}
\begin{lstlisting}[language=EasyCrypt,basicstyle=\ttfamily\scriptsize]
lemma signing_attempt_transcript_valid md pk sk m ctx coins sig :
  signing_attempt_relation md pk sk m ctx coins sig =>
  transcript_valid md (transcript_of_signature md pk m ctx sig).
\end{lstlisting}
\noindent\textbf{Formal mathematical reading.}
Mathematical theorem. This is an implication, normally turning one invariant, validity predicate, or proof obligation into another.

\noindent\textbf{Role in the proof.}
Use in this chapter. It proves a structural correctness fact needed by later extraction or verification arguments.

\noindent\textbf{Proof-script interpretation.}
Proof-script reading. The machine proof decomposes the theorem into rewrites, program-logic obligations, relational equivalences, and arithmetic side conditions.
\emph{Main tactics visible in the proof script:} \ec{rewrite}, \ec{split}, \ec{apply}.
\emph{Named identifiers syntactically cited in the proof block:}
\begin{lstlisting}[language=EasyCrypt,basicstyle=\ttfamily\scriptsize]
pkE
signing_attempt_relation
transcript_context
transcript_fields_match_signature
transcript_message
transcript_of_signature
transcript_pk
transcript_signature
transcript_valid
transcript_verifies
valid_keypair
valid_signature_sign_internal
verify_internal_sign_internal
\end{lstlisting}

\end{formaltheorem}

\begin{formaltheorem}[EasyCrypt line 1753]
\noindent\textbf{EasyCrypt name.}
\begin{lstlisting}[language=EasyCrypt,basicstyle=\ttfamily\scriptsize]
signing_attempt_transcript_relation
\end{lstlisting}
\noindent\textbf{Checked EasyCrypt statement.}
\begin{lstlisting}[language=EasyCrypt,basicstyle=\ttfamily\scriptsize]
lemma signing_attempt_transcript_relation md pk sk m ctx coins sig :
  signing_attempt_relation md pk sk m ctx coins sig =>
  signing_transcript_relation md pk m ctx sig
    (transcript_of_signature md pk m ctx sig).
\end{lstlisting}
\noindent\textbf{Formal mathematical reading.}
Mathematical theorem. This is an implication, normally turning one invariant, validity predicate, or proof obligation into another.

\noindent\textbf{Role in the proof.}
Use in this chapter. It is a reusable checked theorem cited by later proof scripts.

\noindent\textbf{Proof-script interpretation.}
Proof-script reading. The machine proof decomposes the theorem into rewrites, program-logic obligations, relational equivalences, and arithmetic side conditions.
\emph{Main tactics visible in the proof script:} \ec{rewrite}, \ec{split}, \ec{apply}.
\emph{Named identifiers syntactically cited in the proof block:}
\begin{lstlisting}[language=EasyCrypt,basicstyle=\ttfamily\scriptsize]
attempt
coins
ctx
md
pk
sig
signing_attempt_simulates
signing_attempt_transcript_valid
signing_transcript_relation
sk
\end{lstlisting}

\end{formaltheorem}

\begin{formaltheorem}[EasyCrypt line 1766]
\noindent\textbf{EasyCrypt name.}
\begin{lstlisting}[language=EasyCrypt,basicstyle=\ttfamily\scriptsize]
sign_internal_transcript_relation
\end{lstlisting}
\noindent\textbf{Checked EasyCrypt statement.}
\begin{lstlisting}[language=EasyCrypt,basicstyle=\ttfamily\scriptsize]
lemma sign_internal_transcript_relation md sk m ctx coins :
  signing_transcript_relation md
    (public_key_of_secret md sk) m ctx
    (sign_internal md sk m ctx coins)
    (transcript_of_signature md (public_key_of_secret md sk) m ctx
      (sign_internal md sk m ctx coins)).
\end{lstlisting}
\noindent\textbf{Formal mathematical reading.}
Mathematical theorem. This lemma is a universally quantified proposition over the variables shown in the EasyCrypt statement.

\noindent\textbf{Role in the proof.}
Use in this chapter. It is a reusable checked theorem cited by later proof scripts.

\noindent\textbf{Proof-script interpretation.}
Proof-script reading. The machine proof decomposes the theorem into rewrites, program-logic obligations, relational equivalences, and arithmetic side conditions.
\emph{Main tactics visible in the proof script:} \ec{apply}.
\emph{Named identifiers syntactically cited in the proof block:}
\begin{lstlisting}[language=EasyCrypt,basicstyle=\ttfamily\scriptsize]
coins
ctx
md
public_key_of_secret
sign_internal
sign_internal_attempt_relation
signing_attempt_transcript_relation
sk
\end{lstlisting}

\end{formaltheorem}

\begin{formaltheorem}[EasyCrypt line 1779]
\noindent\textbf{EasyCrypt name.}
\begin{lstlisting}[language=EasyCrypt,basicstyle=\ttfamily\scriptsize]
sign_internal_record_relation
\end{lstlisting}
\noindent\textbf{Checked EasyCrypt statement.}
\begin{lstlisting}[language=EasyCrypt,basicstyle=\ttfamily\scriptsize]
lemma sign_internal_record_relation md sk m ctx coins :
  signing_record_relation md (public_key_of_secret md sk)
    (m, ctx,
     sign_internal md sk m ctx coins,
     transcript_of_signature md (public_key_of_secret md sk) m ctx
       (sign_internal md sk m ctx coins)).
\end{lstlisting}
\noindent\textbf{Formal mathematical reading.}
Mathematical theorem. This lemma is a universally quantified proposition over the variables shown in the EasyCrypt statement.

\noindent\textbf{Role in the proof.}
Use in this chapter. It is a reusable checked theorem cited by later proof scripts.

\noindent\textbf{Proof-script interpretation.}
Proof-script reading. The machine proof decomposes the theorem into rewrites, program-logic obligations, relational equivalences, and arithmetic side conditions.
\emph{Main tactics visible in the proof script:} \ec{rewrite}, \ec{apply}.
\emph{Named identifiers syntactically cited in the proof block:}
\begin{lstlisting}[language=EasyCrypt,basicstyle=\ttfamily\scriptsize]
sign_internal_transcript_relation
signing_record_context
signing_record_message
signing_record_relation
signing_record_signature
signing_record_transcript
\end{lstlisting}

\end{formaltheorem}

\begin{formaltheorem}[EasyCrypt line 1792]
\noindent\textbf{EasyCrypt name.}
\begin{lstlisting}[language=EasyCrypt,basicstyle=\ttfamily\scriptsize]
signing_record_relation_sound
\end{lstlisting}
\noindent\textbf{Checked EasyCrypt statement.}
\begin{lstlisting}[language=EasyCrypt,basicstyle=\ttfamily\scriptsize]
lemma signing_record_relation_sound md pk r :
  signing_record_relation md pk r =>
  signing_simulation_relation md pk
    (signing_record_message r)
    (signing_record_context r)
    (signing_record_signature r) /\
  transcript_valid md (signing_record_transcript r).
\end{lstlisting}
\noindent\textbf{Formal mathematical reading.}
Mathematical theorem. This is an implication, normally turning one invariant, validity predicate, or proof obligation into another.

\noindent\textbf{Role in the proof.}
Use in this chapter. It proves a structural correctness fact needed by later extraction or verification arguments.

\noindent\textbf{Proof-script interpretation.}
Proof-script reading. The machine proof decomposes the theorem into rewrites, program-logic obligations, relational equivalences, and arithmetic side conditions.
\emph{Main tactics visible in the proof script:} \ec{rewrite}, \ec{case}, \ec{split}, \ec{apply}.
\emph{Named identifiers syntactically cited in the proof block:}
\begin{lstlisting}[language=EasyCrypt,basicstyle=\ttfamily\scriptsize]
hsim
rest
signing_record_relation
signing_transcript_relation
valid
\end{lstlisting}

\end{formaltheorem}

\begin{formaltheorem}[EasyCrypt line 1809]
\noindent\textbf{EasyCrypt name.}
\begin{lstlisting}[language=EasyCrypt,basicstyle=\ttfamily\scriptsize]
signing_record_log_sound_cons
\end{lstlisting}
\noindent\textbf{Checked EasyCrypt statement.}
\begin{lstlisting}[language=EasyCrypt,basicstyle=\ttfamily\scriptsize]
lemma signing_record_log_sound_cons md pk r rs :
  signing_record_relation md pk r =>
  signing_record_log_sound md pk rs =>
  signing_record_log_sound md pk (r :: rs).
\end{lstlisting}
\noindent\textbf{Formal mathematical reading.}
Mathematical theorem. This is an implication, normally turning one invariant, validity predicate, or proof obligation into another.

\noindent\textbf{Role in the proof.}
Use in this chapter. It proves a structural correctness fact needed by later extraction or verification arguments.

\noindent\textbf{Proof-script interpretation.}
Proof-script reading. The machine proof decomposes the theorem into rewrites, program-logic obligations, relational equivalences, and arithmetic side conditions.
\emph{Main tactics visible in the proof script:} \ec{rewrite}.
\emph{Named identifiers syntactically cited in the proof block:}
\begin{lstlisting}[language=EasyCrypt,basicstyle=\ttfamily\scriptsize]
signing_record_log_sound
\end{lstlisting}

\end{formaltheorem}

\begin{formaltheorem}[EasyCrypt line 1817]
\noindent\textbf{EasyCrypt name.}
\begin{lstlisting}[language=EasyCrypt,basicstyle=\ttfamily\scriptsize]
signing_record_log_sound_transcripts_valid
\end{lstlisting}
\noindent\textbf{Checked EasyCrypt statement.}
\begin{lstlisting}[language=EasyCrypt,basicstyle=\ttfamily\scriptsize]
lemma signing_record_log_sound_transcripts_valid md pk rs :
  signing_record_log_sound md pk rs =>
  transcript_log_valid md (signing_record_transcripts rs).
\end{lstlisting}
\noindent\textbf{Formal mathematical reading.}
Mathematical theorem. This is an implication, normally turning one invariant, validity predicate, or proof obligation into another.

\noindent\textbf{Role in the proof.}
Use in this chapter. It proves a structural correctness fact needed by later extraction or verification arguments.

\noindent\textbf{Proof-script interpretation.}
Proof-script reading. The machine proof decomposes the theorem into rewrites, program-logic obligations, relational equivalences, and arithmetic side conditions.
\emph{Main tactics visible in the proof script:} \ec{rewrite}, \ec{elim}, \ec{split}, \ec{case}, \ec{apply}.
\emph{Named identifiers syntactically cited in the proof block:}
\begin{lstlisting}[language=EasyCrypt,basicstyle=\ttfamily\scriptsize]
ih
md
pk
rel
rs
signing_record_log_sound
signing_record_relation_sound
signing_record_transcripts
sound
transcript_log_valid
\end{lstlisting}

\end{formaltheorem}

\begin{formaltheorem}[EasyCrypt line 1830]
\noindent\textbf{EasyCrypt name.}
\begin{lstlisting}[language=EasyCrypt,basicstyle=\ttfamily\scriptsize]
signing_record_queries_cons
\end{lstlisting}
\noindent\textbf{Checked EasyCrypt statement.}
\begin{lstlisting}[language=EasyCrypt,basicstyle=\ttfamily\scriptsize]
lemma signing_record_queries_cons r rs :
  signing_record_queries (r :: rs) =
  signing_record_query r :: signing_record_queries rs.
\end{lstlisting}
\noindent\textbf{Formal mathematical reading.}
Mathematical theorem. This is an equality or definitional expansion used for rewriting and normalization.

\noindent\textbf{Role in the proof.}
Use in this chapter. It is a reusable checked theorem cited by later proof scripts.

\noindent\textbf{Proof-script interpretation.}
No proof block was collected for this item.

\end{formaltheorem}

\begin{formaltheorem}[EasyCrypt line 1835]
\noindent\textbf{EasyCrypt name.}
\begin{lstlisting}[language=EasyCrypt,basicstyle=\ttfamily\scriptsize]
signing_record_transcripts_cons
\end{lstlisting}
\noindent\textbf{Checked EasyCrypt statement.}
\begin{lstlisting}[language=EasyCrypt,basicstyle=\ttfamily\scriptsize]
lemma signing_record_transcripts_cons r rs :
  signing_record_transcripts (r :: rs) =
  signing_record_transcript r :: signing_record_transcripts rs.
\end{lstlisting}
\noindent\textbf{Formal mathematical reading.}
Mathematical theorem. This is an equality or definitional expansion used for rewriting and normalization.

\noindent\textbf{Role in the proof.}
Use in this chapter. It is a reusable checked theorem cited by later proof scripts.

\noindent\textbf{Proof-script interpretation.}
No proof block was collected for this item.

\end{formaltheorem}

\begin{formaltheorem}[EasyCrypt line 1840]
\noindent\textbf{EasyCrypt name.}
\begin{lstlisting}[language=EasyCrypt,basicstyle=\ttfamily\scriptsize]
rejection_bound_parts_nonnegative
\end{lstlisting}
\noindent\textbf{Checked EasyCrypt statement.}
\begin{lstlisting}[language=EasyCrypt,basicstyle=\ttfamily\scriptsize]
lemma rejection_bound_parts_nonnegative :
  rejection_sampling_bound_obligation =>
  0%r <= rejection_sampling_loss_term.
\end{lstlisting}
\noindent\textbf{Formal mathematical reading.}
Mathematical theorem. This is an inequality, usually an arithmetic loss bound, event inclusion bound, or monotonicity fact used to compose probabilities.

\noindent\textbf{Role in the proof.}
Use in this chapter. It contributes directly to an advantage bound, bad-event bound, or real-arithmetic composition step.

\noindent\textbf{Proof-script interpretation.}
No proof block was collected for this item.

\end{formaltheorem}

\begin{formaltheorem}[EasyCrypt line 1845]
\noindent\textbf{EasyCrypt name.}
\begin{lstlisting}[language=EasyCrypt,basicstyle=\ttfamily\scriptsize]
rejection_sampling_bound_obligation_holds
\end{lstlisting}
\noindent\textbf{Checked EasyCrypt statement.}
\begin{lstlisting}[language=EasyCrypt,basicstyle=\ttfamily\scriptsize]
lemma rejection_sampling_bound_obligation_holds :
  rejection_sampling_bound_obligation.
\end{lstlisting}
\noindent\textbf{Formal mathematical reading.}
Mathematical theorem. This lemma is a universally quantified proposition over the variables shown in the EasyCrypt statement.

\noindent\textbf{Role in the proof.}
Use in this chapter. It contributes directly to an advantage bound, bad-event bound, or real-arithmetic composition step.

\noindent\textbf{Proof-script interpretation.}
Proof-script reading. The machine proof decomposes the theorem into rewrites, program-logic obligations, relational equivalences, and arithmetic side conditions.
\emph{Main tactics visible in the proof script:} \ec{rewrite}, \ec{split}, \ec{apply}.
\emph{Named identifiers syntactically cited in the proof block:}
\begin{lstlisting}[language=EasyCrypt,basicstyle=\ttfamily\scriptsize]
fs_with_aborts_min_entropy_term_nonnegative
fs_with_aborts_reprogramming_term_nonnegative
rejection_sampling_bound_obligation
rejection_sampling_loss_term_nonnegative
\end{lstlisting}

\end{formaltheorem}

\begin{formaltheorem}[EasyCrypt line 1856]
\noindent\textbf{EasyCrypt name.}
\begin{lstlisting}[language=EasyCrypt,basicstyle=\ttfamily\scriptsize]
signing_attempt_state_of_sample_pair_sampled_yE
\end{lstlisting}
\noindent\textbf{Checked EasyCrypt statement.}
\begin{lstlisting}[language=EasyCrypt,basicstyle=\ttfamily\scriptsize]
lemma signing_attempt_state_of_sample_pair_sampled_yE md sk m ctx coins sample :
  signing_attempt_state_sampled_y
    (signing_attempt_state_of_sample_pair md sk m ctx coins sample) =
  commitment_raw md sk m ctx coins.
\end{lstlisting}
\noindent\textbf{Formal mathematical reading.}
Mathematical theorem. This is an equality or definitional expansion used for rewriting and normalization.

\noindent\textbf{Role in the proof.}
Use in this chapter. It contributes directly to an advantage bound, bad-event bound, or real-arithmetic composition step.

\noindent\textbf{Proof-script interpretation.}
Proof-script reading. The machine proof decomposes the theorem into rewrites, program-logic obligations, relational equivalences, and arithmetic side conditions.
\emph{Main tactics visible in the proof script:} \ec{rewrite}.
\emph{Named identifiers syntactically cited in the proof block:}
\begin{lstlisting}[language=EasyCrypt,basicstyle=\ttfamily\scriptsize]
signing_attempt_sample_commitment_raw
signing_attempt_sample_of_pair
signing_attempt_state_of_sample_pair
signing_attempt_state_sample
signing_attempt_state_sampled_y
\end{lstlisting}

\end{formaltheorem}

\begin{formaltheorem}[EasyCrypt line 1868]
\noindent\textbf{EasyCrypt name.}
\begin{lstlisting}[language=EasyCrypt,basicstyle=\ttfamily\scriptsize]
signing_attempt_state_of_sample_pair_highbitsE
\end{lstlisting}
\noindent\textbf{Checked EasyCrypt statement.}
\begin{lstlisting}[language=EasyCrypt,basicstyle=\ttfamily\scriptsize]
lemma signing_attempt_state_of_sample_pair_highbitsE md sk m ctx coins sample :
  signing_attempt_state_highbits
    (signing_attempt_state_of_sample_pair md sk m ctx coins sample) =
  commitment_highbits md sk m ctx coins.
\end{lstlisting}
\noindent\textbf{Formal mathematical reading.}
Mathematical theorem. This is an equality or definitional expansion used for rewriting and normalization.

\noindent\textbf{Role in the proof.}
Use in this chapter. It contributes directly to an advantage bound, bad-event bound, or real-arithmetic composition step.

\noindent\textbf{Proof-script interpretation.}
Proof-script reading. The machine proof decomposes the theorem into rewrites, program-logic obligations, relational equivalences, and arithmetic side conditions.
\emph{Main tactics visible in the proof script:} \ec{rewrite}.
\emph{Named identifiers syntactically cited in the proof block:}
\begin{lstlisting}[language=EasyCrypt,basicstyle=\ttfamily\scriptsize]
signing_attempt_state_highbits
signing_attempt_state_of_sample_pair
\end{lstlisting}

\end{formaltheorem}

\begin{formaltheorem}[EasyCrypt line 1877]
\noindent\textbf{EasyCrypt name.}
\begin{lstlisting}[language=EasyCrypt,basicstyle=\ttfamily\scriptsize]
signing_attempt_state_of_sample_pair_lowbitsE
\end{lstlisting}
\noindent\textbf{Checked EasyCrypt statement.}
\begin{lstlisting}[language=EasyCrypt,basicstyle=\ttfamily\scriptsize]
lemma signing_attempt_state_of_sample_pair_lowbitsE md sk m ctx coins sample :
  signing_attempt_state_lowbits
    (signing_attempt_state_of_sample_pair md sk m ctx coins sample) =
  commitment_lowbits md sk m ctx coins.
\end{lstlisting}
\noindent\textbf{Formal mathematical reading.}
Mathematical theorem. This is an equality or definitional expansion used for rewriting and normalization.

\noindent\textbf{Role in the proof.}
Use in this chapter. It contributes directly to an advantage bound, bad-event bound, or real-arithmetic composition step.

\noindent\textbf{Proof-script interpretation.}
Proof-script reading. The machine proof decomposes the theorem into rewrites, program-logic obligations, relational equivalences, and arithmetic side conditions.
\emph{Main tactics visible in the proof script:} \ec{rewrite}.
\emph{Named identifiers syntactically cited in the proof block:}
\begin{lstlisting}[language=EasyCrypt,basicstyle=\ttfamily\scriptsize]
signing_attempt_state_lowbits
signing_attempt_state_of_sample_pair
\end{lstlisting}

\end{formaltheorem}

\begin{formaltheorem}[EasyCrypt line 1886]
\noindent\textbf{EasyCrypt name.}
\begin{lstlisting}[language=EasyCrypt,basicstyle=\ttfamily\scriptsize]
signing_attempt_state_of_sample_pair_challengeE
\end{lstlisting}
\noindent\textbf{Checked EasyCrypt statement.}
\begin{lstlisting}[language=EasyCrypt,basicstyle=\ttfamily\scriptsize]
lemma signing_attempt_state_of_sample_pair_challengeE md sk m ctx coins sample :
  signing_attempt_state_challenge
    (signing_attempt_state_of_sample_pair md sk m ctx coins sample) =
  commitment_challenge md sk m ctx coins.
\end{lstlisting}
\noindent\textbf{Formal mathematical reading.}
Mathematical theorem. This is an equality or definitional expansion used for rewriting and normalization.

\noindent\textbf{Role in the proof.}
Use in this chapter. It contributes directly to an advantage bound, bad-event bound, or real-arithmetic composition step.

\noindent\textbf{Proof-script interpretation.}
Proof-script reading. The machine proof decomposes the theorem into rewrites, program-logic obligations, relational equivalences, and arithmetic side conditions.
\emph{Main tactics visible in the proof script:} \ec{rewrite}.
\emph{Named identifiers syntactically cited in the proof block:}
\begin{lstlisting}[language=EasyCrypt,basicstyle=\ttfamily\scriptsize]
signing_attempt_state_challenge
signing_attempt_state_of_sample_pair
\end{lstlisting}

\end{formaltheorem}

\begin{formaltheorem}[EasyCrypt line 1895]
\noindent\textbf{EasyCrypt name.}
\begin{lstlisting}[language=EasyCrypt,basicstyle=\ttfamily\scriptsize]
signing_attempt_state_of_sample_pair_responseE
\end{lstlisting}
\noindent\textbf{Checked EasyCrypt statement.}
\begin{lstlisting}[language=EasyCrypt,basicstyle=\ttfamily\scriptsize]
lemma signing_attempt_state_of_sample_pair_responseE md sk m ctx coins sample :
  signing_attempt_state_response
    (signing_attempt_state_of_sample_pair md sk m ctx coins sample) =
  response_vector md sk m ctx coins.
\end{lstlisting}
\noindent\textbf{Formal mathematical reading.}
Mathematical theorem. This is an equality or definitional expansion used for rewriting and normalization.

\noindent\textbf{Role in the proof.}
Use in this chapter. It contributes directly to an advantage bound, bad-event bound, or real-arithmetic composition step.

\noindent\textbf{Proof-script interpretation.}
Proof-script reading. The machine proof decomposes the theorem into rewrites, program-logic obligations, relational equivalences, and arithmetic side conditions.
\emph{Main tactics visible in the proof script:} \ec{rewrite}.
\emph{Named identifiers syntactically cited in the proof block:}
\begin{lstlisting}[language=EasyCrypt,basicstyle=\ttfamily\scriptsize]
signing_attempt_state_of_sample_pair
signing_attempt_state_response
\end{lstlisting}

\end{formaltheorem}

\begin{formaltheorem}[EasyCrypt line 1904]
\noindent\textbf{EasyCrypt name.}
\begin{lstlisting}[language=EasyCrypt,basicstyle=\ttfamily\scriptsize]
signing_attempt_state_of_sample_pair_response_auxE
\end{lstlisting}
\noindent\textbf{Checked EasyCrypt statement.}
\begin{lstlisting}[language=EasyCrypt,basicstyle=\ttfamily\scriptsize]
lemma signing_attempt_state_of_sample_pair_response_auxE
   md sk m ctx coins sample :
  signing_attempt_state_response_aux
    (signing_attempt_state_of_sample_pair md sk m ctx coins sample) =
  response_aux_vector md sk m ctx coins.
\end{lstlisting}
\noindent\textbf{Formal mathematical reading.}
Mathematical theorem. This is an equality or definitional expansion used for rewriting and normalization.

\noindent\textbf{Role in the proof.}
Use in this chapter. It contributes directly to an advantage bound, bad-event bound, or real-arithmetic composition step.

\noindent\textbf{Proof-script interpretation.}
Proof-script reading. The machine proof decomposes the theorem into rewrites, program-logic obligations, relational equivalences, and arithmetic side conditions.
\emph{Main tactics visible in the proof script:} \ec{rewrite}.
\emph{Named identifiers syntactically cited in the proof block:}
\begin{lstlisting}[language=EasyCrypt,basicstyle=\ttfamily\scriptsize]
signing_attempt_state_of_sample_pair
signing_attempt_state_response_aux
\end{lstlisting}

\end{formaltheorem}

\begin{formaltheorem}[EasyCrypt line 1914]
\noindent\textbf{EasyCrypt name.}
\begin{lstlisting}[language=EasyCrypt,basicstyle=\ttfamily\scriptsize]
signing_attempt_state_of_sample_pair_hintE
\end{lstlisting}
\noindent\textbf{Checked EasyCrypt statement.}
\begin{lstlisting}[language=EasyCrypt,basicstyle=\ttfamily\scriptsize]
lemma signing_attempt_state_of_sample_pair_hintE md sk m ctx coins sample :
  signing_attempt_state_hint
    (signing_attempt_state_of_sample_pair md sk m ctx coins sample) =
  hint_vector md sk m ctx coins.
\end{lstlisting}
\noindent\textbf{Formal mathematical reading.}
Mathematical theorem. This is an equality or definitional expansion used for rewriting and normalization.

\noindent\textbf{Role in the proof.}
Use in this chapter. It contributes directly to an advantage bound, bad-event bound, or real-arithmetic composition step.

\noindent\textbf{Proof-script interpretation.}
Proof-script reading. The machine proof decomposes the theorem into rewrites, program-logic obligations, relational equivalences, and arithmetic side conditions.
\emph{Main tactics visible in the proof script:} \ec{rewrite}.
\emph{Named identifiers syntactically cited in the proof block:}
\begin{lstlisting}[language=EasyCrypt,basicstyle=\ttfamily\scriptsize]
signing_attempt_state_hint
signing_attempt_state_of_sample_pair
\end{lstlisting}

\end{formaltheorem}

\begin{formaltheorem}[EasyCrypt line 1923]
\noindent\textbf{EasyCrypt name.}
\begin{lstlisting}[language=EasyCrypt,basicstyle=\ttfamily\scriptsize]
signing_attempt_state_of_sample_pair_lowbits_carrierE
\end{lstlisting}
\noindent\textbf{Checked EasyCrypt statement.}
\begin{lstlisting}[language=EasyCrypt,basicstyle=\ttfamily\scriptsize]
lemma signing_attempt_state_of_sample_pair_lowbits_carrierE
   md sk m ctx coins sample :
  signing_attempt_state_lowbits_carrier
    (signing_attempt_state_of_sample_pair md sk m ctx coins sample) =
  nth poly_zero (commitment_raw md sk m ctx coins) 0.
\end{lstlisting}
\noindent\textbf{Formal mathematical reading.}
Mathematical theorem. This is an equality or definitional expansion used for rewriting and normalization.

\noindent\textbf{Role in the proof.}
Use in this chapter. It contributes directly to an advantage bound, bad-event bound, or real-arithmetic composition step.

\noindent\textbf{Proof-script interpretation.}
Proof-script reading. The machine proof decomposes the theorem into rewrites, program-logic obligations, relational equivalences, and arithmetic side conditions.
\emph{Main tactics visible in the proof script:} \ec{rewrite}.
\emph{Named identifiers syntactically cited in the proof block:}
\begin{lstlisting}[language=EasyCrypt,basicstyle=\ttfamily\scriptsize]
signing_attempt_state_lowbits_carrier
signing_attempt_state_of_sample_pair_sampled_yE
\end{lstlisting}

\end{formaltheorem}

\begin{formaltheorem}[EasyCrypt line 1933]
\noindent\textbf{EasyCrypt name.}
\begin{lstlisting}[language=EasyCrypt,basicstyle=\ttfamily\scriptsize]
signing_attempt_state_of_sample_pair_tokenE
\end{lstlisting}
\noindent\textbf{Checked EasyCrypt statement.}
\begin{lstlisting}[language=EasyCrypt,basicstyle=\ttfamily\scriptsize]
lemma signing_attempt_state_of_sample_pair_tokenE md sk m ctx coins sample :
  signing_attempt_state_token
    (signing_attempt_state_of_sample_pair md sk m ctx coins sample) =
  signature_token md sk m ctx coins.
\end{lstlisting}
\noindent\textbf{Formal mathematical reading.}
Mathematical theorem. This is an equality or definitional expansion used for rewriting and normalization.

\noindent\textbf{Role in the proof.}
Use in this chapter. It contributes directly to an advantage bound, bad-event bound, or real-arithmetic composition step.

\noindent\textbf{Proof-script interpretation.}
Proof-script reading. The machine proof decomposes the theorem into rewrites, program-logic obligations, relational equivalences, and arithmetic side conditions.
\emph{Main tactics visible in the proof script:} \ec{rewrite}.
\emph{Named identifiers syntactically cited in the proof block:}
\begin{lstlisting}[language=EasyCrypt,basicstyle=\ttfamily\scriptsize]
signature_token
signing_attempt_state_of_sample_pair_hintE
signing_attempt_state_of_sample_pair_responseE
signing_attempt_state_of_sample_pair_response_auxE
signing_attempt_state_token
\end{lstlisting}

\end{formaltheorem}

\begin{formaltheorem}[EasyCrypt line 1945]
\noindent\textbf{EasyCrypt name.}
\begin{lstlisting}[language=EasyCrypt,basicstyle=\ttfamily\scriptsize]
signing_attempt_state_of_sample_pair_programmed_lowbitsE
\end{lstlisting}
\noindent\textbf{Checked EasyCrypt statement.}
\begin{lstlisting}[language=EasyCrypt,basicstyle=\ttfamily\scriptsize]
lemma signing_attempt_state_of_sample_pair_programmed_lowbitsE
   md sk m ctx coins sample :
  signing_attempt_state_programmed_lowbits
    (signing_attempt_state_of_sample_pair md sk m ctx coins sample) =
  commitment_lowbits md sk m ctx coins.
\end{lstlisting}
\noindent\textbf{Formal mathematical reading.}
Mathematical theorem. This is an equality or definitional expansion used for rewriting and normalization.

\noindent\textbf{Role in the proof.}
Use in this chapter. It contributes directly to an advantage bound, bad-event bound, or real-arithmetic composition step.

\noindent\textbf{Proof-script interpretation.}
Proof-script reading. The machine proof decomposes the theorem into rewrites, program-logic obligations, relational equivalences, and arithmetic side conditions.
\emph{Main tactics visible in the proof script:} \ec{rewrite}.
\emph{Named identifiers syntactically cited in the proof block:}
\begin{lstlisting}[language=EasyCrypt,basicstyle=\ttfamily\scriptsize]
commitment_lowbits
signing_attempt_state_coins
signing_attempt_state_of_sample_pair
signing_attempt_state_of_sample_pair_lowbits_carrierE
signing_attempt_state_programmed_lowbits
\end{lstlisting}

\end{formaltheorem}

\begin{formaltheorem}[EasyCrypt line 1957]
\noindent\textbf{EasyCrypt name.}
\begin{lstlisting}[language=EasyCrypt,basicstyle=\ttfamily\scriptsize]
signing_attempt_state_of_sample_pair_lowbits_consistent
\end{lstlisting}
\noindent\textbf{Checked EasyCrypt statement.}
\begin{lstlisting}[language=EasyCrypt,basicstyle=\ttfamily\scriptsize]
lemma signing_attempt_state_of_sample_pair_lowbits_consistent
   md sk m ctx coins sample :
  signing_attempt_state_lowbits_consistent
    (signing_attempt_state_of_sample_pair md sk m ctx coins sample).
\end{lstlisting}
\noindent\textbf{Formal mathematical reading.}
Mathematical theorem. This lemma is a universally quantified proposition over the variables shown in the EasyCrypt statement.

\noindent\textbf{Role in the proof.}
Use in this chapter. It contributes directly to an advantage bound, bad-event bound, or real-arithmetic composition step.

\noindent\textbf{Proof-script interpretation.}
Proof-script reading. The machine proof decomposes the theorem into rewrites, program-logic obligations, relational equivalences, and arithmetic side conditions.
\emph{Main tactics visible in the proof script:} \ec{rewrite}.
\emph{Named identifiers syntactically cited in the proof block:}
\begin{lstlisting}[language=EasyCrypt,basicstyle=\ttfamily\scriptsize]
signing_attempt_state_lowbits_consistent
signing_attempt_state_of_sample_pair_lowbitsE
signing_attempt_state_of_sample_pair_programmed_lowbitsE
\end{lstlisting}

\end{formaltheorem}

\begin{formaltheorem}[EasyCrypt line 1967]
\noindent\textbf{EasyCrypt name.}
\begin{lstlisting}[language=EasyCrypt,basicstyle=\ttfamily\scriptsize]
signing_attempt_state_of_sample_pair_programmed_challengeE
\end{lstlisting}
\noindent\textbf{Checked EasyCrypt statement.}
\begin{lstlisting}[language=EasyCrypt,basicstyle=\ttfamily\scriptsize]
lemma signing_attempt_state_of_sample_pair_programmed_challengeE
   md sk m ctx coins sample :
  signing_attempt_state_programmed_challenge md (public_key_of_secret md sk)
    m ctx (signing_attempt_state_of_sample_pair md sk m ctx coins sample) =
  commitment_challenge md sk m ctx coins.
\end{lstlisting}
\noindent\textbf{Formal mathematical reading.}
Mathematical theorem. This is an equality or definitional expansion used for rewriting and normalization.

\noindent\textbf{Role in the proof.}
Use in this chapter. It contributes directly to an advantage bound, bad-event bound, or real-arithmetic composition step.

\noindent\textbf{Proof-script interpretation.}
Proof-script reading. The machine proof decomposes the theorem into rewrites, program-logic obligations, relational equivalences, and arithmetic side conditions.
\emph{Main tactics visible in the proof script:} \ec{rewrite}.
\emph{Named identifiers syntactically cited in the proof block:}
\begin{lstlisting}[language=EasyCrypt,basicstyle=\ttfamily\scriptsize]
commitment_challenge
signing_attempt_state_of_sample_pair_lowbitsE
signing_attempt_state_of_sample_pair_response_auxE
signing_attempt_state_programmed_challenge
\end{lstlisting}

\end{formaltheorem}

\begin{formaltheorem}[EasyCrypt line 1979]
\noindent\textbf{EasyCrypt name.}
\begin{lstlisting}[language=EasyCrypt,basicstyle=\ttfamily\scriptsize]
signing_attempt_state_of_sample_pair_challenge_consistent
\end{lstlisting}
\noindent\textbf{Checked EasyCrypt statement.}
\begin{lstlisting}[language=EasyCrypt,basicstyle=\ttfamily\scriptsize]
lemma signing_attempt_state_of_sample_pair_challenge_consistent
   md sk m ctx coins sample :
  signing_attempt_state_challenge_consistent md (public_key_of_secret md sk)
    m ctx (signing_attempt_state_of_sample_pair md sk m ctx coins sample).
\end{lstlisting}
\noindent\textbf{Formal mathematical reading.}
Mathematical theorem. This lemma is a universally quantified proposition over the variables shown in the EasyCrypt statement.

\noindent\textbf{Role in the proof.}
Use in this chapter. It contributes directly to an advantage bound, bad-event bound, or real-arithmetic composition step.

\noindent\textbf{Proof-script interpretation.}
Proof-script reading. The machine proof decomposes the theorem into rewrites, program-logic obligations, relational equivalences, and arithmetic side conditions.
\emph{Main tactics visible in the proof script:} \ec{rewrite}.
\emph{Named identifiers syntactically cited in the proof block:}
\begin{lstlisting}[language=EasyCrypt,basicstyle=\ttfamily\scriptsize]
signing_attempt_state_challenge_consistent
signing_attempt_state_of_sample_pair_challengeE
signing_attempt_state_of_sample_pair_programmed_challengeE
\end{lstlisting}

\end{formaltheorem}

\begin{formaltheorem}[EasyCrypt line 1989]
\noindent\textbf{EasyCrypt name.}
\begin{lstlisting}[language=EasyCrypt,basicstyle=\ttfamily\scriptsize]
signing_attempt_state_of_sample_pair_reconstructed_highbitsE
\end{lstlisting}
\noindent\textbf{Checked EasyCrypt statement.}
\begin{lstlisting}[language=EasyCrypt,basicstyle=\ttfamily\scriptsize]
lemma signing_attempt_state_of_sample_pair_reconstructed_highbitsE
   md sk m ctx coins sample :
  signing_attempt_state_reconstructed_highbits md (public_key_of_secret md sk)
    (signing_attempt_state_of_sample_pair md sk m ctx coins sample) =
  commitment_highbits md sk m ctx coins.
\end{lstlisting}
\noindent\textbf{Formal mathematical reading.}
Mathematical theorem. This is an equality or definitional expansion used for rewriting and normalization.

\noindent\textbf{Role in the proof.}
Use in this chapter. It contributes directly to an advantage bound, bad-event bound, or real-arithmetic composition step.

\noindent\textbf{Proof-script interpretation.}
Proof-script reading. The machine proof decomposes the theorem into rewrites, program-logic obligations, relational equivalences, and arithmetic side conditions.
\emph{Main tactics visible in the proof script:} \ec{rewrite}.
\emph{Named identifiers syntactically cited in the proof block:}
\begin{lstlisting}[language=EasyCrypt,basicstyle=\ttfamily\scriptsize]
commitment_highbits
signing_attempt_state_of_sample_pair_challengeE
signing_attempt_state_of_sample_pair_response_auxE
signing_attempt_state_reconstructed_highbits
\end{lstlisting}

\end{formaltheorem}

\begin{formaltheorem}[EasyCrypt line 2001]
\noindent\textbf{EasyCrypt name.}
\begin{lstlisting}[language=EasyCrypt,basicstyle=\ttfamily\scriptsize]
signing_attempt_state_of_sample_pair_public_equation_consistent
\end{lstlisting}
\noindent\textbf{Checked EasyCrypt statement.}
\begin{lstlisting}[language=EasyCrypt,basicstyle=\ttfamily\scriptsize]
lemma signing_attempt_state_of_sample_pair_public_equation_consistent
   md sk m ctx coins sample :
  signing_attempt_state_public_equation_consistent md
    (public_key_of_secret md sk)
    (signing_attempt_state_of_sample_pair md sk m ctx coins sample).
\end{lstlisting}
\noindent\textbf{Formal mathematical reading.}
Mathematical theorem. This lemma is a universally quantified proposition over the variables shown in the EasyCrypt statement.

\noindent\textbf{Role in the proof.}
Use in this chapter. It contributes directly to an advantage bound, bad-event bound, or real-arithmetic composition step.

\noindent\textbf{Proof-script interpretation.}
Proof-script reading. The machine proof decomposes the theorem into rewrites, program-logic obligations, relational equivalences, and arithmetic side conditions.
\emph{Main tactics visible in the proof script:} \ec{rewrite}.
\emph{Named identifiers syntactically cited in the proof block:}
\begin{lstlisting}[language=EasyCrypt,basicstyle=\ttfamily\scriptsize]
signing_attempt_state_of_sample_pair_highbitsE
signing_attempt_state_of_sample_pair_reconstructed_highbitsE
signing_attempt_state_public_equation_consistent
\end{lstlisting}

\end{formaltheorem}

\begin{formaltheorem}[EasyCrypt line 2012]
\noindent\textbf{EasyCrypt name.}
\begin{lstlisting}[language=EasyCrypt,basicstyle=\ttfamily\scriptsize]
signing_attempt_state_of_sample_pair_signature_of_fieldsE
\end{lstlisting}
\noindent\textbf{Checked EasyCrypt statement.}
\begin{lstlisting}[language=EasyCrypt,basicstyle=\ttfamily\scriptsize]
lemma signing_attempt_state_of_sample_pair_signature_of_fieldsE
   md sk m ctx coins sample :
  signing_attempt_state_signature_of_fields
    (public_key_of_secret md sk) m ctx
    (signing_attempt_state_of_sample_pair md sk m ctx coins sample) =
  sign_internal md sk m ctx coins.
\end{lstlisting}
\noindent\textbf{Formal mathematical reading.}
Mathematical theorem. This is an equality or definitional expansion used for rewriting and normalization.

\noindent\textbf{Role in the proof.}
Use in this chapter. It contributes directly to an advantage bound, bad-event bound, or real-arithmetic composition step.

\noindent\textbf{Proof-script interpretation.}
Proof-script reading. The machine proof decomposes the theorem into rewrites, program-logic obligations, relational equivalences, and arithmetic side conditions.
\emph{Main tactics visible in the proof script:} \ec{rewrite}.
\emph{Named identifiers syntactically cited in the proof block:}
\begin{lstlisting}[language=EasyCrypt,basicstyle=\ttfamily\scriptsize]
sign_internal
signing_attempt_state_of_sample_pair_challengeE
signing_attempt_state_of_sample_pair_highbitsE
signing_attempt_state_of_sample_pair_lowbitsE
signing_attempt_state_of_sample_pair_tokenE
signing_attempt_state_signature_of_fields
\end{lstlisting}

\end{formaltheorem}

\begin{formaltheorem}[EasyCrypt line 2026]
\noindent\textbf{EasyCrypt name.}
\begin{lstlisting}[language=EasyCrypt,basicstyle=\ttfamily\scriptsize]
signing_attempt_state_of_sample_pair_fields_match_signature
\end{lstlisting}
\noindent\textbf{Checked EasyCrypt statement.}
\begin{lstlisting}[language=EasyCrypt,basicstyle=\ttfamily\scriptsize]
lemma signing_attempt_state_of_sample_pair_fields_match_signature
   md sk m ctx coins sample :
  signing_attempt_state_fields_match_signature
    (public_key_of_secret md sk) m ctx
    (signing_attempt_state_of_sample_pair md sk m ctx coins sample).
\end{lstlisting}
\noindent\textbf{Formal mathematical reading.}
Mathematical theorem. This lemma is a universally quantified proposition over the variables shown in the EasyCrypt statement.

\noindent\textbf{Role in the proof.}
Use in this chapter. It contributes directly to an advantage bound, bad-event bound, or real-arithmetic composition step.

\noindent\textbf{Proof-script interpretation.}
Proof-script reading. The machine proof decomposes the theorem into rewrites, program-logic obligations, relational equivalences, and arithmetic side conditions.
\emph{Main tactics visible in the proof script:} \ec{rewrite}.
\emph{Named identifiers syntactically cited in the proof block:}
\begin{lstlisting}[language=EasyCrypt,basicstyle=\ttfamily\scriptsize]
signing_attempt_state_fields_match_signature
signing_attempt_state_of_sample_pair_signatureE
signing_attempt_state_of_sample_pair_signature_of_fieldsE
\end{lstlisting}

\end{formaltheorem}

\begin{formaltheorem}[EasyCrypt line 2037]
\noindent\textbf{EasyCrypt name.}
\begin{lstlisting}[language=EasyCrypt,basicstyle=\ttfamily\scriptsize]
signing_attempt_state_of_sample_pair_field_accepts
\end{lstlisting}
\noindent\textbf{Checked EasyCrypt statement.}
\begin{lstlisting}[language=EasyCrypt,basicstyle=\ttfamily\scriptsize]
lemma signing_attempt_state_of_sample_pair_field_accepts
   md sk m ctx coins sample :
  signing_attempt_state_field_accepts md (public_key_of_secret md sk)
    m ctx (signing_attempt_state_of_sample_pair md sk m ctx coins sample).
\end{lstlisting}
\noindent\textbf{Formal mathematical reading.}
Mathematical theorem. This lemma is a universally quantified proposition over the variables shown in the EasyCrypt statement.

\noindent\textbf{Role in the proof.}
Use in this chapter. It contributes directly to an advantage bound, bad-event bound, or real-arithmetic composition step.

\noindent\textbf{Proof-script interpretation.}
Proof-script reading. The machine proof decomposes the theorem into rewrites, program-logic obligations, relational equivalences, and arithmetic side conditions.
\emph{Main tactics visible in the proof script:} \ec{rewrite}, \ec{split}, \ec{apply}.
\emph{Named identifiers syntactically cited in the proof block:}
\begin{lstlisting}[language=EasyCrypt,basicstyle=\ttfamily\scriptsize]
signing_attempt_state_field_accepts
signing_attempt_state_of_sample_pair_challenge_consistent
signing_attempt_state_of_sample_pair_fields_match_signature
signing_attempt_state_of_sample_pair_lowbits_consistent
signing_attempt_state_of_sample_pair_public_equation_consistent
\end{lstlisting}

\end{formaltheorem}

\begin{formaltheorem}[EasyCrypt line 2052]
\noindent\textbf{EasyCrypt name.}
\begin{lstlisting}[language=EasyCrypt,basicstyle=\ttfamily\scriptsize]
signing_attempt_state_of_sample_pair_valid_signature_accepts
\end{lstlisting}
\noindent\textbf{Checked EasyCrypt statement.}
\begin{lstlisting}[language=EasyCrypt,basicstyle=\ttfamily\scriptsize]
lemma signing_attempt_state_of_sample_pair_valid_signature_accepts
   md sk m ctx coins sample :
  signing_attempt_state_valid_signature_accepts md
    (signing_attempt_state_of_sample_pair md sk m ctx coins sample).
\end{lstlisting}
\noindent\textbf{Formal mathematical reading.}
Mathematical theorem. This lemma is a universally quantified proposition over the variables shown in the EasyCrypt statement.

\noindent\textbf{Role in the proof.}
Use in this chapter. It contributes directly to an advantage bound, bad-event bound, or real-arithmetic composition step.

\noindent\textbf{Proof-script interpretation.}
Proof-script reading. The machine proof decomposes the theorem into rewrites, program-logic obligations, relational equivalences, and arithmetic side conditions.
\emph{Main tactics visible in the proof script:} \ec{rewrite}.
\emph{Named identifiers syntactically cited in the proof block:}
\begin{lstlisting}[language=EasyCrypt,basicstyle=\ttfamily\scriptsize]
signing_attempt_state_of_sample_pair_signatureE
signing_attempt_state_valid_signature_accepts
valid_signature_sign_internal
\end{lstlisting}

\end{formaltheorem}

\begin{formaltheorem}[EasyCrypt line 2062]
\noindent\textbf{EasyCrypt name.}
\begin{lstlisting}[language=EasyCrypt,basicstyle=\ttfamily\scriptsize]
signing_attempt_state_of_sample_pair_accepts_current
\end{lstlisting}
\noindent\textbf{Checked EasyCrypt statement.}
\begin{lstlisting}[language=EasyCrypt,basicstyle=\ttfamily\scriptsize]
lemma signing_attempt_state_of_sample_pair_accepts_current
   md sk m ctx coins sample :
  signing_attempt_state_accepts md
    (signing_attempt_state_of_sample_pair md sk m ctx coins sample).
\end{lstlisting}
\noindent\textbf{Formal mathematical reading.}
Mathematical theorem. This lemma is a universally quantified proposition over the variables shown in the EasyCrypt statement.

\noindent\textbf{Role in the proof.}
Use in this chapter. It contributes directly to an advantage bound, bad-event bound, or real-arithmetic composition step.

\noindent\textbf{Proof-script interpretation.}
Proof-script reading. The machine proof decomposes the theorem into rewrites, program-logic obligations, relational equivalences, and arithmetic side conditions.
\emph{Main tactics visible in the proof script:} \ec{rewrite}, \ec{split}, \ec{apply}.
\emph{Named identifiers syntactically cited in the proof block:}
\begin{lstlisting}[language=EasyCrypt,basicstyle=\ttfamily\scriptsize]
signing_attempt_state_accepts
signing_attempt_state_of_sample_pair_hint_norm_accepts
signing_attempt_state_of_sample_pair_response_norm_accepts
signing_attempt_state_of_sample_pair_valid_signature_accepts
\end{lstlisting}

\end{formaltheorem}

\begin{formaltheorem}[EasyCrypt line 2075]
\noindent\textbf{EasyCrypt name.}
\begin{lstlisting}[language=EasyCrypt,basicstyle=\ttfamily\scriptsize]
signing_attempt_state_of_sample_pair_verification_accepts_current
\end{lstlisting}
\noindent\textbf{Checked EasyCrypt statement.}
\begin{lstlisting}[language=EasyCrypt,basicstyle=\ttfamily\scriptsize]
lemma signing_attempt_state_of_sample_pair_verification_accepts_current
   md sk m ctx coins sample :
  signing_attempt_state_verification_accepts md (public_key_of_secret md sk)
    m ctx (signing_attempt_state_of_sample_pair md sk m ctx coins sample).
\end{lstlisting}
\noindent\textbf{Formal mathematical reading.}
Mathematical theorem. This lemma is a universally quantified proposition over the variables shown in the EasyCrypt statement.

\noindent\textbf{Role in the proof.}
Use in this chapter. It contributes directly to an advantage bound, bad-event bound, or real-arithmetic composition step.

\noindent\textbf{Proof-script interpretation.}
Proof-script reading. The machine proof decomposes the theorem into rewrites, program-logic obligations, relational equivalences, and arithmetic side conditions.
\emph{Main tactics visible in the proof script:} \ec{rewrite}, \ec{split}, \ec{apply}.
\emph{Named identifiers syntactically cited in the proof block:}
\begin{lstlisting}[language=EasyCrypt,basicstyle=\ttfamily\scriptsize]
signing_attempt_state_of_sample_pair_accepts_current
signing_attempt_state_of_sample_pair_field_accepts
signing_attempt_state_verification_accepts
\end{lstlisting}

\end{formaltheorem}

\begin{formaltheorem}[EasyCrypt line 2086]
\noindent\textbf{EasyCrypt name.}
\begin{lstlisting}[language=EasyCrypt,basicstyle=\ttfamily\scriptsize]
signing_attempt_state_of_sample_pair_reference_rejection_accepts
\end{lstlisting}
\noindent\textbf{Checked EasyCrypt statement.}
\begin{lstlisting}[language=EasyCrypt,basicstyle=\ttfamily\scriptsize]
lemma signing_attempt_state_of_sample_pair_reference_rejection_accepts
   (md : mode) (sk : skey) (m : message) (ctx : context)
   (coins : random_coins) (sample : signing_sample_pair) :
  signing_attempt_state_reference_rejection_accepts md
    (signing_attempt_state_of_sample_pair md sk m ctx coins sample).
\end{lstlisting}
\noindent\textbf{Formal mathematical reading.}
Mathematical theorem. This lemma is a universally quantified proposition over the variables shown in the EasyCrypt statement.

\noindent\textbf{Role in the proof.}
Use in this chapter. It contributes directly to an advantage bound, bad-event bound, or real-arithmetic composition step.

\noindent\textbf{Proof-script interpretation.}
Proof-script reading. The machine proof decomposes the theorem into rewrites, program-logic obligations, relational equivalences, and arithmetic side conditions.
\emph{Main tactics visible in the proof script:} \ec{rewrite}, \ec{split}, \ec{apply}.
\emph{Named identifiers syntactically cited in the proof block:}
\begin{lstlisting}[language=EasyCrypt,basicstyle=\ttfamily\scriptsize]
signing_attempt_state_of_sample_pair_hint_norm_accepts
signing_attempt_state_of_sample_pair_pack_accepts
signing_attempt_state_of_sample_pair_reject1_accepts
signing_attempt_state_of_sample_pair_reject2_accepts
signing_attempt_state_reference_rejection_accepts
\end{lstlisting}

\end{formaltheorem}

\begin{formaltheorem}[EasyCrypt line 2102]
\noindent\textbf{EasyCrypt name.}
\begin{lstlisting}[language=EasyCrypt,basicstyle=\ttfamily\scriptsize]
signing_attempt_state_of_sample_pair_reference_rejection_no_abort
\end{lstlisting}
\noindent\textbf{Checked EasyCrypt statement.}
\begin{lstlisting}[language=EasyCrypt,basicstyle=\ttfamily\scriptsize]
lemma signing_attempt_state_of_sample_pair_reference_rejection_no_abort
   (md : mode) (sk : skey) (m : message) (ctx : context)
   (coins : random_coins) (sample : signing_sample_pair) :
  ! signing_attempt_state_reference_rejection_aborts md
      (signing_attempt_state_of_sample_pair md sk m ctx coins sample).
\end{lstlisting}
\noindent\textbf{Formal mathematical reading.}
Mathematical theorem. This lemma is a universally quantified proposition over the variables shown in the EasyCrypt statement.

\noindent\textbf{Role in the proof.}
Use in this chapter. It contributes directly to an advantage bound, bad-event bound, or real-arithmetic composition step.

\noindent\textbf{Proof-script interpretation.}
Proof-script reading. The machine proof decomposes the theorem into rewrites, program-logic obligations, relational equivalences, and arithmetic side conditions.
\emph{Main tactics visible in the proof script:} \ec{rewrite}.
\emph{Named identifiers syntactically cited in the proof block:}
\begin{lstlisting}[language=EasyCrypt,basicstyle=\ttfamily\scriptsize]
signing_attempt_state_of_sample_pair_reference_rejection_accepts
signing_attempt_state_reference_rejection_aborts
\end{lstlisting}

\end{formaltheorem}

\begin{formaltheorem}[EasyCrypt line 2112]
\noindent\textbf{EasyCrypt name.}
\begin{lstlisting}[language=EasyCrypt,basicstyle=\ttfamily\scriptsize]
signing_attempt_state_of_sample_pair_rejection_accepts_current
\end{lstlisting}
\noindent\textbf{Checked EasyCrypt statement.}
\begin{lstlisting}[language=EasyCrypt,basicstyle=\ttfamily\scriptsize]
lemma signing_attempt_state_of_sample_pair_rejection_accepts_current
   md sk m ctx coins sample :
  signing_attempt_state_rejection_accepts md (public_key_of_secret md sk)
    m ctx (signing_attempt_state_of_sample_pair md sk m ctx coins sample).
\end{lstlisting}
\noindent\textbf{Formal mathematical reading.}
Mathematical theorem. This lemma is a universally quantified proposition over the variables shown in the EasyCrypt statement.

\noindent\textbf{Role in the proof.}
Use in this chapter. It contributes directly to an advantage bound, bad-event bound, or real-arithmetic composition step.

\noindent\textbf{Proof-script interpretation.}
Proof-script reading. The machine proof decomposes the theorem into rewrites, program-logic obligations, relational equivalences, and arithmetic side conditions.
\emph{Main tactics visible in the proof script:} \ec{rewrite}, \ec{split}, \ec{apply}.
\emph{Named identifiers syntactically cited in the proof block:}
\begin{lstlisting}[language=EasyCrypt,basicstyle=\ttfamily\scriptsize]
signing_attempt_state_of_sample_pair_reference_rejection_accepts
signing_attempt_state_of_sample_pair_verification_accepts_current
signing_attempt_state_rejection_accepts
\end{lstlisting}

\end{formaltheorem}

\begin{formaltheorem}[EasyCrypt line 2123]
\noindent\textbf{EasyCrypt name.}
\begin{lstlisting}[language=EasyCrypt,basicstyle=\ttfamily\scriptsize]
signing_attempt_state_of_sample_pair_rejection_no_abort_current
\end{lstlisting}
\noindent\textbf{Checked EasyCrypt statement.}
\begin{lstlisting}[language=EasyCrypt,basicstyle=\ttfamily\scriptsize]
lemma signing_attempt_state_of_sample_pair_rejection_no_abort_current
   md sk m ctx coins sample :
  ! signing_attempt_state_rejection_aborts md (public_key_of_secret md sk)
      m ctx (signing_attempt_state_of_sample_pair md sk m ctx coins sample).
\end{lstlisting}
\noindent\textbf{Formal mathematical reading.}
Mathematical theorem. This lemma is a universally quantified proposition over the variables shown in the EasyCrypt statement.

\noindent\textbf{Role in the proof.}
Use in this chapter. It contributes directly to an advantage bound, bad-event bound, or real-arithmetic composition step.

\noindent\textbf{Proof-script interpretation.}
Proof-script reading. The machine proof decomposes the theorem into rewrites, program-logic obligations, relational equivalences, and arithmetic side conditions.
\emph{Main tactics visible in the proof script:} \ec{rewrite}.
\emph{Named identifiers syntactically cited in the proof block:}
\begin{lstlisting}[language=EasyCrypt,basicstyle=\ttfamily\scriptsize]
signing_attempt_state_of_sample_pair_rejection_accepts_current
signing_attempt_state_rejection_aborts
\end{lstlisting}

\end{formaltheorem}

\begin{formaltheorem}[EasyCrypt line 2132]
\noindent\textbf{EasyCrypt name.}
\begin{lstlisting}[language=EasyCrypt,basicstyle=\ttfamily\scriptsize]
dexact_hyperball_signing_attempt_state_reference_rejection_accepts
\end{lstlisting}
\noindent\textbf{Checked EasyCrypt statement.}
\begin{lstlisting}[language=EasyCrypt,basicstyle=\ttfamily\scriptsize]
lemma dexact_hyperball_signing_attempt_state_reference_rejection_accepts
   (md : mode) (sk : skey) (m : message) (ctx : context)
   (st : signing_attempt_state) :
  st \in dexact_hyperball_signing_attempt_state md sk m ctx =>
  signing_attempt_state_reference_rejection_accepts md st.
\end{lstlisting}
\noindent\textbf{Formal mathematical reading.}
Mathematical theorem. This is an implication, normally turning one invariant, validity predicate, or proof obligation into another.

\noindent\textbf{Role in the proof.}
Use in this chapter. It proves an exact game replacement or simulator equivalence.

\noindent\textbf{Proof-script interpretation.}
Proof-script reading. The machine proof decomposes the theorem into rewrites, program-logic obligations, relational equivalences, and arithmetic side conditions.
\emph{Main tactics visible in the proof script:} \ec{rewrite}.
\emph{Named identifiers syntactically cited in the proof block:}
\begin{lstlisting}[language=EasyCrypt,basicstyle=\ttfamily\scriptsize]
coins
dexact_hyperball_signing_attempt_state_supportP
sample
signing_attempt_state_of_sample_pair_reference_rejection_accepts
stE
st_supp
\end{lstlisting}

\end{formaltheorem}

\begin{formaltheorem}[EasyCrypt line 2144]
\noindent\textbf{EasyCrypt name.}
\begin{lstlisting}[language=EasyCrypt,basicstyle=\ttfamily\scriptsize]
dexact_hyperball_signing_attempt_state_reference_rejection_no_abort
\end{lstlisting}
\noindent\textbf{Checked EasyCrypt statement.}
\begin{lstlisting}[language=EasyCrypt,basicstyle=\ttfamily\scriptsize]
lemma dexact_hyperball_signing_attempt_state_reference_rejection_no_abort
   (md : mode) (sk : skey) (m : message) (ctx : context)
   (st : signing_attempt_state) :
  st \in dexact_hyperball_signing_attempt_state md sk m ctx =>
  ! signing_attempt_state_reference_rejection_aborts md st.
\end{lstlisting}
\noindent\textbf{Formal mathematical reading.}
Mathematical theorem. This is an implication, normally turning one invariant, validity predicate, or proof obligation into another.

\noindent\textbf{Role in the proof.}
Use in this chapter. It proves an exact game replacement or simulator equivalence.

\noindent\textbf{Proof-script interpretation.}
Proof-script reading. The machine proof decomposes the theorem into rewrites, program-logic obligations, relational equivalences, and arithmetic side conditions.
\emph{Main tactics visible in the proof script:} \ec{rewrite}.
\emph{Named identifiers syntactically cited in the proof block:}
\begin{lstlisting}[language=EasyCrypt,basicstyle=\ttfamily\scriptsize]
ctx
dexact_hyperball_signing_attempt_state_reference_rejection_accepts
md
signing_attempt_state_reference_rejection_aborts
sk
st
st_supp
\end{lstlisting}

\end{formaltheorem}

\begin{formaltheorem}[EasyCrypt line 2156]
\noindent\textbf{EasyCrypt name.}
\begin{lstlisting}[language=EasyCrypt,basicstyle=\ttfamily\scriptsize]
dexact_hyperball_signing_attempt_state_rejection_accepts
\end{lstlisting}
\noindent\textbf{Checked EasyCrypt statement.}
\begin{lstlisting}[language=EasyCrypt,basicstyle=\ttfamily\scriptsize]
lemma dexact_hyperball_signing_attempt_state_rejection_accepts
   (md : mode) (sk : skey) (m : message) (ctx : context)
   (st : signing_attempt_state) :
  st \in dexact_hyperball_signing_attempt_state md sk m ctx =>
  signing_attempt_state_rejection_accepts md (public_key_of_secret md sk)
    m ctx st.
\end{lstlisting}
\noindent\textbf{Formal mathematical reading.}
Mathematical theorem. This is an implication, normally turning one invariant, validity predicate, or proof obligation into another.

\noindent\textbf{Role in the proof.}
Use in this chapter. It proves an exact game replacement or simulator equivalence.

\noindent\textbf{Proof-script interpretation.}
Proof-script reading. The machine proof decomposes the theorem into rewrites, program-logic obligations, relational equivalences, and arithmetic side conditions.
\emph{Main tactics visible in the proof script:} \ec{rewrite}.
\emph{Named identifiers syntactically cited in the proof block:}
\begin{lstlisting}[language=EasyCrypt,basicstyle=\ttfamily\scriptsize]
coins
dexact_hyperball_signing_attempt_state_supportP
sample
signing_attempt_state_of_sample_pair_rejection_accepts_current
stE
st_supp
\end{lstlisting}

\end{formaltheorem}

\begin{formaltheorem}[EasyCrypt line 2169]
\noindent\textbf{EasyCrypt name.}
\begin{lstlisting}[language=EasyCrypt,basicstyle=\ttfamily\scriptsize]
dexact_hyperball_signing_attempt_state_rejection_no_abort
\end{lstlisting}
\noindent\textbf{Checked EasyCrypt statement.}
\begin{lstlisting}[language=EasyCrypt,basicstyle=\ttfamily\scriptsize]
lemma dexact_hyperball_signing_attempt_state_rejection_no_abort
   (md : mode) (sk : skey) (m : message) (ctx : context)
   (st : signing_attempt_state) :
  st \in dexact_hyperball_signing_attempt_state md sk m ctx =>
  ! signing_attempt_state_rejection_aborts md (public_key_of_secret md sk)
      m ctx st.
\end{lstlisting}
\noindent\textbf{Formal mathematical reading.}
Mathematical theorem. This is an implication, normally turning one invariant, validity predicate, or proof obligation into another.

\noindent\textbf{Role in the proof.}
Use in this chapter. It proves an exact game replacement or simulator equivalence.

\noindent\textbf{Proof-script interpretation.}
Proof-script reading. The machine proof decomposes the theorem into rewrites, program-logic obligations, relational equivalences, and arithmetic side conditions.
\emph{Main tactics visible in the proof script:} \ec{rewrite}.
\emph{Named identifiers syntactically cited in the proof block:}
\begin{lstlisting}[language=EasyCrypt,basicstyle=\ttfamily\scriptsize]
ctx
dexact_hyperball_signing_attempt_state_rejection_accepts
md
signing_attempt_state_rejection_aborts
sk
st
st_supp
\end{lstlisting}

\end{formaltheorem}

\begin{formaltheorem}[EasyCrypt line 2182]
\noindent\textbf{EasyCrypt name.}
\begin{lstlisting}[language=EasyCrypt,basicstyle=\ttfamily\scriptsize]
dexact_hyperball_signing_attempt_state_rejection_abort_zero
\end{lstlisting}
\noindent\textbf{Checked EasyCrypt statement.}
\begin{lstlisting}[language=EasyCrypt,basicstyle=\ttfamily\scriptsize]
lemma dexact_hyperball_signing_attempt_state_rejection_abort_zero
  md sk m ctx :
  mu (dexact_hyperball_signing_attempt_state md sk m ctx)
    (signing_attempt_state_rejection_aborts md
       (public_key_of_secret md sk) m ctx) =
  0%r.
\end{lstlisting}
\noindent\textbf{Formal mathematical reading.}
Mathematical theorem. This is an equality or definitional expansion used for rewriting and normalization.

\noindent\textbf{Role in the proof.}
Use in this chapter. It proves an exact game replacement or simulator equivalence.

\noindent\textbf{Proof-script interpretation.}
Proof-script reading. The machine proof decomposes the theorem into rewrites, program-logic obligations, relational equivalences, and arithmetic side conditions.
\emph{Main tactics visible in the proof script:} \ec{have}, \ec{apply}, \ec{rewrite}, \ec{smt}.
\emph{Named identifiers syntactically cited in the proof block:}
\begin{lstlisting}[language=EasyCrypt,basicstyle=\ttfamily\scriptsize]
coins
ctx
dexact_hyperball_signing_attempt_state
dexact_hyperball_signing_attempt_state_supportP
le0
md
mu
mu0
mu_bounded
mu_le
pred0
public_key_of_secret
sample
signing_attempt_state_of_sample_pair_rejection_no_abort_current
signing_attempt_state_rejection_aborts
sk
st
stE
st_supp
\end{lstlisting}

\end{formaltheorem}

\medskip
\noindent\emph{End of generated formal inventory for \file{HAETAE_Rejection.ec}.}
